\documentclass[9pt]{article}

% --- Page layout (tight and readable) ---
\usepackage[a4paper,margin=0.25in]{geometry}
% --- Math ---
\usepackage{amsmath,amssymb,amsfonts,mathtools}
\setlength{\jot}{4pt} % controls spacing between aligned equations (default is large)

% --- Physics (optional but useful) ---
\usepackage{braket}
\usepackage{slashed}
\usepackage{cancel}
\newcommand{\eqname}[1]{\tag*{#1}}
% --- Graphics ---
\usepackage{graphicx}
\usepackage{tikz}
\usepackage{pgfplots}

% --- Utilities ---
\usepackage{xcolor}
\usepackage{hyperref}
\usepackage{comment}

% --- Custom commands (keep yours) ---
\newcommand{\CGCket}{|{\rm CGC}\rangle}
\newcommand{\CGCbra}{\langle {\rm CGC}|}

\def\wt#1{\widetilde{#1}}
\def\mf#1{\mathbb{#1}}

\newcommand{\A}[3]{\mf{A}^{#2}_{#3}(#1)}

% --- added: corrections in red (set \markedfalse for a clean copy) ---
\newif\ifmarked \markedtrue
\newcommand{\chg}[1]{\ifmarked{\color{red}#1}\else{#1}\fi}
\numberwithin{equation}{section}

\title{Trijet cross section at large $N_{c}$ (Massive case)}
\author{}
\date{}

\begin{document}

\maketitle
\noindent\ifmarked{\color{red}\textbf{Marked version.}} Corrections to the supplied \texttt{Trijet\_massive.tex} are printed in {\color{red}red}; Appendix~\ref{app:corr} lists each with its equation number, page and original line number (all clickable). \fi Section~\ref{sec:el} (new) sets up and evaluates the Fourier transform of the elastic (dipole-only, no $F$-ratio) terms of the cross sections; the results are written as terms without mass $+$ terms with mass, Eq.~\eqref{eq:resel}.
\section{Setup}
The cross section for trijet production is given by
\begin{align}
    \frac{d^{3}N}{d^{3}k_{1}\,d^{3}k_{2}\, d^{3}k_{3}}\, (2\pi)\,\delta\left(q^{+}-k^{+}_{1}-k^{+}_{2}-k^{+}_{3}\right) & =  \frac{1}{2}\,\int_{w,w'}\, e^{-iq\cdot (w-w')} 
\bra{\gamma^{i}_{T}(Q,q^{+},w')}\,\mathcal{N}_{q}(k_{1})\,\mathcal{N}_{\Bar{q}}(k_{2})\,\mathcal{N}_{g}(k_{3})\,\ket{\gamma^{i}_{T}(Q,q^{+},w)}_{q\Bar{q}g}
\end{align}
\subsection{Kinematics}
The virtual photon with the four momentum as $q^{\mu} = \left(q^{+},-\frac{Q^{2}}{2q^{+}},q\right)$. Let us define the kinematics in the following:
\begin{table}[h!]
\centering
\begin{tabular}{||c c c c||} 
 \hline
 Component & Longitudinal momentum fraction & Transverse momentum & mass \\ [0.5ex] 
 \hline\hline
 Quark & $\zeta = \frac{k^{+}_{1}}{q^{+}}$  & $k_{1}$ & $m_{f}$ \\ 
 Anti Quark & $\eta = 1-\zeta-\xi$  &  $k_{2}$ & $m_{f}$ \\
 Gluon & $\xi = \frac{k^{+}_{3}}{q^{+}}$ & $k_{3}$ & $0$ \\
 \hline
\end{tabular}
\label{table:1}
\end{table}
The onshell denominators are $E_{q}(k) = \frac{k^{2}+m_{f}^{2}}{2k^{+}}$, $E_{\Bar{q}}(p) = \frac{p^{2}+m_{f}^{2}}{2p^{+}}$, $E_{g}(p) = \frac{p^{2}}{2p^{+}}$.  
\subsection{Matrix Elements}
The matrix elements that are used in this process are as follows:
\begin{enumerate}
\item Virtual photon splitting into a quark and anti quark pair
\begin{align}
    \bra{q^{\alpha}_{\lambda_{1}}(k)\,\Bar{q}^{\beta}_{\lambda_{2}}(p)}H_{\gamma\rightarrow q\Bar{q}}\ket{\gamma_{\lambda}(q)} & = (2\pi)^{3}\, \delta^{(3)}(q-k-p)\, \frac{e\,e_{f}\,\delta_{\alpha\beta}}{2\,\sqrt{2q^{+}}}\, \chi^{\dagger}_{\lambda_{1}}\,\left[\frac{2q^{\lambda}}{q^{+}}-\sigma^{\lambda}\,\frac{\sigma\cdot p+m_{f}\,\sigma^{3}}{p^{+}}- \frac{\sigma\cdot k-m_{f}\,\sigma^{3}}{k^{+}}\sigma^{\lambda}\right]\,\chi_{\lambda_{2}}
\end{align}
rewriting it using the delta function,
$$k^{+} = \zeta\,q^{+}, \, p^{+} = (1-\zeta)\,q^{+}, \, \Tilde{k} = k-\zeta\,q$$
we get
$$\left[\frac{2q^{\lambda}}{q^{+}}-\frac{\sigma\cdot k - m_{f}\,\sigma^{3}}{k^{+}}\,\sigma^{\lambda}-\sigma^{\lambda}\, \frac{\sigma\cdot (q-k)+m_{f}\,\sigma^{3}}{(q^{+}-k^{+})}\right] = \frac{\left[(2\zeta -1)\,\delta^{\lambda\,j} + i\epsilon^{\lambda\, j}\sigma^{3}\right]\,\Tilde{k}^{j}+m_{f}\,\sigma^{3}\,\sigma^{\lambda}}{\zeta\,(1-\zeta)\,q^{+}}$$
where we used $\sigma^{i}\sigma^{j} = \delta^{ij} +i\epsilon^{ij}\sigma^{3}$ and $\sigma^{\lambda}\sigma^{3} = - \sigma^{3}\sigma^{\lambda}$
and $$\chi^{\dagger}_{\lambda_{1}}\left[(2\zeta -1)\,\delta^{\lambda\,j} + i\epsilon^{\lambda\, j}\sigma^{3}\right]\chi_{\lambda_{2}} = \delta_{\lambda_{1}\,\lambda_{2}}\left[(2\zeta - 1)\delta^{\lambda\,j} +2i\epsilon^{\lambda\,j}\,\lambda_{1}\right]$$
$$\chi^{\dagger}_{\lambda_{1}}\,\sigma^{3}\sigma^{\lambda}\,\chi_{\lambda_{2}} = \delta_{\lambda_{2},-\lambda_{1}} \left[2\lambda_{1}\delta^{\lambda,1} - i\delta^{\lambda,2}\right] \equiv \mu^{\lambda}_{\lambda_{1}\lambda_{2}}$$
The mass term flips the helicity $(\lambda_{2} = -\lambda_{1})$. The matrix element defined above can be rewritten as follows:
\begin{align}
    \bra{q^{\alpha}_{\lambda_{1}}(k)\,\Bar{q}^{\beta}_{\lambda_{2}}(p)}H_{\gamma\rightarrow q\Bar{q}}\ket{\gamma_{\lambda}(q)} & = (2\pi)^{3}\, \delta^{(3)}(q-k-p)\, \frac{e\,e_{f}\,\delta_{\alpha\beta}\,\bigg\{\left[(2\zeta -1)\,\delta^{\lambda\,j}+2i\epsilon^{\lambda\,j}\,\lambda_{1}\right]\,\delta_{\lambda_{1}\,\lambda_{2}}\Tilde{k}^{j}+m_{f}\,\mu^{\lambda}_{\lambda_{1}\lambda_{2}}\bigg\}}{2\,\sqrt{2\,q^{+\,3}}\, \zeta\,(1-\zeta)}
\end{align}
\item  Gluon emission from anti - quark
\begin{align}
    \bra{q^{\alpha'}_{\lambda'_{1}}(k')\,\Bar{q}^{\beta'}_{\lambda'_{2}}(p')\,g^{a'}_{i'}(q')}H_{\Bar{q}g}\ket{q^{\alpha}_{\lambda_{1}}(k)\,\Bar{q}^{\beta}_{\lambda_{2}}(p)} & = -\,(2\pi)^{6}\, \delta^{(3)}(k-k')\, \delta^{(3)}(p-p'-q')\,\delta_{\alpha'\,\alpha}\,\delta_{\lambda_{1}\,\lambda'_{1}}\, \frac{g\,t^{a'}_{\beta\,\beta'}}{2\sqrt{2\,q'^{+}}} \notag \\ & \chi^{\dagger}_{\lambda_{2}}\,\left[\frac{2\,q'^{i'}}{q'^{+}}-\sigma^{i'}\,\frac{\sigma\cdot p' +m_{f}\,\sigma^{3}}{p'^{+}}-\frac{\sigma\cdot p +m_{f}\,\sigma^{3}}{p^{+}}\,\sigma^{i'}\right]\,\chi_{\lambda'_{2}}
\end{align}
with $p^{+} = (1-\zeta)\, q^{+},\, p'^{+} = \eta\,q^{+}, \, q'^{+} = \xi\,q^{+}$ and $\Tilde{p} = q' - \frac{\xi}{(1-\zeta)}\, p $. The terms proportional to $p$ cancel and from here onwards we term $\omega_{f} = \frac{\xi\, m_{f}}{(1-\zeta)}$. The reduced vertex is 
\begin{align}
\chi^{\dagger}_{\lambda}\,\left[\frac{2\,q'^{i'}}{q'^{+}}-\sigma^{i'}\,\frac{\sigma\cdot p' +m_{f}\,\sigma^{3}}{p'^{+}}-\frac{\sigma\cdot p +m_{f}\,\sigma^{3}}{p^{+}}\,\sigma^{i'}\right]\,\chi_{\lambda_{2}} & = \frac{\tau^{i'n}_{\lambda\lambda_{2}}(\xi,\eta)\,\Tilde{p}^{n}+\xi\,\omega_{f}\,\mu^{i'}_{\lambda\lambda_{2}}}{\xi\,\eta\,q^{+}}
\end{align}
and 
\begin{align}
m_{f}\,\sigma^{3}\,\sigma^{i'}\,\left[\frac{1}{p'^{+}}-\frac{1}{p^{+}}\right] & = \frac{\xi^{2}\,m_{f}}{(1-\zeta)}\,\frac{\sigma^{3}\sigma^{i'}}{\xi\,\eta\,q^{+}}
\end{align}
\item  Instantaneous contribution
\begin{align}
    \bra{q^{\alpha}_{\lambda_{1}}(k)\,\Bar{q}^{\beta}_{\lambda_{2}}(s)\,g^{b}_{j}(p)}\,H_{\gamma\rightarrow q\Bar{q}g}\ket{\gamma} & = -(2\pi)^{3}\, \delta^{(3)}(q-k-p-s)\, \frac{g\,e\,e_{f}\,\delta_{\lambda_{1}\lambda_{2}}\,t^{b}_{\alpha\beta}}{4\,\sqrt{q^{+}p^{+}}}\, \frac{\delta^{ij}+2i\,\epsilon^{ij}\lambda_{1}}{p^{+}+s^{+}}
\end{align}
\end{enumerate}
\subsection{Notations}
$$\Tilde{k} = k - \zeta\,q,\hspace{2em}\, \Tilde{p} = p - \frac{\xi}{1-\zeta}\, (q-k),\hspace{2em}\,\Tilde{Q}^{2} = \zeta(1-\zeta)\, Q^{2}$$
$$\epsilon^{2}_{f} = \Tilde{Q}^{2}+ m^{2}_{f} = \zeta\, (1-\zeta)\,Q^{2} +m^{2}_{f}, \hspace{2em} \omega_{f} = \frac{\xi\,m_{f}}{(1-\zeta)}, \hspace{2em} E^{2}_{f} = Q^{2} +m^{2}_{f}\,\left(\frac{1}{\zeta} + \frac{1}{\eta}\right),\hspace{2em} a = \frac{\xi\,\eta}{\zeta\,(1-\zeta)^{2}}$$
$$\varphi^{ij}_{\lambda_{1}\lambda}(\zeta) = \delta_{\lambda_{1}\lambda}\left[(2\zeta -1)\delta^{ij} +2i\epsilon^{ij}\lambda_{1}\right],\hspace{3em} \tau^{mn}_{\lambda \lambda_{2}}(\xi,\eta) = \delta_{\lambda \lambda_{2}}\,\left[(2(1-\zeta)-\xi)\delta^{mn}+2i\,\xi\,\epsilon^{mn}\,\lambda\right] $$
$$\mu^{i}_{\lambda\lambda'} = \delta_{\lambda',-\lambda}\left[2\lambda\delta^{i1} - i\delta^{i2}\right] = \chi^{\dagger}_{\lambda}\sigma^{3}\sigma^{i}\chi_{\lambda'},\hspace{2em} \sum_{i}\,\mu^{i}_{\lambda\lambda'}\mu^{i*}_{\lambda\lambda'} = 2\delta_{\lambda',-\lambda}$$
Here $\epsilon_{f}$ is the usual mass shifted photon scale. The photon vertex and the anti quark - gluon vertex are:
$$\mathrm{photon}: \varphi^{ij}_{\lambda_{1}\lambda}(\zeta)\Tilde{k}^{j} + m_{f}\mu^{i}_{\lambda_{1}\lambda},\hspace{2em} \mathrm{gluon\, off\, the\, antiquark}: \tau^{mn}_{\lambda\lambda_{2}}(\xi,\eta)\, \Tilde{p}^{n} + \xi\, \omega_{f}\,\mu^{m}_{\lambda\lambda_{2}}$$
The energy denominator is as follows:
$$E_{q}+E_{\Bar{q}}-E_{\gamma} = \frac{\Tilde{k}^{2}+\Tilde{Q}^{2}+m^{2}_{f}}{2\zeta\, (1-\zeta)\,q^{+}} = \frac{\Tilde{k}^{2}+\epsilon^{2}_{f}}{2\,\zeta(1-\zeta)\,q^{+}}$$
$$E_{q_{1}}+E_{\Bar{q}_{1}}+E_{g_{1}}-E_{\gamma} = \frac{(1-\zeta) (\Tilde{p}^{2}+\omega^{2}_{f})}{2\,\xi\, \eta\,q^{+}} + \frac{(\Tilde{k}^{2}+\epsilon_{f}^{2})}{2\zeta\,(1-\zeta)\,q^{+}} = \frac{\zeta\,(1-\zeta)^{2} (\Tilde{p}^{2}+\omega^{2}_{f}) + \xi\,\eta\, (\Tilde{k}^{2}+\epsilon^{2}_{f})}{2\,\zeta\,(1-\zeta)\,\xi\,\eta\,q^{+}}$$
\begin{align}
    E_{q_{1}}+E_{\Bar{q}_{1}}+E_{g_{1}}-E_{q}-E_{\Bar{q}} & = \frac{p^{2}}{2q^{+}\,\xi}+\frac{(q-k-p)^{2}+m^{2}_{f}}{2q^{+}\, \eta}-\frac{(q-k)^{2}+m^{2}_{f}}{2q^{+}\, (1-\zeta)} = \frac{\left[(1-\zeta)\,p - \xi\,(q-k)\right]^{2}}{2\xi\,(1-\zeta)\,\eta\,q^{+}} +\frac{m^{2}_{f}}{2q^{+}}\,\left(\frac{1}{\eta} - \frac{1}{1-\zeta}\right)\notag \\ & = \frac{(1-\zeta)^{2}\,\Tilde{p}^{2} + \xi^{2}\,m^{2}_{f}}{2\,\xi(1-\zeta)\,\eta\,q^{+}}
\end{align}
Simplifying it we get
$$E_{q_{1}}+E_{\Bar{q}_{1}}+E_{g_{1}}-E_{q}-E_{\Bar{q}} = \frac{(1-\zeta)\, (\Tilde{p}^{2}+\omega^{2}_{f})}{2\xi\,\eta\,q^{+}}$$
Also, for instantaneous term involving gluon emission is given by
\begin{align}
    \zeta(1-\zeta)^{2}\,\omega^{2}_{f} + \xi\,\eta\,\epsilon^{2}_{f} = \zeta\,(1-\zeta)\, \xi\, \eta\, E^{2}_{f}, \hspace{2em} E^{2}_{f} = Q^{2} + \frac{m^{2}_{f}}{\zeta} + \frac{m^{2}_{f}}{\eta}
\end{align}

\section{Outgoing component of wave function through LCPT}
\begin{enumerate}
\item Anti quark emits a gluon before the shock wave
\begin{figure}[h]
    \centering
    \includegraphics[width=0.2\linewidth]{Wave functions/1.jpeg}
    \label{fig:1}
\end{figure}
\begin{align}
    \ket{\gamma}^{(1)}_{q\Bar{q}g} & = \ket{q_{2}\Bar{q}_{2}g_{2}}\,\bra{q_{2}\Bar{q}_{2}g_{2}}\,S\,\ket{q_{1}\Bar{q}_{1}g_{1}}\times \frac{\bra{q_{1}\Bar{q}_{1}g_{1}}H_{\Bar{q}g}\ket{q\Bar{q}}}{E_{\Bar{q}_{1}}+E_{q_{1}}+E_{g_{1}}-E_{\gamma}} \times \frac{\bra{q\Bar{q}}H_{\gamma\rightarrow q\Bar{q}}\ket{\gamma}}{E_{q}+E_{\Bar{q}}-E_{\gamma}}
\end{align}
\begin{align}
    \ket{\gamma}^{(1)}_{q\Bar{q}g} & = \int\, d^{3}k''\, d^{3}p''\, d^{3}r\,d^{3}k'\, d^{3}p'\, d^{3}q'\, d^{3}k\, d^{3}p\, \ket{q^{\alpha''}_{\lambda''_{1}}(k'')\,\Bar{q}^{\beta''}_{\lambda''_{2}}(p'')\,g^{b}_{j}(r)}\,\bra{q^{\alpha''}_{\lambda''_{1}}(k'')\,\Bar{q}^{\beta''}_{\lambda''_{2}}(p'')\,g^{b}_{j}(r)}\,S\,\ket{q^{\alpha'}_{\lambda'_{1}}(k')\,\Bar{q}^{\beta'}_{\lambda'_{2}}(p')\,g^{a'}_{i'}(q')}\times\notag \\ & \frac{\bra{q^{\alpha'}_{\lambda'_{1}}(k')\,\Bar{q}^{\beta'}_{\lambda'_{2}}(p')\,g^{a'}_{i'}(q')}H_{\Bar{q}g}\ket{q^{\alpha}_{\lambda_{1}}(k)\,\Bar{q}^{\beta}_{\lambda_{2}}(p)}}{E_{\Bar{q}}(p')+E_{q}(k')+E_{g}(q')-E_{\gamma}} \times \frac{\bra{q^{\alpha}_{\lambda_{1}}(k)\,\Bar{q}^{\beta}_{\lambda_{2}}(p)}H_{\gamma\rightarrow q\Bar{q}}\ket{\gamma_{\lambda}(q)}}{E_{\Bar{q}}(p)+E_{q}(k)-E_{\gamma}(q,Q)}
\end{align} 
We will first evaluate the matrix elements without scattering (i.e., $S = 1$). With the massive vertices and energy denominators of Sec. 1, we get the following
\begin{align}
    \ket{\gamma_{i}\,(Q,q^{+},q)}^{(1\,a)}_{q\Bar{q}g} & = -\,\frac{1}{(2\pi)^{6}}\, \int\, d^{2}\Tilde{k}\, d^{2}\Tilde{p}\, \int_{0}^{1}\, d\zeta\, \int_{0}^{1-\zeta}\, d\xi\,\ket{q^{\alpha}_{\lambda_{1}}(\zeta q^{+},k)\Bar{q}^{\beta}_{\lambda_{2}}(\eta\,q^{+},q-k-p)g^{a}_{m}(\xi q^{+},p)} \frac{e\,e_{f}\,g\,q^{+}\,\zeta(1-\zeta)}{2\sqrt{\xi}\,(\Tilde{k}^{2}+\epsilon^{2}_{f})}\notag \\ & \frac{\left[\varphi^{ij}_{\lambda_{1}\lambda}(\zeta)\, \Tilde{k}^{j}+m_{f}\,\mu^{i}_{\lambda_{1}\lambda}\right]\,\left[\tau^{mn}_{\lambda \lambda_{2}}(\xi,\eta)\, \Tilde{p}^{n}\,+ \xi\, \omega_{f}\,\mu^{m}_{\lambda\lambda_{2}}\right] t^{a}_{\alpha\beta}}{\left[\zeta(1-\zeta)^{2}\, (\Tilde{p}^{2} +\omega^{2}_{f}) +\xi\,\eta\,(\Tilde{k}^{2}+\epsilon_{f}^{2})\right]}
\end{align} 
However, Wilson lines can be obtained only if the wave function is expressed in transverse coordinates. Therefore, we perform a Fourier transform to rewrite the wave function in transverse coordinate space. Performing Fourier transform for transverse co-ordinates,
\begin{align}
   \ket{q^{\alpha}_{\lambda_{1}}(k)\Bar{q}^{\beta}_{\lambda_{2}}(q-k-p)g^{a}_{m}(p)} & = \int_{x,y,z}\, e^{ik\cdot x}\, e^{ip\cdot z}\, e^{i(q-k-p)\cdot y}\,\ket{q^{\alpha}_{\lambda_{1}}(x)\Bar{q}^{\beta}_{\lambda_{2}}(y)g^{a}_{m}(z)}
\end{align}
Rewriting $k$, $p$ in terms of $\Tilde{k}$ and $\Tilde{p}$ respectively, we get
\begin{align}
    e^{i\left[\zeta\,x +\eta\, y + \xi\,z\right]\cdot q}\,e^{-i\,(y-z)\cdot \Tilde{p}}\, e^{i\left[(1-\zeta)x-\eta y-\xi z\right]\cdot \frac{\Tilde{k}}{(1-\zeta)}} & = e^{i\,\left(\zeta\,x+\eta\,y + \xi\,z\right)\cdot q}\,e^{-i\,(y-z)\cdot \Tilde{p}}\, e^{i\,\left(x-\frac{\eta\,y + \xi\,z}{(1-\zeta)}\right)\cdot \Tilde{k}} \notag \\ & = e^{i\,w\cdot q}\, e^{-i\,Y\cdot \Tilde{p}}\, e^{i\,R\cdot \Tilde{k}}.
\end{align}
The transverse coordinates of the photon $(w)$ and the intermediate antiquark ($y'$, before emitting the gluon) are fixed by centre of energy conservation at the two vertices:
\begin{align}
    w & = \zeta\,x +(1-\zeta)\,y' = \zeta\,x + \eta\,y + \xi\,z,
\end{align}
and 
\begin{align}
    y' & = \frac{\eta\,y + \xi\,z}{(1-\zeta)}.
\end{align}
The relative distances are defined as 
$$R = x - y'; \hspace{2em} Y = y-z$$
where we used 
$$(1-\zeta)\,x - \eta\,y -\xi\,z = (1-\zeta)\, (x-y') = (1-\zeta)\,R$$
We will use the following Fourier transforms:
\begin{align}
    \int\, \frac{d^{2}k}{2\pi}\, \frac{d^{2}p}{2\pi}\, \frac{k^{i}p^{j}}{k^{2}\, \left[ak^{2}+bp^{2}\right]}\,e^{ik\cdot X + ip\cdot Z} & = -\frac{X^{i}Z^{j}}{Z^{2}\,\left[bX^{2}+aZ^{2}\right]}
\end{align}
With masses the two denominators contain two different scales $\epsilon_{f}$ and $\omega_{f}$ and the transform is no longer a single Bessel function. We need four transforms, one with each helicity path. Define
\begin{align}
    (1-\zeta)\, \int\, \frac{d^{2}\Tilde{k}}{2\pi}\, \frac{d^{2}\Tilde{p}}{2\pi}\, \frac{e^{i\Tilde{k}\cdot R}\, e^{-i\Tilde{p}\cdot Y}\, \{\Tilde{k}^{j}\Tilde{p}^{n}, \Tilde{k}^{j}, \Tilde{p}^{n}, 1\}}{(\Tilde{k}^{2}+\epsilon^{2}_{f})\, \left[\zeta\,(1-\zeta)^{2}\, (\Tilde{p}^{2}+\omega_{f}^{2}) + \xi\,\eta\,(\Tilde{k}^{2} + \epsilon^{2}_{f})\right]} & = \Bigg\{R^{j}Y^{n}\mathcal{G}^{(11)}, iR^{j}\mathcal{G}^{(10)}, -iY^{n}\,\mathcal{G}^{(01)}, \mathcal{G}^{(00)}\Bigg\}
\end{align}
Using Schwinger parametrization, we get
\begin{align}
    \frac{1}{(\Tilde{k}^{2}+\epsilon^{2}_{f})\,\left[\zeta\,(1-\zeta)^{2}\, (\Tilde{p}^{2}+\omega_{f}^{2}) + \xi\,\eta\,(\Tilde{k}^{2} + \epsilon^{2}_{f})\right]} & = \frac{1}{\xi\,\eta}\, \int_{0}^{\infty}\, du \, \int_{0}^{\infty}\, dv\, e^{-(u+v)\, (\Tilde{k}^{2}+\epsilon^{2}_{f})-\frac{v}{a}(\Tilde{p}^{2}+\omega^{2}_{f})}
\end{align}
where $a = \frac{\xi\, \eta}{\left[\zeta\,(1-\zeta)^{2}\right]},\, K_{\nu}\, \text{is the Modified Bessel function of the second kind}$. Using the Gaussian integrals and including massive terms, we define the following notations, with $x$ as a Feynman parameter.
$$\mathcal{G}^{(\alpha\beta)}(R,Y) = \frac{a^{\beta}}{2\,\zeta(1-\zeta)}\,\int_{0}^{1}\, \frac{dx}{x^{1+\beta}}\, \left(\frac{A_{x}}{\Delta_{x}}\right)^{\frac{(\alpha + \beta)}{2}}\,K_{\alpha + \beta} (\sqrt{A_{x}\Delta_{x}}), \alpha, \beta \in \{0,1\}$$
$$A_{x} = \epsilon^{2}_{f} + \frac{x\,\omega^{2}_{f}}{a} = \zeta\,(1-\zeta)\,Q^{2} +m^{2}_{f}\,\left(1+\frac{x\zeta\xi}{\eta}\right)$$
$$\Delta_{x} = R^{2} + \frac{a\,Y^{2}}{x}$$
Explicitly
\begin{align}
    \mathcal{G}^{(11)} & = \frac{a}{2\,\zeta\, (1-\zeta)}\, \int_{0}^{1}\, \frac{dx}{x^{2}}\, \frac{A_{x}}{\Delta_{x}}\,K_{2}\left(\sqrt{A_{x}\Delta_{x}}\right)
\end{align}
\begin{align}
    \mathcal{G}^{(10)} & = \frac{1}{2\, \zeta\, (1-\zeta)}\, \int_{0}^{1}\, \frac{dx}{x}\, \sqrt{\frac{A_{x}}{\Delta_{x}}}\, K_{1}\left(\sqrt{A_{x}\Delta_{x}}\right)
\end{align}
\begin{align}
     \mathcal{G}^{(01)} & = \frac{a}{2\, \zeta\, (1-\zeta)}\, \int_{0}^{1}\, \frac{dx}{x^{2}}\, \sqrt{\frac{A_{x}}{\Delta_{x}}}\, K_{1}\left(\sqrt{A_{x}\,\Delta_{x}}\right)
\end{align}
\begin{align}
     \mathcal{G}^{(00)} & = \frac{1}{2\, \zeta\, (1-\zeta)}\, \int_{0}^{1}\, \frac{dx}{x}\,K_{0}\left(\sqrt{A_{x}\Delta_{x}}\right)
\end{align}
Now, implementing all these steps in the outgoing component of wave function, we get
\begin{align}
    \ket{\gamma_{i}\,(Q,q^{+},q)}^{(1\,a)}_{q\Bar{q}g} &  = -\,\frac{e\,e_{f}\,g\,q^{+}\,}{2\,(2\pi)^{4}}\, \int_{x,y,z} \int_{0}^{1}\, d\zeta\, \int_{0}^{1-\zeta}\, d\xi\,\ket{q^{\alpha}_{\lambda_{1}}(\zeta q^{+},x)\Bar{q}^{\beta}_{\lambda_{2}}(\eta\,q^{+},y)g^{a}_{m}(\xi q^{+},z)}\, \frac{\zeta\, t^{a}_{\alpha\beta}}{\sqrt{\xi}}\, \delta^{(2)}(w-c)\,\times \Psi^{im}_{\lambda_{1}\lambda_{2}}[\mathcal{G}]
\end{align}
where 
\begin{align}
  \Psi^{im}_{\lambda_{1}\lambda_{2}}[\mathcal{G}] & = \sum_{\lambda}\, \Bigg[\varphi^{ij}_{\lambda_{1}\lambda}(\zeta)\,\tau^{mn}_{\lambda\lambda_{2}}(\xi,\eta)\,R^{j}\,Y^{n}\,\mathcal{G}^{(11)}+i\,\xi\,\omega_{f}\,\varphi^{ij}_{\lambda_{1}\lambda}(\zeta)\,\mu^{m}_{\lambda\lambda_{2}}\,R^{j}\,\mathcal{G}^{(10)}\notag \\ & -i\,m_{f}\,\mu^{i}_{\lambda_{1}\lambda}\,\tau^{mn}_{\lambda\lambda_{2}}(\xi,\eta)\,Y^{n}\,\mathcal{G}^{(01)} + m_{f}\,\xi\,\omega_{f}\,\mu^{i}_{\lambda_{1}\lambda}\,\mu^{m}_{\lambda\lambda_{2}}\,\mathcal{G}^{(00)}\Bigg]  
\end{align}
$$\mathcal{D}^{2}(x,y,z,\zeta,\xi) = \zeta\,(1-\zeta)\,R^{2} + \frac{\xi\,\eta\,Y^{2}}{(1-\zeta)} = \zeta\,\eta\,(x-y)^{2} + \zeta\,\xi\,(x-z)^{2} + \eta\,\xi\,(y-z)^{2}$$
$$\mu^{i}_{\lambda\lambda'} = \delta_{\lambda',-\lambda}\left[2\lambda\delta^{i1} - i\delta^{i2}\right] = \chi^{\dagger}_{\lambda}\sigma^{3}\sigma^{i}\chi_{\lambda'}$$
$$\epsilon^{2}_{f} = \Tilde{Q}^{2}+ m^{2}_{f} = \zeta\, (1-\zeta)\,Q^{2} +m^{2}_{f}, \hspace{2em} \omega_{f} = \frac{\xi\,m_{f}}{(1-\zeta)}, \hspace{2em} E^{2}_{f} = Q^{2} +m^{2}_{f}\,\left(\frac{1}{\zeta} + \frac{1}{\eta}\right),\hspace{2em} a = \frac{\xi\,\eta}{\zeta\,(1-\zeta)^{2}}$$
$$\varphi^{ij}_{\lambda_{1}\lambda}(\zeta) = \delta_{\lambda_{1}\lambda}\left[(2\zeta -1)\delta^{ij} +2i\epsilon^{ij}\lambda_{1}\right],\hspace{3em} \tau^{mn}_{\lambda \lambda_{2}}(\xi,\eta) = \delta_{\lambda \lambda_{2}}\,\left[(2(1-\zeta)-\xi)\delta^{mn}+2i\,\xi\,\epsilon^{mn}\,\lambda\right] $$
and 
$$c = \zeta\,x +\eta\, y + \xi\,z$$ is the centre of energy of trijet system. Attaching  Wilson lines, we get 
\begin{align} \label{eq:1}
     \ket{\gamma_{i}\,(Q,q^{+},q)}^{(1\,a)}_{q\Bar{q}g} &  = -\,\frac{e\,e_{f}\,g\,q^{+}\,}{2\,(2\pi)^{4}}\, \int_{x,y,z} \int_{0}^{1}\, d\zeta\, \int_{0}^{1-\zeta}\, d\xi\,\ket{q^{\alpha}_{\lambda_{1}}(\zeta q^{+},x)\Bar{q}^{\beta}_{\lambda_{2}}(\eta\,q^{+},y)g^{a}_{m}(\xi q^{+},z)}\, \frac{\zeta}{\sqrt{\xi}}\, \delta^{(2)}(w-c)\,\times \Psi^{im}_{\lambda_{1}\lambda_{2}}[\mathcal{G}](R,Y) \notag \\ & \left[U^{ab}(z)\,V(x)\,t^{b}\,V^{\dagger}(y)-t^{a}\right]_{\alpha\beta}
\end{align}
\item Anti quark emits a gluon after the shock wave
\begin{figure}[h]
    \centering
    \includegraphics[width=0.2\linewidth]{Wave functions/2.jpeg}
    \label{fig:2}
\end{figure}
\begin{align}
    \ket{\gamma}^{(2)}_{q\Bar{q}g} & = \ket{q_{2}\Bar{q}_{2}g_{2}}\,\bra{q_{2}\Bar{q}_{2}g_{2}}\,S\,\ket{q_{1}\Bar{q}_{1}g_{1}}\times \frac{\bra{q_{1}\Bar{q}_{1}g_{1}}H_{\Bar{q}g}\ket{q\Bar{q}}}{E_{\Bar{q}_{1}}+E_{q_{1}}-E_{q_{2}}-E_{\Bar{q}_{2}}-E_{g_{2}}} \times \frac{\bra{q\Bar{q}}H_{\gamma\rightarrow q\Bar{q}}\ket{\gamma}}{E_{q}+E_{\Bar{q}}-E_{\gamma}}
\end{align}
\begin{align}
    \ket{\gamma}^{(2)}_{q\Bar{q}g} & =  \int\, d^{3}k''\, d^{3}p''\, d^{3}r\, d^{3}k'\, d^{3}p'\, d^{3}k\, d^{3}p\,\ket{q^{\alpha''}_{\lambda''_{1}}(k'')\,\Bar{q}^{\beta''}_{\lambda''_{2}}(p'')\,g^{b}_{j}(r)}\frac{\bra{q^{\alpha''}_{\lambda''_{1}}(k'')\,\Bar{q}^{\beta''}_{\lambda''_{2}}(p'')\,g^{b}_{j}(r)}H_{\Bar{q}g}\ket{q^{\alpha'}_{\lambda'_{1}}(k')\,\Bar{q}^{\beta'}_{\lambda'_{2}}(p')}}{E_{q}(k')+E_{\Bar{q}}(p')-E_{q}(k'')\chg{-}E_{\Bar{q}}(p'')-E_{g}(r)}\,\times \notag \\ & \bra{q^{\alpha'}_{\lambda'_{1}}(k')\,\Bar{q}^{\beta'}_{\lambda'_{2}}(p')}S\ket{q^{\alpha}_{\lambda_{1}}(k)\,\Bar{q}^{\beta}_{\lambda_{2}}(p)}\times  \frac{\bra{q^{\alpha}_{\lambda_{1}}(k)\,\Bar{q}^{\beta}_{\lambda_{2}}(p)}H_{\gamma\rightarrow q\Bar{q}}\ket{\gamma_{\lambda}(q)}}{E_{\Bar{q}}(p)+E_{q}(k)-E_{\gamma}(q,Q)}
\label{chg:dena}\end{align}
The only difference between the first and the second wave function component is the energy denominator and the Wilson line structure. The changes are as follows:
$$\frac{1}{E_{q\Bar{q}g}-E_{\gamma}}\rightarrow \frac{1}{E_{q\Bar{q}}-E_{q\Bar{q}g}}$$ 
Equivalently,
$$\frac{1}{\zeta\,(1-\zeta)^{2}(\Tilde{p}^{2}+\omega^{2}_{f})+\xi\,\eta\, (\Tilde{k}^{2}+\epsilon_{f}^{2})}\rightarrow -\frac{1}{\zeta\,(1-\zeta)^{2}\,(\Tilde{p}^{2}+\omega^{2}_{f})}$$
Also, the Fourier transform which we will be using here is the following:
$$\int\, \frac{d^{2}\Tilde{k}}{2\pi}\, e^{i\Tilde{k}\cdot R}\, \frac{\{\Tilde{k}^{j},1\}}{\Tilde{k}^{2}+\epsilon^{2}_{f}} = \Bigg\{\frac{i\,R^{j}}{R}\,\epsilon_{f}\,K_{1}(\epsilon_{f}\,R),K_{0}(\epsilon_{f} R)\Bigg\}$$
and $$\int\, \frac{d^{2}\Tilde{p}}{2\pi}\, e^{-i\Tilde{p}\cdot Y}\,\frac{\{\Tilde{p}^{n},1\}}{(\Tilde{p}^{2}+\omega^{2}_{f})} = \Bigg\{-\frac{i\,Y^{n}}{Y}\,\omega_{f}\,K_{1}(\omega_{f}\,Y),K_{0}(\omega_{f}\,Y)\Bigg\} $$
With the same normalization, we get
\begin{align}
    (1-\zeta)\, \int\, \frac{d^{2}\Tilde{k}}{2\pi}\, \frac{d^{2}\Tilde{p}}{2\pi}\, e^{i\Tilde{k}\cdot R}\,e^{\chg{-}i\Tilde{p}\cdot Y}\, \frac{\Bigg\{\Tilde{k}^{j}\Tilde{p}^{n},\,\Tilde{k}^{j},\Tilde{p}^{n},\,1\,\Bigg\}}{(\Tilde{k}^{2}+\epsilon^{2}_{f})\,\zeta(1-\zeta)^{2}\,(\Tilde{p}^{2} +\omega^{2}_{f})} & = \Bigg\{R^{j}Y^{n}\mathcal{A}^{(11)},\, i\,R^{j}\,\mathcal{A}^{(10)},\, -iY^{n}\,\mathcal{A}^{(01)},\, \mathcal{A}^{(00)}\Bigg\}
\label{chg:pFT2}\end{align}
where: 
$$\mathcal{A}^{(11)} = \frac{\epsilon_{f}\,K_{1}(\epsilon_{f} R)}{\zeta(1-\zeta)\,R}\,\frac{\omega_{f}\,K_{1}(\omega_{f}\,Y)}{Y}$$
$$\mathcal{A}^{(10)} = \frac{\epsilon_{f}\,K_{1}(\epsilon_{f}R)}{\zeta(1-\zeta)\,R}\, K_{0}(\omega_{f}Y)$$
$$\mathcal{A}^{(01)} = \frac{K_{0}(\epsilon_{f}R)}{\zeta\,(1-\zeta)}\,\frac{\omega_{f}\,K_{1}(\omega_{f}Y)}{Y}$$
$$\mathcal{A}^{(00)} = \frac{K_{0}(\epsilon_{f}R)}{\zeta\,(1-\zeta)}\,K_{0}(\omega_{f}Y)$$
The Wilson line will not have any component of gluon because the gluons are not getting scattered as the gluon is emitted after the scattering. 
$$\left[U^{ab}(z)\,V(x)\,t^{b}\,V^{\dagger}(y)-t^{a}\right]_{\alpha\beta}\rightarrow \left[V(x)\,V^{\dagger}(y')t^{a}-t^{a}\right]_{\alpha\beta} $$
Substituting all these changes in eq. \eqref{eq:1}, we get
\begin{align} \label{eq:2}
     \ket{\gamma(w)}^{(2)}_{q\Bar{q}g} & = +\,\frac{e\,e_{f}\,g\,q^{+}\,}{2\,(2\pi)^{4}}\, \int_{x,y,z} \int_{0}^{1}\, d\zeta\, \int_{0}^{1-\zeta}\, d\xi\,\ket{q^{\alpha}_{\lambda_{1}}(\zeta q^{+},x)\Bar{q}^{\beta}_{\lambda_{2}}(\eta\,q^{+},y)g^{a}_{m}(\xi q^{+},z)}\, \frac{\zeta}{\sqrt{\xi}}\, \delta^{(2)}(w-c)\,\times \notag \\ & \Psi^{im}_{\lambda_{1}\lambda_{2}}[\mathcal{A}](R,Y)\times \left[V(x)\,V^{\dagger}(y')\,t^{a}-t^{a}\right]_{\alpha\beta}
\end{align}
where
\begin{align}
  \Psi^{im}_{\lambda_{1}\lambda_{2}}[\mathcal{A}] & = \sum_{\lambda}\, \Bigg[\varphi^{ij}_{\lambda_{1}\lambda}(\zeta)\,\tau^{mn}_{\lambda\lambda_{2}}(\xi,\eta)\,R^{j}\,Y^{n}\,\mathcal{A}^{(11)}+i\,\xi\,\omega_{f}\,\varphi^{ij}_{\lambda_{1}\lambda}(\zeta)\,\mu^{m}_{\lambda\lambda_{2}}\,R^{j}\,\mathcal{A}^{(10)}\notag \\ & -i\,m_{f}\,\mu^{i}_{\lambda_{1}\lambda}\,\tau^{mn}_{\lambda\lambda_{2}}(\xi,\eta)\,Y^{n}\,\mathcal{A}^{(01)} + m_{f}\,\xi\,\omega_{f}\,\mu^{i}_{\lambda_{1}\lambda}\,\mu^{m}_{\lambda\lambda_{2}}\,\mathcal{A}^{(00)}\Bigg]  
\end{align}
\label{chg:ba}\chg{In the massless case Eq. \eqref{eq:2} follows from Eq. \eqref{eq:1} by the formal substitution $(b) = (a)|_{y,z\rightarrow y'}$ of the paper (keeping the prefactor $1/Y^{2}$). With masses this is no longer literally true: the functions $\mathcal{A}^{(\alpha\beta)}$ are different functions from $\mathcal{G}^{(\alpha\beta)}$ (after the shock wave the gluon is emitted from an on-shell antiquark and sees only the dead-cone scale $\omega_{f}$). We keep the label $[y,z\rightarrow y']$ for this term, meaning: Wilson lines at $y'$ and $\mathcal{G}\rightarrow\mathcal{A}$. The short-distance behaviour is the same, $Y^{2}\mathcal{G}^{(11)},\,Y^{2}\mathcal{A}^{(11)}\xrightarrow{Y\rightarrow 0}\frac{\epsilon_{f}K_{1}(\epsilon_{f}R)}{\zeta(1-\zeta)R}$, $Y^{2}\mathcal{G}^{(01)},\,Y^{2}\mathcal{A}^{(01)}\xrightarrow{Y\rightarrow 0}\frac{K_{0}(\epsilon_{f}R)}{\zeta(1-\zeta)}$, and $\mathcal{G}^{(10)}-\mathcal{A}^{(10)}$, $\mathcal{G}^{(00)}-\mathcal{A}^{(00)}$ stay finite, so the UV cancellation between Eq. \eqref{eq:1} and Eq. \eqref{eq:2} works as in the massless case.}
%and $\mathcal{D}(x,y',y',\zeta,\xi) = \sqrt{\zeta\,(1-\zeta)}\,|x-y'| = \frac{\Tilde{Q}\,R}{Q}$
\item Instantaneous contribution (Anti quark emission of gluon when an instantaneous vertex acts before the shockwave)
\begin{figure}[h]
    \centering
    \includegraphics[width=0.2\linewidth]{Wave functions/5.png}
    \label{fig:5}
\end{figure}
\begin{align}
\ket{\gamma}^{(1\,\text{inst.})}_{q\Bar{q}g} & =  \ket{q\Bar{q}g}\,\bra{q\Bar{q}g}\,S\,\ket{q\Bar{q}g}\,\frac{\bra{q\Bar{q}g}\, H_{\gamma\rightarrow q\Bar{q}g}\ket{\gamma}}{E_{q\Bar{q}g}-E_{\gamma}}
\end{align}
Using the relevant matrix  we get (In momentum space - without Wilson lines)
\begin{align}
\ket{\gamma}^{(1\,\text{inst.})}_{q\Bar{q}g} & = \frac{1}{(2\pi)^{6}}\, \int\, d^{2}\Tilde{k}\, d^{2}\Tilde{p}\, \int_{0}^{1}\, d\zeta\, \int_{0}^{1-\zeta}\, d\xi\,\ket{\Bar{q}^{\beta}_{\lambda_{2}}\,(\eta\,q^{+},q-k-p),g^{a}_{m}(\xi q^{+},p),q^{\alpha}_{\lambda_{1}}(\zeta q^{+},k)}\, \frac{e\,e_{f}\,g\,q^{+}\,\zeta}{2\sqrt{\xi}}\notag \\ & \frac{(\delta^{im} +2i\,\epsilon^{im}\,\lambda_{1})\,\delta_{\lambda_{1}\lambda_{2}}\,t^{a}_{\alpha\beta}}{\Tilde{k}^{2}+\epsilon^{2}_{f}+\frac{\zeta\,(1-\zeta)^{2}}{\xi\,\eta}\,(\Tilde{p}^{2} +\omega^{2}_{f})}
\end{align}
Performing Fourier transform and inserting $S-1$ component, we get the following
\begin{align}
    (1-\zeta)\, \int\, \frac{d^{2}\Tilde{k}}{2\pi}\, \frac{d^{2}\Tilde{p}}{2\pi}\, e^{i\,\Tilde{k}\cdot R}\, e^{-i\Tilde{p}\cdot Y}\, \frac{1}{\zeta\, (1-\zeta)^{2}\,(\Tilde{p}^{2}+\omega^{2}_{f}) + \xi\,\eta\,(\Tilde{k}^{2}+\epsilon^{2}_{f})} = \frac{E_{f}\,K_{1}(E_{f}\mathcal{D})}{\mathcal{D}} = \mathcal{G}_{E}
\end{align}
The outgoing component of wave function is obtained as follows:  
\begin{align}
\ket{\gamma}^{(1\,\text{inst.})}_{q\Bar{q}g} & = +\frac{e\,e_{f}\,g\,q^{+}}{2\,(2\pi)^{4}}\, \int_{x,y,z}\, \int_{0}^{1}\, d\zeta\, \int_{0}^{1-\zeta}\, d\xi\, \ket{\Bar{q}^{\beta}_{\lambda_{2}}(\eta\,q^{+},y),g^{a}_{m}(\xi\,q^{+},z),q^{\alpha}_{\lambda_{1}}(\zeta\,q^{+},x)}\, \frac{\zeta\,\eta\, \sqrt{\xi}}{(1-\zeta)}\, \delta^{(2)}(w-c)\, \left(\delta^{im}+2i\,\epsilon^{im}\lambda_{1}\right)\,\delta_{\lambda_{1}\lambda_{2}}\notag \\ & \frac{E_{f}\,K_{1}(E_{f}\,\mathcal{D})}{\mathcal{D}} \left[U^{ab}(z)V(x)t^{b}V^{\dagger}(y)-t^{a}\right]_{\alpha\beta}
\end{align}
This has the same Wilson line structure as the regular i.e., Eq. \eqref{eq:1}. \label{chg:Y2RY}\chg{There is no instantaneous term after the shock wave (the gluon emission after the shock wave is a single vertex), so the instantaneous term has no $[y,z\rightarrow y']$ counterpart (in the massless note this was expressed by the factor $\frac{Y^{2}}{R\cdot Y}$ of the effective vertex, which vanishes for $y,z\rightarrow y'$).} Now, we have all the wave functions corresponding to anti quark emission. All four anti quark emission can be written as a single expression with an effective vertex. The effective vertex is written as
\begin{align}
\mathcal{M}^{im}_{1\,\lambda_{1}\lambda_{2}} (x,y,z,\zeta,\xi) & = \Psi^{im}_{\lambda_{1}\lambda_{2}}[\mathcal{G}](R,Y) - \delta_{\lambda_{1}\lambda_{2}}\,\frac{\xi\,\eta}{(1-\zeta)}\,\left(\delta^{im}+2i\epsilon^{im}\,\lambda_{1}\right)\,\mathcal{G}_{E}
\end{align}
\begin{align}
\mathcal{M}^{im}_{2\,\lambda_{1}\lambda_{2}} (x,y,z,\zeta,\xi) & = \Psi^{im}_{\lambda_{1}\lambda_{2}}[\mathcal{A}](R,Y)
\end{align}
All four anti quark emission are then written using a single expression: 
\begin{align}
    \ket{\gamma^{i}(w)}_{q\Bar{q}g} & = -\frac{e\,e_{f}\,g\,q^{+}}{2\,(2\pi)^{4}}\,\int_{x,y,z}\, \int_{0}^{1}\, d\zeta\, \int_{0}^{1-\zeta}\,d\xi\, \ket{\Bar{q}^{\beta}_{\lambda_{2}}(\eta\,q^{+},y),g^{a}_{m}(\xi\,q^{+},z),q^{\alpha}_{\lambda_{1}}(\zeta\,q^{+},x)}\,\frac{\zeta}{\sqrt{\xi}}\,\delta^{(2)}(w-c)\notag \\ & \bigg[\mathcal{M}^{im}_{1\lambda_{1}\lambda_{2}}(x,y,z,\zeta,\xi)\left[U^{ab}(z)V(x)t^{b}V^{\dagger}(y)-t^{a}\right]_{\alpha\beta} - \mathcal{M}^{im}_{2\lambda_{1}\lambda_{2}}\,\left[V(x)\,V^{\dagger}(y')t^{a}-t^{a}\right]_{\alpha\beta}\bigg]
\end{align}
\item Quark emits a gluon before the shock wave 
\begin{figure}[h]
    \centering
    \includegraphics[width=0.2\linewidth]{Wave functions/3.jpeg}
    \label{fig:3}
\end{figure}\\ 

When we move from anti quark emission to quark emission, we will have to perform the following changes from anti quark wave functions:
\begin{itemize}
    \item $\zeta \leftrightarrow \eta$
    \item $x\leftrightarrow y$
    \item $\alpha \leftrightarrow \beta$
    \item $\lambda_{1}\leftrightarrow \lambda_{2}$
    \item Change overall sign
    \item $y'\rightarrow x' = \frac{\zeta\,x + \xi\, z}{\zeta + \xi}$
    \item $R \rightarrow y-x' = -\Tilde{R}$
    \item $Y \rightarrow x-z  = X$
    \item $\varphi^{ij}_{\lambda_{1}\lambda}(\zeta) \rightarrow \varphi^{ij\,*}_{\lambda_{2}\lambda}(\eta) = -\,\varphi^{ij}_{\lambda\lambda_{2}}(\zeta+\xi)$
    \item $\tau^{mn}_{\lambda\lambda_{2}}(\xi,\eta) \rightarrow \tau^{mn\,*}_{\lambda\,\lambda_{1}}(\xi,\zeta) = \tau^{nm}_{\lambda_{1}\lambda}(\xi,\zeta)$
    \item $D(x,y,z,\zeta, \xi) \rightarrow D(y,x,z,\eta,\xi)$
    \label{chg:Drule}\item \chg{after the shock wave (massless limit): $\mathcal{D}\rightarrow\sqrt{\eta(1-\eta)}\,\Tilde{R}$, i.e.\ $\Tilde{Q}\rightarrow\sqrt{\eta(1-\eta)}\,Q$}
    \item $m_{f}\,\mu^{i}_{\lambda_{1}\lambda} \rightarrow m_{f}\,\mu^{i}_{\lambda\lambda_{2}} $
    \item $\xi\, \omega_{f}\,\mu^{m}_{\lambda\lambda_{2}} \rightarrow \xi\, \Tilde{\omega}_{f}\,\mu^{m}_{\lambda_{1}\lambda}$ with $\Tilde{\omega}_{f} = \frac{\xi\,m_{f}}{(1-\eta)}$
    \item $\epsilon^{2}_{f} \rightarrow \Tilde{\epsilon}^{2}_{f} = \eta(1-\eta)\, Q^{2} + m^{2}_{f}$, $a \rightarrow \Tilde{a} = \frac{\xi\,\zeta}{\eta\, (1-\eta)^{2}}$
    \item $A_{x} \rightarrow \Tilde{A}_{x} = \eta\,(1-\eta)\,Q^{2} + m^{2}_{f}\left(1+\frac{x\eta\xi}{\zeta}\right)$ and $\Delta_{x} \rightarrow \Tilde{\Delta}_{x} = \Tilde{R}^{2} + \frac{\Tilde{a}\,X^{2}}{x}$
    \item $\mathcal{G}^{(\alpha\beta)}(R,Y) \rightarrow \mathcal{\Tilde{G}}^{(\alpha\beta)}(\Tilde{R},X) $ and $\mathcal{A}^{(\alpha,\beta)}(R,Y) \rightarrow \mathcal{\Tilde{A}}^{(\alpha\beta)}(\Tilde{R},X)$ also, $\zeta\,(1-\zeta) \rightarrow \eta\,(1-\eta)$
\end{itemize}
\begin{align}
    \ket{\gamma}^{(3)}_{q\Bar{q}g} & = \ket{q_{2}\Bar{q}_{2}g_{2}}\,\bra{q_{2}\Bar{q}_{2}g_{2}}\,S\,\ket{q_{1}\Bar{q}_{1}g_{1}}\times \frac{\bra{q_{1}\Bar{q}_{1}g_{1}}H_{qg}\ket{q\Bar{q}}}{E_{\Bar{q}_{1}}+E_{q_{1}}+E_{g_{1}}-E_{\gamma}} \times \frac{\bra{q\Bar{q}}H_{\gamma\rightarrow q\Bar{q}}\ket{\gamma}}{E_{q}+E_{\Bar{q}}-E_{\gamma}}
\end{align}
\begin{align}
    \ket{\gamma}^{(3)}_{q\Bar{q}g} & = \int\, d^{3}k''\, d^{3}p''\, d^{3}r\,d^{3}k'\, d^{3}p'\, d^{3}q'\, d^{3}k\, d^{3}p\, \ket{q^{\alpha''}_{\lambda''_{1}}(k'')\,\Bar{q}^{\beta''}_{\lambda''_{2}}(p'')\,g^{b}_{j}(r)}\,\bra{q^{\alpha''}_{\lambda''_{1}}(k'')\,\Bar{q}^{\beta''}_{\lambda''_{2}}(p'')\,g^{b}_{j}(r)}\,S\,\ket{q^{\alpha'}_{\lambda'_{1}}(k')\,\Bar{q}^{\beta'}_{\lambda'_{2}}(p')\,g^{a'}_{i'}(q')}\times\notag \\ & \frac{\bra{q^{\alpha'}_{\lambda'_{1}}(k')\,\Bar{q}^{\beta'}_{\lambda'_{2}}(p')\,g^{a'}_{i'}(q')}H_{qg}\ket{q^{\alpha}_{\lambda_{1}}(k)\,\Bar{q}^{\beta}_{\lambda_{2}}(p)}}{E_{\Bar{q}}(p')+E_{q}(k')+E_{g}(q')-E_{\gamma}} \times \frac{\bra{q^{\alpha}_{\lambda_{1}}(k)\,\Bar{q}^{\beta}_{\lambda_{2}}(p)}H_{\gamma\rightarrow q\Bar{q}}\ket{\gamma_{\lambda}(q)}}{E_{\Bar{q}}(p)+E_{q}(k)-E_{\gamma}(q,Q)}
\end{align} 
\begin{align} \label{eq:5}
     \ket{\gamma^{i}(w)}^{(3)}_{q\Bar{q}g} & = \,\frac{e\,e_{f}\,g\,q^{+}\,}{2\,(2\pi)^{4}}\, \int_{x,y,z} \int_{0}^{1}\, d\zeta\, \int_{0}^{1-\zeta}\, d\xi\,\ket{q^{\alpha}_{\lambda_{1}}(\zeta q^{+},x)\Bar{q}^{\beta}_{\lambda_{2}}(\eta\,q^{+},y)g^{a}_{m}(\xi q^{+},z)}\, \frac{\eta}{\sqrt{\xi}}\, \delta^{(2)}(w-c)\,\times \notag \\ & \bigg[\Tilde{\Psi}^{im}_{\lambda_{1}\lambda_{2}}[\mathcal{\Tilde{G}}](\Tilde{R},X)\times \left[U^{ab}(z)\,V(x)\,t^{b}\,V^{\dagger}(y)-t^{a}\right]_{\alpha\beta}\bigg]
\end{align}
with 
\begin{align}
    \Tilde{\Psi}^{im}_{\lambda_{1}\lambda_{2}}[\mathcal{\Tilde{G}}] & = \sum_{\lambda}\, \Bigg[\tau^{nm}_{\lambda_{1}\lambda}(\xi,\zeta)\varphi^{ij}_{\lambda\lambda_{2}}(\zeta + \xi)\,\Tilde{R}^{j}X^{n}\,\mathcal{\Tilde{G}}^{(11)}+i\xi\,\Tilde{\omega}_{f}\mu^{m}_{\lambda_{1}\lambda}\varphi^{ij}_{\lambda\lambda_{2}}(\zeta + \xi)\Tilde{R}^{j}\mathcal{\Tilde{G}}^{(10)}\notag \\ & -im_{f}\,\tau^{nm}_{\lambda_{1}\lambda}(\xi,\zeta)\,\mu^{i}_{\lambda\lambda_{2}}\,X^{n}\mathcal{\Tilde{G}}^{(01)} + m_{f}\xi\,\Tilde{\omega}_{f}\mu^{m}_{\lambda_{1}\lambda}\mu^{i}_{\lambda\lambda_{2}}\mathcal{\Tilde{G}}^{(00)}\Bigg]
\end{align}
where
$$\mathcal{\Tilde{G}}^{(\alpha\beta)}(\Tilde{R},X) = \frac{\Tilde{a}^{\beta}}{2\,\eta(1-\eta)}\,\int_{0}^{1}\, \frac{dx}{x^{1+\beta}}\, \left(\frac{\Tilde{A}_{x}}{\Tilde{\Delta}_{x}}\right)^{\frac{(\alpha + \beta)}{2}}\,K_{\alpha + \beta} (\sqrt{\Tilde{A_{x}}\Tilde{\Delta_{x}}}), \alpha, \beta \in \{0,1\}$$
$$\Tilde{A}_{x} = \Tilde{\epsilon}^{2}_{f} + \frac{x\,\Tilde{\omega}^{2}_{f}}{\Tilde{a}} = \eta\,(1-\eta)\,Q^{2} +m^{2}_{f}\,\left(1+\frac{x\eta\xi}{\zeta}\right)$$
$$\Tilde{\Delta}_{x} = \Tilde{R}^{2} + \frac{\Tilde{a}\,X^{2}}{x}$$
Explicitly
\begin{align}
    \mathcal{\Tilde{G}}^{(11)} & = \frac{\Tilde{a}}{2\,\eta\, (1-\eta)}\, \int_{0}^{1}\, \frac{dx}{x^{2}}\, \frac{\Tilde{A}_{x}}{\Tilde{\Delta}_{x}}\,K_{2}\left(\sqrt{\Tilde{A}_{x}\Tilde{\Delta}_{x}}\right)
\end{align}
\begin{align}
    \mathcal{\Tilde{G}}^{(10)} & = \frac{1}{2\, \eta\, (1-\eta)}\, \int_{0}^{1}\, \frac{dx}{x}\, \sqrt{\frac{\Tilde{A}_{x}}{\Tilde{\Delta}_{x}}}\, K_{1}\left(\sqrt{\Tilde{A}_{x}\Tilde{\Delta}_{x}}\right)
\end{align}
\begin{align}
     \mathcal{\Tilde{G}}^{(01)} & = \frac{\Tilde{a}}{2\, \eta\, (1-\eta)}\, \int_{0}^{1}\, \frac{dx}{x^{2}}\, \sqrt{\frac{\Tilde{A}_{x}}{\Tilde{\Delta}_{x}}}\, K_{1}\left(\sqrt{\Tilde{A}_{x}\,\Tilde{\Delta}_{x}}\right)
\end{align}
\begin{align}
     \mathcal{\Tilde{G}}^{(00)} & = \frac{1}{2\, \eta\, (1-\eta)}\, \int_{0}^{1}\, \frac{dx}{x}\, K_{0}\left(\sqrt{\Tilde{A}_{x}\Tilde{\Delta}_{x}}\right)
\end{align}
\item Quark emits a gluon after the shock wave 
\begin{figure}[h]
    \centering
    \includegraphics[width=0.2\linewidth]{Wave functions/4.jpeg}
    \label{fig:4}
\end{figure}
\begin{align}
    \ket{\gamma}^{(4)}_{q\Bar{q}g} & = \ket{q_{2}\Bar{q}_{2}g_{2}}\,\bra{q_{2}\Bar{q}_{2}g_{2}}\,S\,\ket{q_{1}\Bar{q}_{1}g_{1}}\times \frac{\bra{q_{1}\Bar{q}_{1}g_{1}}H_{qg}\ket{q\Bar{q}}}{E_{\Bar{q}_{1}}+E_{q_{1}}-E_{q_{2}}-E_{\Bar{q}_{2}}-E_{g_{2}}} \times \frac{\bra{q\Bar{q}}H_{\gamma\rightarrow q\Bar{q}}\ket{\gamma}}{E_{q}+E_{\Bar{q}}-E_{\gamma}}
\end{align}
\begin{align}
    \ket{\gamma}^{(4)}_{q\Bar{q}g} & =  \int\, d^{3}k''\, d^{3}p''\, d^{3}r\, d^{3}k'\, d^{3}p'\, d^{3}k\, d^{3}p\,\ket{q^{\alpha''}_{\lambda''_{1}}(k'')\,\Bar{q}^{\beta''}_{\lambda''_{2}}(p'')\,g^{b}_{j}(r)}\frac{\bra{q^{\alpha''}_{\lambda''_{1}}(k'')\,\Bar{q}^{\beta''}_{\lambda''_{2}}(p'')\,g^{b}_{j}(r)}H_{qg}\ket{q^{\alpha'}_{\lambda'_{1}}(k')\,\Bar{q}^{\beta'}_{\lambda'_{2}}(p')}}{E_{q}(k')+E_{\Bar{q}}(p')-E_{q}(k'')\chg{-}E_{\Bar{q}}(p'')-E_{g}(r)}\,\times \notag \\ & \bra{q^{\alpha'}_{\lambda'_{1}}(k')\,\Bar{q}^{\beta'}_{\lambda'_{2}}(p')}S\ket{q^{\alpha}_{\lambda_{1}}(k)\,\Bar{q}^{\beta}_{\lambda_{2}}(p)}\times  \frac{\bra{q^{\alpha}_{\lambda_{1}}(k)\,\Bar{q}^{\beta}_{\lambda_{2}}(p)}H_{\gamma\rightarrow q\Bar{q}}\ket{\gamma_{\lambda}(q)}}{E_{\Bar{q}}(p)+E_{q}(k)-E_{\gamma}(q,Q)}
\label{chg:denb}\end{align}
Substituting the relations we get
\begin{align}
    \ket{\gamma^{i}(w)}^{(4)}_{q\Bar{q}g} & = -\,\frac{e\,e_{f}\,g\,q^{+}\,}{2\,(2\pi)^{4}}\, \int_{x,y,z} \int_{0}^{1}\, d\zeta\, \int_{0}^{1-\zeta}\, d\xi\,\ket{q^{\alpha}_{\lambda_{1}}(\zeta q^{+},x)\Bar{q}^{\beta}_{\lambda_{2}}(\eta\,q^{+},y)g^{a}_{m}(\xi q^{+},z)}\, \frac{\eta}{\sqrt{\xi}}\, \delta^{(2)}(w-c)\,\times \notag \\ & \Tilde{\Psi}^{im}_{\lambda_{1}\lambda_{2}}[\chg{\mathcal{\Tilde{A}}}](\Tilde{R},X)\times \left[t^{a}V(x')V^{\dagger}(y)-t^a\right]_{\alpha\beta}
\label{chg:PsiAq}\end{align}
where 
\begin{align}
    \mathcal{\Tilde{A}}^{(11)} & = \frac{\Tilde{\epsilon}_{f}\,K_{1}(\Tilde{\epsilon}_{f}\,\Tilde{R})}{\eta\,(1-\eta)\,\Tilde{R}}\, \frac{\Tilde{\omega}_{f}\,K_{1}(\Tilde{\omega}_{f}\,X)}{X}
\end{align}
\begin{align}
    \mathcal{\Tilde{A}}^{(10)} & = \frac{\Tilde{\epsilon}_{f}\,K_{1}(\Tilde{\epsilon}_{f}\Tilde{R})}{\eta\,(1-\eta)\,\Tilde{R}}\, K_{0}(\Tilde{\omega}_{f}\,X)
\end{align}
\begin{align}
    \mathcal{\Tilde{A}}^{(01)} & =   \frac{K_{0}(\Tilde{\epsilon}_{f}\Tilde{R})}{\eta\,(1-\eta)}\,\frac{\Tilde{\omega}_{f}\,K_{1}(\Tilde{\omega}_{f}\,X)}{X}
\end{align}
\begin{align}
    \mathcal{\Tilde{A}}^{(00)} & = \frac{K_{0}(\Tilde{\epsilon}_{f}\,\Tilde{R})}{\eta\,(1-\eta)}\, K_{0}(\Tilde{\omega}_{f}\,X)
\end{align}
\item Instantaneous contribution (Quark emission of gluon when an instantaneous vertex acts before the shockwave)
\begin{figure}[h]
    \centering
    \includegraphics[width=0.2\linewidth]{Wave functions/6.png}
    \label{fig:6}
\end{figure}
\begin{align}
\ket{\gamma}^{(3\,\text{inst.})}_{q\Bar{q}g} & =  \ket{q\Bar{q}g}\,\bra{q\Bar{q}g}\,S\,\ket{q\Bar{q}g}\,\frac{\bra{q\Bar{q}g}\, H_{\gamma\rightarrow q\Bar{q}g}\ket{\gamma}}{E_{q\Bar{q}g}-E_{\gamma}}
\end{align}
Using the transformations from anti quark to quark with the insertions of Wilson lines, we get
\begin{align}
\ket{\gamma}^{(3\,\text{inst.})}_{q\Bar{q}g} & = -\frac{e\,e_{f}\,g\,q^{+}}{2\,(2\pi)^{4}}\, \int_{x,y,z}\, \int_{0}^{1}\, d\zeta\, \int_{0}^{1-\zeta}\, d\xi\, \ket{\Bar{q}^{\beta}_{\lambda_{2}}(\eta\,q^{+},y),g^{a}_{m}(\xi\,q^{+},z),q^{\alpha}_{\lambda_{1}}(\zeta\,q^{+},x)}\, \frac{\zeta\,\eta\, \sqrt{\xi}}{(1-\eta)}\, \delta^{(2)}(w-c)\, \left(\delta^{im}-2i\,\epsilon^{im}\lambda_{2}\right)\,\delta_{\lambda_{1}\lambda_{2}}\notag \\ & \frac{E_{f}\,K_{1}(E_{f}\mathcal{D})}{\mathcal{D}}\,\left[U^{ab}(z)V(x)t^{b}V^{\dagger}(y)-t^{a}\right]_{\alpha\beta}
\end{align}
Now, we will write the quark emissions as how we wrote the anti quark section. We will use the effective vertex formalism.
\begin{align} 
     \ket{\gamma^{i}(w)}_{q\Bar{q}g} & = \,\frac{e\,e_{f}\,g\,q^{+}\,}{2\,(2\pi)^{4}}\, \int_{x,y,z} \int_{0}^{1}\, d\zeta\, \int_{0}^{1-\zeta}\, d\xi\,\ket{q^{\alpha}_{\lambda_{1}}(\zeta q^{+},x)\Bar{q}^{\beta}_{\lambda_{2}}(\eta\,q^{+},y)g^{a}_{m}(\xi q^{+},z)}\, \frac{\eta}{\sqrt{\xi}}\, \delta^{(2)}(w-c)\,\times \notag \\ & \Bigg[\mathcal{\Tilde{M}}^{im}_{1\lambda_{1}\lambda_{2}}(x,y,z,\eta,\xi) \times\, \left[U^{ab}(z)\,V(x)\,t^{b}\,V^{\dagger}(y)-t^{a}\right]_{\alpha\beta}-\mathcal{\Tilde{M}}^{im}_{2\lambda_{1}\lambda_{2}}\left[t^{a}\,V(x')\,V^{\dagger}(y) - t^{a}\right]_{\alpha\beta}\Bigg]
\end{align}
with the effective vertex
\begin{align}
\mathcal{\Tilde{M}}^{im}_{1\lambda_{1}\lambda_{2}}(x,y,z,\eta,\xi) &  = \Tilde{\Psi}^{im}_{\lambda_{1}\lambda_{2}}[\mathcal{\Tilde{G}}](\Tilde{R},X) - \delta_{\lambda_{1}\lambda_{2}}\,\frac{\zeta\,\xi}{(\xi + \zeta)}\,(\delta^{im}-2i\,\epsilon^{im}\lambda_{1})\,\mathcal{G}_{E} 
\end{align}
\begin{align}
\mathcal{\Tilde{M}}^{im}_{2\lambda_{1}\lambda_{2}}(x,y,z,\eta,\xi) &  = \Tilde{\Psi}^{im}_{\lambda_{1}\lambda_{2}}[\mathcal{\Tilde{A}}] (\Tilde{R},X)
\end{align}
\end{enumerate}
\section{Trijet production cross section}
\subsection{Direct contributions - Anti quark contribution}
\begin{align}
     \frac{d^{3}N^{\text{direct}}}{d^{3}k_{1}\,d^{3}k_{2}\, d^{3}k_{3}} & = \frac{e^{2}\,g^{2}\,\zeta^{2}\,C_{F}\,N_{c}}{8\, (2\pi)^{12}\,\xi\,q^{+\,2}}\,  \delta\left(q^{+}-k^{+}_{1}-k^{+}_{2}-k^{+}_{3}\right)\, \int_{x,y,z,\Bar{x},\Bar{y},\Bar{z}}\, e^{-i\,k_{1}\cdot (x-\Bar{x})}\, e^{-i\,k_{2}\cdot (y-\Bar{y})}\, e^{-i\,k_{3}\cdot (z-\Bar{z})}\,\left(\sum_{f}\, e^{2}_{f}\right)\,\notag \\ & \Bigg[\mathcal{H}[\mathcal{G};\mathcal{\Bar{G}}]\,\mathcal{W}(x,y,z,\Bar{x},\Bar{y},\Bar{z})-\mathcal{H}[\mathcal{A};\mathcal{\Bar{G}}]\,\mathcal{W}(x,y',y',\Bar{x},\Bar{y},\Bar{z})-\mathcal{H}[\mathcal{G};\mathcal{\Bar{A}}]\,\mathcal{W}(x,y,z,\Bar{x},\Bar{y}',\Bar{y}')+\mathcal{H}[\mathcal{A};\mathcal{\Bar{A}}]\,\mathcal{W}(x,y',y',\Bar{x},\Bar{y}',\Bar{y}')\Bigg]
\end{align}
where
\begin{align}
    \mathcal{H}[F;\Bar{F}] & = \sum_{\lambda_{1},\lambda_{2}}\,\sum_{l,r}\, \mathcal{M}^{lr}_{\lambda_{1}\lambda_{2}}[F](x,y,z)\,\mathcal{M}^{lr\,*}_{\lambda_{1}\lambda_{2}}[\Bar{F}](\Bar{x},\Bar{y},\Bar{z})
\end{align}
\begin{itemize}
    \item Regular $\times$ Regular
    \begin{align}
        \mathcal{H}_{\mathrm{reg.}\times \mathrm{reg.}}[F;\Bar{F}] & = R^{i}\,Y^{m}\,\Bar{R}^{j}\,\Bar{Y}^{n}\,\Bigg\{8\,\left[\eta^{2}+(1-\zeta)^{2}\right]\,\left[\zeta^{2} + (1-\zeta)^{2}\right]\,\delta^{ij}\delta^{mn}-8\xi(2\zeta - 1)\, (2(1-\zeta)-\xi)\epsilon^{ij}\,\epsilon^{mn}\Bigg\}F^{(11)}\,\Bar{F}^{(11)}\notag \\ & +8\,m^{2}_{f}\,\left[\eta^{2} + (1-\zeta)^{2}\right]\, (Y\cdot \Bar{Y})\,F^{(01)}\,\Bar{F}^{(01)}+8\,\xi^{2}\,\omega^{2}_{f}\,\left[\zeta^{2}+(1-\zeta)^{2}\right]\, (R\cdot \Bar{R})\, F^{(10)}\,\Bar{F}^{(10)}\notag \\ & + 8\,m^{2}_{f}\,\xi^{2}\,\omega^{2}_{f}\,F^{(00)}\,\Bar{F}^{(00)}\notag \\ & -8\zeta\eta\, m_{f}\,\xi\,\omega_{f}\, \left[(R\cdot Y)\,F^{(11)}\,\Bar{F}^{(00)} + (\Bar{R}\cdot \Bar{Y})\,F^{(00)}\,\Bar{F}^{(11)} + (R\cdot \Bar{Y})\,F^{(10)}\,\Bar{F}^{(01)} + (Y\cdot \Bar{R})\,F^{(01)}\,\Bar{F}^{(10)}\right]
    \end{align}
    \item Regular $\times$ Instantaneous $+$ h.c.
    \begin{align}
        \mathcal{H}_{\mathrm{reg.}\times \mathrm{inst.}}[F;\Bar{F}] & = -\frac{8\zeta\xi\eta^{2}}{(1-\zeta)}\, \left[(R\cdot Y)\,F^{(11)}\,\mathcal{\Bar{G}}_{E} + (\Bar{R}\cdot \Bar{Y})\mathcal{G}_{E}\,\Bar{F}^{(11)}\right]\notag \\ & +\frac{8\xi^{2}\,\eta\,m_{f}\,\omega_{f}}{(1-\zeta)}\,\left[F^{(00)}\,\mathcal{\Bar{G}}_{E} + \mathcal{G}_{E}\,\Bar{F}^{(00)}\right]
    \end{align}
    \item Instantaneous $\times$ Instantaneous
    \begin{align}
       \mathcal{H}_{\mathrm{inst.}\times \mathrm{inst.}}[F;\Bar{F}] & = 8\,\left[\frac{\xi\,\eta}{(1-\zeta)}\right]^{2}\, \mathcal{G}_{E}\, \mathcal{\Bar{G}}_{E}
    \end{align}
\end{itemize}
where: 
\begin{align}
\mathcal{M}^{im}_{\lambda_{1}\lambda_{2}} [\mathcal{G}] & = \mathcal{M}^{im}_{1\,\lambda_{1}\lambda_{2}} [\mathcal{G}] = \Psi^{im}_{\lambda_{1}\lambda_{2}}[\mathcal{G}](R,Y) - \delta_{\lambda_{1}\lambda_{2}}\,\frac{\xi\,\eta}{(1-\zeta)}\,\left(\delta^{im}+2i\epsilon^{im}\,\lambda_{1}\right)\,\mathcal{G}_{E}
\end{align}
\begin{align}
\mathcal{M}^{im}_{\lambda_{1}\lambda_{2}} [\mathcal{A}] & = \mathcal{M}^{im}_{2\lambda_{1}\lambda_{2}} [\mathcal{A}] = \Psi^{im}_{\lambda_{1}\lambda_{2}}[\mathcal{A}](R,Y)
\end{align}
\begin{align}
  \Psi^{im}_{\lambda_{1}\lambda_{2}}[\mathcal{G}] & = \sum_{\lambda}\, \Bigg[\varphi^{ij}_{\lambda_{1}\lambda}(\zeta)\,\tau^{mn}_{\lambda\lambda_{2}}(\xi,\eta)\,R^{j}\,Y^{n}\,\mathcal{G}^{(11)}+i\,\xi\,\omega_{f}\,\varphi^{ij}_{\lambda_{1}\lambda}(\zeta)\,\mu^{m}_{\lambda\lambda_{2}}\,R^{j}\,\mathcal{G}^{(10)}\notag \\ & -i\,m_{f}\,\mu^{i}_{\lambda_{1}\lambda}\,\tau^{mn}_{\lambda\lambda_{2}}(\xi,\eta)\,Y^{n}\,\mathcal{G}^{(01)} + m_{f}\,\xi\,\omega_{f}\,\mu^{i}_{\lambda_{1}\lambda}\,\mu^{m}_{\lambda\lambda_{2}}\,\mathcal{G}^{(00)}\Bigg]  
\end{align}
Explicitly: 
$$\mathcal{G}^{(\alpha\beta)}(R,Y) = \frac{a^{\beta}}{2\,\zeta(1-\zeta)}\,\int_{0}^{1}\, \frac{dx}{x^{1+\beta}}\, \left(\frac{A_{x}}{\Delta_{x}}\right)^{\frac{(\alpha + \beta)}{2}}\,K_{\alpha + \beta} (\sqrt{A_{x}\Delta_{x}}), \alpha, \beta \in \{0,1\}$$
$$A_{x} = \epsilon^{2}_{f} + \frac{x\,\omega^{2}_{f}}{a} = \zeta\,(1-\zeta)\,Q^{2} +m^{2}_{f}\,\left(1+\frac{x\zeta\xi}{\eta}\right)$$
$$\Delta_{x} = R^{2} + \frac{a\,Y^{2}}{x}$$
\begin{align}
    \mathcal{G}^{(11)} & = \frac{a}{2\,\zeta\, (1-\zeta)}\, \int_{0}^{1}\, \frac{dx}{x^{2}}\, \frac{A_{x}}{\Delta_{x}}\,K_{2}\left(\sqrt{A_{x}\Delta_{x}}\right)
\end{align}
\begin{align}
    \mathcal{G}^{(10)} & = \frac{1}{2\, \zeta\, (1-\zeta)}\, \int_{0}^{1}\, \frac{dx}{x}\, \sqrt{\frac{A_{x}}{\Delta_{x}}}\, K_{1}\left(\sqrt{A_{x}\Delta_{x}}\right)
\end{align}
\begin{align}
     \mathcal{G}^{(01)} & = \frac{a}{2\, \zeta\, (1-\zeta)}\, \int_{0}^{1}\, \frac{dx}{x^{2}}\, \sqrt{\frac{A_{x}}{\Delta_{x}}}\, K_{1}\left(\sqrt{A_{x}\,\Delta_{x}}\right)
\end{align}
\begin{align}
     \mathcal{G}^{(00)} & = \frac{1}{2\, \zeta\, (1-\zeta)}\, \int_{0}^{1}\, \frac{dx}{x}\,K_{0}\left(\sqrt{A_{x}\Delta_{x}}\right)
\end{align}
and
\begin{align}
  \Psi^{im}_{\lambda_{1}\lambda_{2}}[\mathcal{A}] & = \sum_{\lambda}\, \Bigg[\varphi^{ij}_{\lambda_{1}\lambda}(\zeta)\,\tau^{mn}_{\lambda\lambda_{2}}(\xi,\eta)\,R^{j}\,Y^{n}\,\mathcal{A}^{(11)}+i\,\xi\,\omega_{f}\,\varphi^{ij}_{\lambda_{1}\lambda}(\zeta)\,\mu^{m}_{\lambda\lambda_{2}}\,R^{j}\,\mathcal{A}^{(10)}\notag \\ & -i\,m_{f}\,\mu^{i}_{\lambda_{1}\lambda}\,\tau^{mn}_{\lambda\lambda_{2}}(\xi,\eta)\,Y^{n}\,\mathcal{A}^{(01)} + m_{f}\,\xi\,\omega_{f}\,\mu^{i}_{\lambda_{1}\lambda}\,\mu^{m}_{\lambda\lambda_{2}}\,\mathcal{A}^{(00)}\Bigg]  
\end{align}
Explicitly: 
$$\mathcal{A}^{(11)} = \frac{\epsilon_{f}\,K_{1}(\epsilon_{f} R)}{\zeta(1-\zeta)\,R}\,\frac{\omega_{f}\,K_{1}(\omega_{f}\,Y)}{Y}$$
$$\mathcal{A}^{(10)} = \frac{\epsilon_{f}\,K_{1}(\epsilon_{f}R)}{\zeta(1-\zeta)\,R}\, K_{0}(\omega_{f}Y)$$
$$\mathcal{A}^{(01)} = \frac{K_{0}(\epsilon_{f}R)}{\zeta\,(1-\zeta)}\,\frac{\omega_{f}\,K_{1}(\omega_{f}Y)}{Y}$$
$$\mathcal{A}^{(00)} = \frac{K_{0}(\epsilon_{f}R)}{\zeta\,(1-\zeta)}\,K_{0}(\omega_{f}Y)$$
Finally:
$$\mathcal{G}_{E} = \frac{E_{f}\,K_{1}(E_{f}\mathcal{D})}{\mathcal{D}}$$
where: 
$$\mathcal{D}^{2}(x,y,z,\zeta,\xi) = \zeta\,(1-\zeta)\,R^{2} + \frac{\xi\,\eta\,Y^{2}}{(1-\zeta)} = \zeta\,\eta\,(x-y)^{2} + \zeta\,\xi\,(x-z)^{2} + \eta\,\xi\,(y-z)^{2}$$
$$\mu^{i}_{\lambda\lambda'} = \delta_{\lambda',-\lambda}\left[2\lambda\delta^{i1} - i\delta^{i2}\right] = \chi^{\dagger}_{\lambda}\sigma^{3}\sigma^{i}\chi_{\lambda'}$$
$$\epsilon^{2}_{f} = \Tilde{Q}^{2}+ m^{2}_{f} = \zeta\, (1-\zeta)\,Q^{2} +m^{2}_{f}, \hspace{2em} \omega_{f} = \frac{\xi\,m_{f}}{(1-\zeta)}, \hspace{2em} E^{2}_{f} = Q^{2} +m^{2}_{f}\,\left(\frac{1}{\zeta} + \frac{1}{\eta}\right),\hspace{2em} a = \frac{\xi\,\eta}{\zeta\,(1-\zeta)^{2}}$$
$$\varphi^{ij}_{\lambda_{1}\lambda}(\zeta) = \delta_{\lambda_{1}\lambda}\left[(2\zeta -1)\delta^{ij} +2i\epsilon^{ij}\lambda_{1}\right],\hspace{3em} \tau^{mn}_{\lambda \lambda_{2}}(\xi,\eta) = \delta_{\lambda \lambda_{2}}\,\left[(2(1-\zeta)-\xi)\delta^{mn}+2i\,\xi\,\epsilon^{mn}\,\lambda\right] $$
and 
$$c = \zeta\,x +\eta\, y + \xi\,z$$ is the centre of energy of trijet system.
$\mathcal{W}(x,y,z,\Bar{x},\Bar{y},\Bar{z})$ is the Wilson line structure. This has 4 terms:
\begin{itemize}
    \item $\mathcal{W}(x,y,z,\Bar{x},\Bar{y},\Bar{z})$
    \begin{align}
        \mathcal{W}(x,y,z,\Bar{x},\Bar{y},\Bar{z}) & = S_{q\Bar{q}g\,q\Bar{q}g}(x,y,z,\Bar{x},\Bar{y},\Bar{z}) -S_{q\Bar{q}g}(x,y,z) - S^{\dagger}_{q\Bar{q}g}(\Bar{x},\Bar{y},\Bar{z}) + 1
    \end{align}
At large $N_{c}$ limit, we get
$$S_{q\Bar{q}g\,q\Bar{q}g}(x,y,z,\Bar{x},\Bar{y},\Bar{z}) = Q(x,z,\Bar{z},\Bar{x})\, Q(z,y,\Bar{y},\Bar{z})$$
$$S_{q\Bar{q}g}(x,y,z) = S(x,z)\, S(z,y)$$
$$S^{\dagger}_{q\Bar{q}g}(\Bar{x},\Bar{y},\Bar{z}) = S(\Bar{z},\Bar{x})\, S(\Bar{y},\Bar{z})$$
\boxed{$$\mathcal{W}(x,y,z,\Bar{x},\Bar{y},\Bar{z}) = Q(x,z,\Bar{z},\Bar{x})\, Q(z,y,\Bar{y},\Bar{z}) - S(x,z)\, S(z,y) - S(\Bar{z},\Bar{x})\, S(\Bar{y},\Bar{z}) +1 $$}\\ \\ 
Using Gaussian approximation, we get \\  \\ 
\begin{align}
\mathcal{W}(x,y,z,\Bar{x},\Bar{y},\Bar{z})  & = \left[S(x,z)\,S(\Bar{z},\Bar{x}) - \frac{F(x,z;\Bar{z},\Bar{x})}{F(x,\Bar{z};z,\Bar{x})}\, \left[S(x,z)\,S(\Bar{z},\Bar{x})-S(x,\Bar{x})\, S(\Bar{z},z)\right]\right]\times \notag \\ & \left[S(z,y)\,S(\Bar{y},\Bar{z}) - \frac{F(z,y;\Bar{y},\Bar{z})}{F(z,\Bar{y};y,\Bar{z})}\,\left[S(z,y)\,S(\Bar{y},\Bar{z})-S(z,\Bar{z})\,S(\Bar{y},y)\right]\right] - S(x,z)\, S(z,y) - S(\Bar{z},\Bar{x})\, S(\Bar{y},\Bar{z}) +1
\end{align}
\item $\mathcal{W}(x,y',y',\Bar{x},\Bar{y},\Bar{z})$
\begin{align}
        \mathcal{W}(x,y',y',\Bar{x},\Bar{y},\Bar{z}) & = S_{q\Bar{q}g\,q\Bar{q}g}(x,y',y',\Bar{x},\Bar{y},\Bar{z}) -S_{q\Bar{q}g}(x,y',y') - S^{\dagger}_{q\Bar{q}g}(\Bar{x},\Bar{y},\Bar{z}) + 1
    \end{align}
At large $N_{c}$ limit, we get
$$S_{q\Bar{q}g\,q\Bar{q}g}(x,y',y',\Bar{x},\Bar{y},\Bar{z}) = Q(x,y',\Bar{z},\Bar{x})\, Q(y',y',\Bar{y},\Bar{z}) = Q(x,y',\Bar{z},\Bar{x})\, S(\Bar{y},\Bar{z})$$
$$S_{q\Bar{q}g}(x,y',y') = S(x,y')\, S(y',y') = S(x,y') $$
$$S^{\dagger}_{q\Bar{q}g}(\Bar{x},\Bar{y},\Bar{z}) = S(\Bar{z},\Bar{x})\, S(\Bar{y},\Bar{z})$$
\boxed{$$\mathcal{W}(x,y',y',\Bar{x},\Bar{y},\Bar{z}) = Q(x,y',\Bar{z},\Bar{x})\, S(\Bar{y},\Bar{z}) -  S(x,y') - S(\Bar{z},\Bar{x})\, S(\Bar{y},\Bar{z}) +1  $$} \\ \\ 
Using Gaussian approximation, we get 
\begin{align}
    \mathcal{W}(x,y',y',\Bar{x},\Bar{y},\Bar{z}) &  =  \left[S(x,y')\,S(\Bar{z},\Bar{x}) - \frac{F(x,y';\Bar{z},\Bar{x})}{F(x,\Bar{z};y',\Bar{x})}\, \left[S(x,y')\,S(\Bar{z},\Bar{x})-S(x,\Bar{x})\, S(\Bar{z},y')\right]\right]\,S(\Bar{y},\Bar{z}) -  S(x,y') - S(\Bar{z},\Bar{x})\, S(\Bar{y},\Bar{z}) +1 
\end{align}
\item $\mathcal{W}(x,y,z,\Bar{x},\Bar{y}',\Bar{y}')$
\begin{align}
        \mathcal{W}(x,y,z,\Bar{x},\Bar{y}',\Bar{y}') & = S_{q\Bar{q}g\,q\Bar{q}g}(x,y,z,\Bar{x},\Bar{y}',\Bar{y}') -S_{q\Bar{q}g}(x,y,z) - S^{\dagger}_{q\Bar{q}g}(\Bar{x},\Bar{y}',\Bar{y}') + 1
    \end{align}
At large $N_{c}$ limit, we get
$$S_{q\Bar{q}g\,q\Bar{q}g}(x,y,z,\Bar{x},\Bar{y}',\Bar{y}') = Q(x,z,\Bar{y}',\Bar{x})\, Q(z,y,\Bar{y}',\Bar{y}') = Q(x,z,\Bar{y}',\Bar{x})\, S(z,y) $$
$$S_{q\Bar{q}g}(x,y,z) = S(x,z)\, S(z,y)$$
$$S^{\dagger}_{q\Bar{q}g}(\Bar{x},\Bar{y}',\Bar{y}') = S(\Bar{y}',\Bar{x})\, S(\Bar{y}',\Bar{y}') = S(\Bar{y}',\Bar{x})$$
\boxed{$$\mathcal{W}(x,y,z,\Bar{x},\Bar{y}',\Bar{y}') = Q(x,z,\Bar{y}',\Bar{x})\, S(z,y) - S(x,z)\, S(z,y) -  S(\Bar{y}',\Bar{x}) +1 $$}\\ \\ 
Using Gaussian approximation, we get
\begin{align}
    \mathcal{W}(x,y,z,\Bar{x},\Bar{y}',\Bar{y}') & =  \left[S(x,z)\,S(\Bar{y}',\Bar{x}) - \frac{F(x,z;\Bar{y}',\Bar{x})}{F(x,\Bar{y}';z,\Bar{x})}\, \left[S(x,z)\,S(\Bar{y}',\Bar{x})-S(x,\Bar{x})\, S(\Bar{y}',z)\right]\right] S(z,y) - S(x,z)\, S(z,y) -  S(\Bar{y}',\Bar{x}) +1
\end{align}
\item $\mathcal{W}(x,y',y',\Bar{x},\Bar{y}',\Bar{y}')$
    \begin{align}
        \mathcal{W}(x,y',y',\Bar{x},\Bar{y}',\Bar{y}') & = S_{q\Bar{q}g\,q\Bar{q}g}(x,y',y',\Bar{x},\Bar{y}',\Bar{y}') -S_{q\Bar{q}g}(x,y',y') - S^{\dagger}_{q\Bar{q}g}(\Bar{x},\Bar{y}',\Bar{y}') + 1
    \end{align}
At large $N_{c}$ limit, we get
$$S_{q\Bar{q}g\,q\Bar{q}g}(x,y',y',\Bar{x},\Bar{y}',\Bar{y}') = Q(x,y',\Bar{y}',\Bar{x})\, Q(y',y',\Bar{y}',\Bar{y'}) = Q(x,y',\Bar{y}',\Bar{x}) $$
$$S_{q\Bar{q}g}(x,y',y') = S(x,y')\, S(y',y') = S(x,y')$$
$$S^{\dagger}_{q\Bar{q}g}(\Bar{x},\Bar{y}',\Bar{y}') = S(\Bar{y}',\Bar{x})\, S(\Bar{y}',\Bar{y}') = S(\Bar{y}',\Bar{x})$$
\boxed{$$ \mathcal{W}(x,y',y',\Bar{x},\Bar{y}',\Bar{y}') = Q(x,y',\Bar{y}',\Bar{x}) -S(x,y') -  S(\Bar{y}',\Bar{x}) +1 $$} \\ \\ 
Using Gaussian approximation, we get
\begin{align}
    \mathcal{W}(x,y',y',\Bar{x},\Bar{y}',\Bar{y}') & = \left[S(x,y')\,S(\Bar{y}',\Bar{x}) - \frac{F(x,y';\Bar{y}',\Bar{x})}{F(x,\Bar{y}';y',\Bar{x})}\, \left[S(x,y')\,S(\Bar{y}',\Bar{x})-S(x,\Bar{x})\, S(\Bar{y}',y')\right]\right] -S(x,y') -  S(\Bar{y}',\Bar{x}) +1
\end{align}
\end{itemize}
\subsection{Direct contribution - Quark contributions}
The following transformations are made from the anti quark sector to quark sector: \
$$\zeta\rightarrow \eta, R\rightarrow \Tilde{R}, Y \rightarrow X, x\leftrightarrow y$$
and 
$$\epsilon_{f}\rightarrow \Tilde{\epsilon}_{f}, \omega_{f}\rightarrow \Tilde{\omega}_{f}, \mathcal{G} \rightarrow \mathcal{\Tilde{G}}, \mathcal{A} \rightarrow \mathcal{\Tilde{A}}$$
\begin{align}
     \frac{d^{3}N^{\text{direct}}}{d^{3}k_{1}\,d^{3}k_{2}\, d^{3}k_{3}} & = \frac{e^{2}\,g^{2}\,\eta^{2}\,C_{F}\,N_{c}}{8\, (2\pi)^{12}\,\xi\,q^{+\,2}}\, \delta\left(q^{+}-k^{+}_{1}-k^{+}_{2}-k^{+}_{3}\right)\, \int_{x,y,z,\Bar{x},\Bar{y},\Bar{z}}\, e^{-i\,k_{1}\cdot (x-\Bar{x})}\, e^{-i\,k_{2}\cdot (y-\Bar{y})}\, e^{-i\,k_{3}\cdot (z-\Bar{z})}\,\left(\sum_{f}\, e^{2}_{f}\right)\notag \\ & \Bigg[\mathcal{\Tilde{H}}[\mathcal{\Tilde{G}};\mathcal{\Bar{\Tilde{G}}}]\,\mathcal{W}(y,x,z,\Bar{y},\Bar{x},\Bar{z}) - \mathcal{\Tilde{H}}[\mathcal{\Tilde{A}};\mathcal{\Bar{\Tilde{G}}}]\,\mathcal{W}(y,x',x',\Bar{y},\Bar{x},\Bar{z})-\mathcal{\Tilde{H}}[\mathcal{\Tilde{G}};\mathcal{\Bar{\Tilde{A}}}]\,\mathcal{W}(y,x,z,\Bar{y},\Bar{x}',\Bar{x}')+ \mathcal{\Tilde{H}}[\mathcal{\Tilde{A}};\mathcal{\Bar{\Tilde{A}}}]\,\mathcal{W}(y,x',x',\Bar{y},\Bar{x}',\Bar{x}')\Bigg]
\end{align}
where:
\begin{align}
    \mathcal{\Tilde{H}}[F;\Bar{F}] & = \sum_{\lambda_{1},\lambda_{2}}\, \sum_{l,r}\, \mathcal{\Tilde{M}}^{lr}_{\lambda_{1}\lambda_{2}}[F](x,y,z)\, \mathcal{\Tilde{M}}^{lr\,*}_{\lambda_{1}\lambda_{2}}[\Bar{F}](\Bar{x},\Bar{y},\Bar{z})
\end{align}
\begin{itemize}
    \item Regular $\times$ Regular
    \begin{align}
\mathcal{\Tilde{H}}_{\mathrm{reg.}\times \mathrm{reg.}} & = \Tilde{R}^{i}\, X^{m}\, \Bar{\Tilde{R}}^{j}\,\Bar{X}^{n}\,\Bigg\{8\,\left[\zeta^{2} + (1-\eta)^{2}\right]\,\left[\eta^{2} + (1-\eta)^{2}\right]\,\delta^{ij}\delta^{mn}- 8\xi\, (2\eta -1)\,(2(1-\eta)-\xi)\,\epsilon^{ij}\,\epsilon^{mn}\Bigg\}F^{(11)}\Bar{F}^{(11)}\notag \\ & +8m^{2}_{f}\left[\zeta^{2} + (1-\eta)^{2}\right]\, (X\cdot \Bar{X})\,F^{(01)}\,\Bar{F}^{(01)} + 8\xi^{2}\Tilde{\omega}^{2}_{f}\left[\eta^{2}+(1-\eta)^{2}\right]\,(\Tilde{R}\cdot \Bar{\Tilde{R}})\,F^{(10)}\Bar{F}^{(10)} +8m^{2}_{f}\xi^{2}\Tilde{\omega}^{2}_{f}\,F^{(00)}\Bar{F}^{(00)}\notag \\ & +8\zeta\,\eta\, m_{f}\, \xi\, \Tilde{\omega}_{f}\, \left[(\Tilde{R}\cdot X)\, F^{(11)}\Bar{F}^{(00)} + (\Bar{\Tilde{R}}\cdot \Bar{X})\,F^{(00)}\Bar{F}^{(11)} + (\Tilde{R}\cdot \Bar{X})\,F^{(10)}\Bar{F}^{(01)} + (X\cdot \Bar{\Tilde{R}})\,F^{(01)}\Bar{F}^{(10)}\right]
    \end{align}
    \item Regular $\times$ Instantaneous $+$ h.c.
    \begin{align}
               \mathcal{\Tilde{H}}_{\mathrm{reg.}\times \mathrm{inst.}} & = \frac{8\,\zeta^{2}\,\xi\,\eta}{(1-\eta)}\,\left[(\Tilde{R}\cdot X)\,F^{(11)}\,\Bar{\mathcal{G}}_{E} + (\Bar{\Tilde{R}}\cdot \Bar{X})\,\mathcal{G}_{E}\,\Bar{F}^{(11)}\right]\notag \\ & + \frac{8\,\xi^{2}\,\zeta\,m_{f}\,\Tilde{\omega}_{f}}{(1-\eta)}\, \left[F^{(00)}\,\Bar{\mathcal{G}}_{E} + \mathcal{G}_{E}\,\Bar{F}^{(00)}\right]
    \end{align}
\item Instantaneous $\times$ Instantaneous
    \begin{align}
\mathcal{\Tilde{H}}_{\mathrm{inst.}\times \mathrm{inst.}} & = 8\, \left[\frac{\xi\,\zeta}{(1-\eta)}\right]^{2}\, \mathcal{G}_{E}\,\Bar{\mathcal{G}}_{E}
    \end{align}
\end{itemize}
where:
\begin{align}
\mathcal{\Tilde{M}}^{im}_{\lambda_{1}\lambda_{2}}\chg{[\Tilde{\mathcal{G}}] = \mathcal{\Tilde{M}}^{im}_{1\lambda_{1}\lambda_{2}}[\Tilde{\mathcal{G}}]} &  = \Tilde{\Psi}^{im}_{\lambda_{1}\lambda_{2}}[\mathcal{\Tilde{G}}](\Tilde{R},X) - \delta_{\lambda_{1}\lambda_{2}}\,\frac{\zeta\,\xi}{(\xi + \zeta)}\,(\delta^{im}-2i\,\epsilon^{im}\lambda_{1})\,\mathcal{G}_{E} 
\label{chg:dupM}\end{align}
\begin{align}
\mathcal{\Tilde{M}}^{im}_{2\lambda_{1}\lambda_{2}}\chg{[\mathcal{\Tilde{A}}]} &  = \Tilde{\Psi}^{im}_{\lambda_{1}\lambda_{2}}[\mathcal{\Tilde{A}}] (\Tilde{R},X)
\label{chg:M2q}\end{align}
with 
\begin{align}
    \Tilde{\Psi}^{im}_{\lambda_{1}\lambda_{2}}[\mathcal{\Tilde{G}}] & = \sum_{\lambda}\, \Bigg[\tau^{nm}_{\lambda_{1}\lambda}(\xi,\zeta)\varphi^{ij}_{\lambda\lambda_{2}}(\zeta + \xi)\,\Tilde{R}^{j}X^{n}\,\mathcal{\Tilde{G}}^{(11)}+i\xi\,\Tilde{\omega}_{f}\mu^{m}_{\lambda_{1}\lambda}\varphi^{ij}_{\lambda\lambda_{2}}(\zeta + \xi)\Tilde{R}^{j}\mathcal{\Tilde{G}}^{(10)}\notag \\ & -im_{f}\,\tau^{nm}_{\lambda_{1}\lambda}(\xi,\zeta)\,\mu^{i}_{\lambda\lambda_{2}}\,X^{n}\mathcal{\Tilde{G}}^{(01)} + m_{f}\xi\,\Tilde{\omega}_{f}\mu^{m}_{\lambda_{1}\lambda}\mu^{i}_{\lambda\lambda_{2}}\mathcal{\Tilde{G}}^{(00)}\Bigg]
\end{align}
where
\label{chg:Gqarg}$$\mathcal{\Tilde{G}}^{(\alpha\beta)}(\chg{\Tilde{R}},X) = \frac{\Tilde{a}^{\beta}}{2\,\eta(1-\eta)}\,\int_{0}^{1}\, \frac{dx}{x^{1+\beta}}\, \left(\frac{\Tilde{A}_{x}}{\Tilde{\Delta}_{x}}\right)^{\frac{(\alpha + \beta)}{2}}\,K_{\alpha + \beta} (\sqrt{\Tilde{A_{x}}\Tilde{\Delta_{x}}}), \alpha, \beta \in \{0,1\}$$
$$\Tilde{A}_{x} = \Tilde{\epsilon}^{2}_{f} + \frac{x\,\Tilde{\omega}^{2}_{f}}{\Tilde{a}} = \eta\,(1-\eta)\,Q^{2} +m^{2}_{f}\,\left(1+\frac{x\eta\xi}{\zeta}\right)$$
$$\Tilde{\Delta}_{x} = \Tilde{R}^{2} + \frac{\Tilde{a}\,X^{2}}{x}$$
\label{chg:ty}\chg{Explicitly}
\begin{align}
    \mathcal{\Tilde{G}}^{(11)} & = \frac{\Tilde{a}}{2\,\eta\, (1-\eta)}\, \int_{0}^{1}\, \frac{dx}{x^{2}}\, \frac{\Tilde{A}_{x}}{\Tilde{\Delta}_{x}}\,K_{2}\left(\sqrt{\Tilde{A}_{x}\Tilde{\Delta}_{x}}\right)
\end{align}
\begin{align}
    \mathcal{\Tilde{G}}^{(10)} & = \frac{1}{2\, \eta\, (1-\eta)}\, \int_{0}^{1}\, \frac{dx}{x}\, \sqrt{\frac{\Tilde{A}_{x}}{\Tilde{\Delta}_{x}}}\, K_{1}\left(\sqrt{\Tilde{A}_{x}\Tilde{\Delta}_{x}}\right)
\end{align}
\begin{align}
     \mathcal{\Tilde{G}}^{(01)} & = \frac{\Tilde{a}}{2\, \eta\, (1-\eta)}\, \int_{0}^{1}\, \frac{dx}{x^{2}}\, \sqrt{\frac{\Tilde{A}_{x}}{\Tilde{\Delta}_{x}}}\, K_{1}\left(\sqrt{\Tilde{A}_{x}\,\Tilde{\Delta}_{x}}\right)
\end{align}
\begin{align}
     \mathcal{\Tilde{G}}^{(00)} & = \frac{1}{2\, \eta\, (1-\eta)}\, \int_{0}^{1}\, \frac{dx}{x}\, K_{0}\left(\sqrt{\Tilde{A}_{x}\Tilde{\Delta}_{x}}\right)
\end{align}
and
\begin{align}
  \Tilde{\Psi}^{im}_{\lambda_{1}\lambda_{2}}[\Tilde{\mathcal{A}}] & = \sum_{\lambda}\, \Bigg[\tau^{nm}_{\lambda_{1}\lambda}(\xi,\zeta)\,\varphi^{ij}_{\lambda\lambda_{2}}(\zeta + \xi)\,\Tilde{R}^{j}X^{n}\Tilde{\mathcal{A}}^{(11)}+i\xi\,\Tilde{\omega}_{f}\,\mu^{m}_{\lambda_{1}\lambda}\,\varphi^{ij}_{\lambda\lambda_{2}}(\zeta + \xi)\,\Tilde{R}^{j}\Tilde{\mathcal{A}}^{(10)}\notag \\ & -i\,m_{f}\,\tau^{nm}_{\lambda_{1}\lambda}(\xi,\zeta)\mu^{i}_{\lambda\lambda_{2}}\,X^{n}\,\Tilde{\mathcal{A}}^{(01)} + m_{f}\,\xi\,\chg{\Tilde{\omega}_{f}}\,\mu^{m}_{\lambda_{1}\lambda}\,\mu^{i}_{\lambda\lambda_{2}}\,\Tilde{\mathcal{A}}^{(00)}\Bigg]  
\label{chg:omq}\end{align}
Explicitly: 
\begin{align}
    \mathcal{\Tilde{A}}^{(11)} & = \frac{\Tilde{\epsilon}_{f}\,K_{1}(\Tilde{\epsilon}_{f}\,\Tilde{R})}{\eta\,(1-\eta)\,\Tilde{R}}\, \frac{\Tilde{\omega}_{f}\,K_{1}(\Tilde{\omega}_{f}\,X)}{X}
\end{align}
\begin{align}
    \mathcal{\Tilde{A}}^{(10)} & = \frac{\Tilde{\epsilon}_{f}\,K_{1}(\Tilde{\epsilon}_{f}\Tilde{R})}{\eta\,(1-\eta)\,\Tilde{R}}\, K_{0}(\Tilde{\omega}_{f}\,X)
\end{align}
\begin{align}
    \mathcal{\Tilde{A}}^{(01)} & =   \frac{K_{0}(\Tilde{\epsilon}_{f}\Tilde{R})}{\eta\,(1-\eta)}\,\frac{\Tilde{\omega}_{f}\,K_{1}(\Tilde{\omega}_{f}\,X)}{X}
\end{align}
\begin{align}
    \mathcal{\Tilde{A}}^{(00)} & = \frac{K_{0}(\Tilde{\epsilon}_{f}\,\Tilde{R})}{\eta\,(1-\eta)}\, K_{0}(\Tilde{\omega}_{f}\,X)
\end{align}
and 
\begin{itemize}
 \item $\xi\, \omega_{f}\,\mu^{m}_{\lambda\lambda_{2}} \rightarrow \xi\, \Tilde{\omega}_{f}\,\mu^{m}_{\lambda_{1}\lambda}$ with $\Tilde{\omega}_{f} = \frac{\xi\,m_{f}}{(1-\eta)}$
    \item $\epsilon^{2}_{f} \rightarrow \Tilde{\epsilon}^{2}_{f} = \eta(1-\eta)\, Q^{2} + m^{2}_{f}$, $a \rightarrow \Tilde{a} = \frac{\xi\,\zeta}{\eta\, (1-\eta)^{2}}$
    \item $A_{x} \rightarrow \Tilde{A}_{x} = \eta\,(1-\eta)\,Q^{2} + m^{2}_{f}\left(1+\frac{x\eta\xi}{\zeta}\right)$ and $\Delta_{x} \rightarrow \Tilde{\Delta}_{x} = \Tilde{R}^{2} + \frac{\Tilde{a}\,X^{2}}{x}$
    \item $\mathcal{G}^{(\alpha\beta)}(R,Y) \rightarrow \mathcal{\Tilde{G}}^{(\alpha\beta)}(\Tilde{R},X) $ and $\mathcal{A}^{(\alpha,\beta)}(R,Y) \rightarrow \mathcal{\Tilde{A}}^{(\alpha\beta)}(\Tilde{R},X)$ also, $\zeta\,(1-\zeta) \rightarrow \eta\,(1-\eta)$
\end{itemize}
$\mathcal{W}(y,x,z,\Bar{y},\Bar{x},\Bar{z})$ is the Wilson line structure. This has 4 terms:
\begin{itemize}
    \item $\mathcal{W}(y,x,z,\Bar{y},\Bar{x},\Bar{z})$
    \begin{align}
        \mathcal{W}(y,x,z,\Bar{y},\Bar{x},\Bar{z}) & = S_{q\Bar{q}g\,q\Bar{q}g}(y,x,z,\Bar{y},\Bar{x},\Bar{z}) -S_{q\Bar{q}g}(y,x,z) - S^{\dagger}_{q\Bar{q}g}(\Bar{y},\Bar{x},\Bar{z}) + 1
    \end{align}
At large $N_{c}$ limit, we get
$$S_{q\Bar{q}g\,q\Bar{q}g}(y,x,z,\Bar{y},\Bar{x},\Bar{z}) = Q(y,z,\Bar{z},\Bar{y})\, Q(z,x,\Bar{x},\Bar{z})$$
$$S_{q\Bar{q}g}(y,x,z) = S(y,z)\, S(z,x)$$
$$S^{\dagger}_{q\Bar{q}g}(\Bar{y},\Bar{x},\Bar{z}) = S(\Bar{z},\Bar{y})\, S(\Bar{x},\Bar{z})$$
\boxed{$$\mathcal{W}(y,x,z,\Bar{y},\Bar{x},\Bar{z}) = Q(y,z,\Bar{z},\Bar{y})\, Q(z,x,\Bar{x},\Bar{z}) - S(y,z)\, S(z,x) - S(\Bar{z},\Bar{y})\, S(\Bar{x},\Bar{z}) +1 $$}\\ \\ 
Using Gaussian approximation
\begin{align}
    \mathcal{W}(y,x,z,\Bar{y},\Bar{x},\Bar{z}) & = \left[S(y,z)\,S(\Bar{z},\Bar{y}) - \frac{F(y,z;\Bar{z},\Bar{y})}{F(y,\Bar{z};z,\Bar{y})}\, \left[S(y,z)\,S(\Bar{z},\Bar{y})-S(y,\Bar{y})\, S(\Bar{z},z)\right]\right]\times\notag \\ & \left[S(z,x)\,S(\Bar{x},\Bar{z}) - \frac{F(z,x;\Bar{x},\Bar{z})}{F(z,\Bar{x};x,\Bar{z})}\, \left[S(z,x)\,S(\Bar{x},\Bar{z})-S(z,\Bar{z})\, S(\Bar{x},x)\right]\right] - S(y,z)\, S(z,x) - S(\Bar{z},\Bar{y})\, S(\Bar{x},\Bar{z}) +1
\end{align}
\item $\mathcal{W}(y,x',x',\Bar{y},\Bar{x},\Bar{z})$
\begin{align}
        \mathcal{W}(y,x',x',\Bar{y},\Bar{x},\Bar{z}) & = S_{q\Bar{q}g\,q\Bar{q}g}(y,x',x',\Bar{y},\Bar{x},\Bar{z}) -S_{q\Bar{q}g}(y,x',x') - S^{\dagger}_{q\Bar{q}g}(\Bar{y},\Bar{x},\Bar{z}) + 1
    \end{align}
At large $N_{c}$ limit, we get
$$S_{q\Bar{q}g\,q\Bar{q}g}(y,x',x',\Bar{y},\Bar{x},\Bar{z}) = Q(y,x',\Bar{z},\Bar{y})\, Q(x',x',\Bar{x},\Bar{z}) = Q(y,x',\Bar{z},\Bar{y})\, S(\Bar{x},\Bar{z})$$
$$S_{q\Bar{q}g}(y,x',x') = S(y,x')\, S(x',x') = S(y,x') $$
$$S^{\dagger}_{q\Bar{q}g}(\Bar{y},\Bar{x},\Bar{z}) = S(\Bar{z},\Bar{y})\, S(\Bar{x},\Bar{z})$$
\boxed{$$\mathcal{W}(y,x',x',\Bar{y},\Bar{x},\Bar{z}) = Q(y,x',\Bar{z},\Bar{y})\, S(\Bar{x},\Bar{z}) -  S(y,x') - S(\Bar{z},\Bar{y})\, S(\Bar{x},\Bar{z}) +1 $$}\\ \\ 
Using Gaussian approximation, we get
\begin{align}
    \mathcal{W}(y,x',x',\Bar{y},\Bar{x},\Bar{z}) &  = \left[S(y,x')\,S(\Bar{z},\Bar{y}) - \frac{F(y,x';\Bar{z},\Bar{y})}{F(y,\Bar{z};x',\Bar{y})}\, \left[S(y,x')\,S(\Bar{z},\Bar{y})-S(y,\Bar{y})\,S(\Bar{z},x')\right]\right] S(\Bar{x},\Bar{z}) -  S(y,x') - S(\Bar{z},\Bar{y})\, S(\Bar{x},\Bar{z}) +1
\end{align}
\item $\mathcal{W}(y,x,z,\Bar{y},\Bar{x}',\Bar{x}')$
\begin{align}
        \mathcal{W}(y,x,z,\Bar{y},\Bar{x}',\Bar{x}') & = S_{q\Bar{q}g\,q\Bar{q}g}(y,x,z,\Bar{y},\Bar{x}',\Bar{x}') -S_{q\Bar{q}g}(y,x,z) - S^{\dagger}_{q\Bar{q}g}(\Bar{y},\Bar{x}',\Bar{x}') + 1
    \end{align}
At large $N_{c}$ limit, we get
$$S_{q\Bar{q}g\,q\Bar{q}g}(y,x,z,\Bar{y},\Bar{x}',\Bar{x}') = Q(y,z,\Bar{x}',\Bar{y})\, Q(z,x,\Bar{x}',\Bar{x}') = Q(y,z,\Bar{x}',\Bar{y})\, S(z,x) $$
$$S_{q\Bar{q}g}(y,x,z) = S(y,z)\, S(z,x)$$
$$S^{\dagger}_{q\Bar{q}g}(\Bar{y},\Bar{x}',\Bar{x}') = S(\Bar{x}',\Bar{y})\, S(\Bar{x}',\Bar{x}') = S(\Bar{x}',\Bar{y})$$
\boxed{$$\mathcal{W}(y,x,z,\Bar{y},\Bar{x}',\Bar{x}') = Q(y,z,\Bar{x}',\Bar{y})\, S(z,x) - S(y,z)\, S(z,x) - S(\Bar{x}',\Bar{y}) + 1  $$}\\ \\ 
Using Gaussian approximation, we get
\begin{align}
    \mathcal{W}(y,x,z,\Bar{y},\Bar{x}',\Bar{x}') & = \left[S(y,z)\,S(\Bar{x}',\Bar{y}) - \frac{F(y,z;\Bar{x}',\Bar{y})}{F(y,\Bar{x}';z,\Bar{y})}\, \left[S(y,z)\,S(\Bar{x}',\Bar{y})-S(y,\Bar{y})\, S(\Bar{x}',z)\right]\right]S(z,x) - S(y,z)\, S(z,x) - S(\Bar{x}',\Bar{y}) + 1
\end{align}
\item $\mathcal{W}(y,x',x',\Bar{y},\Bar{x}',\Bar{x}')$
    \begin{align}
        \mathcal{W}(y,x',x',\Bar{y},\Bar{x}',\Bar{x}') & = S_{q\Bar{q}g\,q\Bar{q}g}(y,x',x',\Bar{y},\Bar{x}',\Bar{x}') -S_{q\Bar{q}g}(y,x',x') - S^{\dagger}_{q\Bar{q}g}(\Bar{y},\Bar{x}',\Bar{x}') + 1
    \end{align}
At large $N_{c}$ limit, we get
$$S_{q\Bar{q}g\,q\Bar{q}g}(y,x',x',\Bar{y},\Bar{x}',\Bar{x}') = Q(y,x',\Bar{x}',\Bar{y})\, Q(x',x',\Bar{x}',\Bar{x'}) = Q(y,x',\Bar{x}',\Bar{y}) $$
$$S_{q\Bar{q}g}(y,x',x') = S(y,x')\, S(x',x') = S(y,x')$$
$$S^{\dagger}_{q\Bar{q}g}(\Bar{y},\Bar{x}',\Bar{x}') = S(\Bar{x}',\Bar{y})\, S(\Bar{x}',\Bar{x}') = S(\Bar{x}',\Bar{y})$$ 
\boxed{$$\mathcal{W}(y,x',x',\Bar{y},\Bar{x}',\Bar{x}') =Q(y,x',\Bar{x}',\Bar{y}) - S(y,x') -  S(\Bar{x}',\Bar{y}) +1  $$}\\ \\ 
Using Gaussian approximation, we get
\begin{align}
    \mathcal{W}(y,x',x',\Bar{y},\Bar{x}',\Bar{x}') & = S(y,x')\,S(\Bar{x}',\Bar{y}) - \frac{F(y,x';\Bar{x}',\Bar{y})}{F(y,\Bar{x}';x',\Bar{y})}\, \left[S(y,x')\,S(\Bar{x}',\Bar{y})-S(y,\Bar{y})\, S(\Bar{x}',x')\right] - S(y,x') -  S(\Bar{x}',\Bar{y}) +1
\end{align}
\end{itemize}
\subsection{Interference contributions} 
\begin{align}
     \frac{d^{3}N^{\text{cross}}}{d^{3}k_{1}\,d^{3}k_{2}\, d^{3}k_{3}} & = -\frac{e^{2}\,g^{2}\,\eta\,\zeta\,C_{F}\,N_{c}}{8\, (2\pi)^{12}\,\xi\,q^{+\,2}}\, \delta\left(q^{+}-k^{+}_{1}-k^{+}_{2}-k^{+}_{3}\right)\times\, 2\,\mathcal{\Re}\,\int_{x,y,z,\Bar{x},\Bar{y},\Bar{z}}\, e^{-i\,k_{1}\cdot (x-\Bar{x})}\, e^{-i\,k_{2}\cdot (y-\Bar{y})}\, e^{-i\,k_{3}\cdot (z-\Bar{z})}\,\left(\sum_{f}\, e^{2}_{f}\right)\notag \\ &  \Bigg[\mathcal{H}^{x}[\mathcal{\Tilde{G}}\chg{;}\Bar{\mathcal{G}}]\,\mathcal{W}(x,y,z,\Bar{x},\Bar{y},\Bar{z})-\mathcal{H}^{x}[\Tilde{\mathcal{A}};\Bar{\mathcal{G}}]\,\mathcal{W}(x',y,x',\Bar{x},\Bar{y},\Bar{z})-\mathcal{H}^{x}[\Tilde{\mathcal{G}};\Bar{\mathcal{A}}]\,\mathcal{W}(x,y,z,\Bar{x},\Bar{y}',\Bar{y}')\notag \\ & +\mathcal{H}^{x}[\Tilde{\mathcal{A}};\Bar{\mathcal{A}}]\,\mathcal{W}(x',y,x',\Bar{x},\Bar{y}',\Bar{y}')\Bigg]
\label{chg:ty2}\end{align}  
where
\begin{align}
    \mathcal{H}^{x} [\Tilde{F}\chg{;}\Bar{F}] & = \sum_{\lambda_{1},\lambda_{2}}\, \sum_{l,r}\, \Tilde{\mathcal{M}}^{lr}_{\lambda_{1}\lambda_{2}}[\Tilde{F}](x,y,z)\,\mathcal{M}^{lr\,*}_{\lambda_{1}\lambda_{2}}[\Bar{F}] (\Bar{x},\Bar{y},\Bar{z})
\label{chg:ty3}\end{align}
\begin{itemize}
    \item Regular $\times$ Regular
    \begin{align}
    \mathcal{H}^{x}_{\mathrm{reg.}\times \mathrm{reg.}} & = \Tilde{R}^{i}\,X^{m}\,\Bar{R}^{j}\,\Bar{Y}^{n}\,\bigg\{8[1-\xi-2\zeta\eta]\,[\xi(1-\xi) + 2\zeta\eta]\,\delta^{ij}\delta^{mn}-8\xi[1-2\zeta-\xi]^{2}\,\epsilon^{ij}\epsilon^{mn}\bigg\}\Tilde{F}^{(11)}\,\Bar{F}^{(11)}\notag \\ & +8\,m^{2}_{f}\,(\xi +2\zeta\eta)\, (X\cdot \Bar{Y})\,\Tilde{F}^{(01)}\Bar{F}^{(01)} - 8(\xi + 2\eta\zeta)\xi^{2}\,\Tilde{\omega}_{f}\omega_{f}\,(\Tilde{R}\cdot\Bar{R})\,\Tilde{F}^{(10)}\Bar{F}^{(10)} - 8(\zeta + \xi)^{2}\,m_{f}\,\xi\,\omega_{f}\,(\Tilde{R}\cdot X)\,\Tilde{F}^{(11)}\Bar{F}^{(00)}\notag \\ & + 8(1-\zeta)^{2}\,m_{f}\,\xi\,\Tilde{\omega}_{f}\,(\Bar{R}\cdot \Bar{Y})\,\Tilde{F}^{(00)}\Bar{F}^{(11)} +8\eta^{2}\,m_{f}\,\xi\, \Tilde{\omega}_{f}\,(\Tilde{R}\cdot \Bar{Y})\,\Tilde{F}^{(10)}\,\Bar{F}^{(01)}-8\zeta^{2}\,m_{f}\,\xi\,\omega_{f}\,(X\cdot \Bar{R})\,\Tilde{F}^{(01)}\Bar{F}^{(10)}
    \end{align}
    \item Regular $\times$ Instantaneous $+$ h.c.
    \begin{align}
    \mathcal{H}^{x}_{\mathrm{reg.}\times \mathrm{inst.}} & = -\frac{8\xi\,\eta (\xi +\zeta)^{2}}{(1-\zeta)}\,(\Tilde{R}\cdot X)\,\Tilde{F}^{(11)}\,\Bar{\mathcal{G}}_{E} + \frac{8\,\xi\,\zeta\,(1-\zeta)^{2}}{(\xi + \zeta)}\,(\Bar{R}\cdot \Bar{Y})\,\mathcal{G}_{E}\,\Bar{F}^{(11)}
    \end{align}
\item Instantaneous $\times$ Instantaneous
    \begin{align}
    \mathcal{H}^{x}_{\mathrm{inst.}\times \mathrm{inst.}} & = \frac{\zeta\,\xi}{(\zeta + \xi)}\,\frac{\xi\,\eta}{(1-\zeta)}\mathcal{G}_{E}\,\Bar{\mathcal{G}_{E}}\,\sum_{\lambda_{1}}\sum_{l,r}\,(\delta^{lr}-2i\epsilon^{lr}\lambda_{1})\,(\delta^{lr}+2i\epsilon^{lr}\lambda_{1})^{*} = 0
    \end{align}
\end{itemize}
$\mathcal{W}(x,y,z,\Bar{x},\Bar{y},\Bar{z})$ is the Wilson line structure. This has 4 terms:
\begin{itemize}
    \item $\mathcal{W}(x,y,z,\Bar{x},\Bar{y},\Bar{z})$
    \begin{align}
        \mathcal{W}(x,y,z,\Bar{x},\Bar{y},\Bar{z}) & = S_{q\Bar{q}g\,q\Bar{q}g}(x,y,z,\Bar{x},\Bar{y},\Bar{z}) -S_{q\Bar{q}g}(x,y,z) - S^{\dagger}_{q\Bar{q}g}(\Bar{x},\Bar{y},\Bar{z}) + 1
    \end{align}
At large $N_{c}$ limit, we get
$$S_{q\Bar{q}g\,q\Bar{q}g}(x,y,z,\Bar{x},\Bar{y},\Bar{z}) = Q(x,z,\Bar{z},\Bar{x})\, Q(z,y,\Bar{y},\Bar{z})$$
$$S_{q\Bar{q}g}(x,y,z) = S(x,z)\, S(z,y)$$
$$S^{\dagger}_{q\Bar{q}g}(\Bar{x},\Bar{y},\Bar{z}) = S(\Bar{z},\Bar{x})\, S(\Bar{y},\Bar{z})$$
\boxed{$$\mathcal{W}(x,y,z,\Bar{x},\Bar{y},\Bar{z}) = Q(x,z,\Bar{z},\Bar{x})\, Q(z,y,\Bar{y},\Bar{z}) - S(x,z)\, S(z,y) - S(\Bar{z},\Bar{x})\, S(\Bar{y},\Bar{z}) +1 $$}\\ \\ 
Using Gaussian approximation, we get
\begin{align}
    \mathcal{W}(x,y,z,\Bar{x},\Bar{y},\Bar{z}) & = \bigg[S(x,z)\,S(\Bar{z},\Bar{x}) - \frac{F(x,z;\Bar{z},\Bar{x})}{F(x,\Bar{z};z,\Bar{x})}\, \left[S(x,z)\,S(\Bar{z},\Bar{x})-S(x,\Bar{x})\, S(\Bar{z},z)\right]\bigg]\notag \\ &  \bigg[S(z,y)\,S(\Bar{y},\Bar{z}) - \frac{F(z,y;\Bar{y},\Bar{z})}{F(z,\Bar{y};y,\Bar{z})}\,\left[S(z,y)\,S(\Bar{y},\Bar{z})-S(z,\Bar{z})\,S(\Bar{y},y)\right]\bigg] - S(x,z)\, S(z,y) - S(\Bar{z},\Bar{x})\, S(\Bar{y},\Bar{z}) +1
\end{align}
\item $\mathcal{W}(x',y,x',\Bar{x},\Bar{y},\Bar{z})$
\begin{align}
        \mathcal{W}(x',y,x',\Bar{x},\Bar{y},\Bar{z}) & = S_{q\Bar{q}g\,q\Bar{q}g}(x',y,x',\Bar{x},\Bar{y},\Bar{z}) -S_{q\Bar{q}g}(x',y,x') - S^{\dagger}_{q\Bar{q}g}(\Bar{x},\Bar{y},\Bar{z}) + 1
    \end{align}
At large $N_{c}$ limit, we get
$$S_{q\Bar{q}g\,q\Bar{q}g}(x',y,x',\Bar{x},\Bar{y},\Bar{z}) = Q(x',x',\Bar{z},\Bar{x})\, Q(x',y,\Bar{y},\Bar{z}) = S(\Bar{z},\Bar{x})\,Q(x',y,\Bar{y},\Bar{z})$$
$$S_{q\Bar{q}g}(x',y,x') = S(x',x')\, S(x',y) = S(x',y) $$
$$S^{\dagger}_{q\Bar{q}g}(\Bar{x},\Bar{y},\Bar{z}) = S(\Bar{z},\Bar{x})\, S(\Bar{y},\Bar{z})$$
\boxed{$$\mathcal{W}(x',y,x',\Bar{x},\Bar{y},\Bar{z}) = S(\Bar{z},\Bar{x})\,Q(x',y,\Bar{y},\Bar{z}) - S(x',y) - S(\Bar{z},\Bar{x})\, S(\Bar{y},\Bar{z}) + 1   $$}\\ \\
Using Gaussian approximation, we get
\begin{align}
   \mathcal{W}(x',y,x',\Bar{x},\Bar{y},\Bar{z}) & = S(\Bar{z},\Bar{x})\,\bigg[S(x',y)\,S(\Bar{y},\Bar{z}) - \frac{F(x',y;\Bar{y},\Bar{z})}{F(x',\Bar{y};y,\Bar{z})}\,\left[S(x',y)\,S(\Bar{y},\Bar{z})-S(x',\Bar{z})\,S(\Bar{y},y)\right]\bigg] - S(x',y) - S(\Bar{z},\Bar{x})\, S(\Bar{y},\Bar{z}) + 1  
\end{align}
\item $\mathcal{W}(x,y,z,\Bar{x},\Bar{y}',\Bar{y}')$
\begin{align}
        \mathcal{W}(x,y,z,\Bar{x},\Bar{y}',\Bar{y}') & = S_{q\Bar{q}g\,q\Bar{q}g}(x,y,z,\Bar{x},\Bar{y}',\Bar{y}') -S_{q\Bar{q}g}(x,y,z) - S^{\dagger}_{q\Bar{q}g}(\Bar{x},\Bar{y}',\Bar{y}') + 1
    \end{align}
At large $N_{c}$ limit, we get
$$S_{q\Bar{q}g\,q\Bar{q}g}(x,y,z,\Bar{x},\Bar{y}',\Bar{y}') = Q(x,z,\Bar{y}',\Bar{x})\, Q(z,y,\Bar{y}',\Bar{y}') = Q(x,z,\Bar{y}',\Bar{x})\, S(z,y) $$
$$S_{q\Bar{q}g}(x,y,z) = S(x,z)\, S(z,y)$$
$$S^{\dagger}_{q\Bar{q}g}(\Bar{x},\Bar{y}',\Bar{y}') = S(\Bar{y}',\Bar{x})\, S(\Bar{y}',\Bar{y}') = S(\Bar{y}',\Bar{x})$$
\boxed{$$\mathcal{W}(x,y,z,\Bar{x},\Bar{y}',\Bar{y}') = Q(x,z,\Bar{y}',\Bar{x})\, S(z,y) -  S(x,z)\, S(z,y) - S(\Bar{y}',\Bar{x}) +1 $$}\\ \\
Using Gaussian approximation, we get
\begin{align}
    \mathcal{W}(x,y,z,\Bar{x},\Bar{y}',\Bar{y}') & = \bigg[S(x,z)\,S(\Bar{y}',\Bar{x}) - \frac{F(x,z;\Bar{y}',\Bar{x})}{F(x,\Bar{y}';z,\Bar{x})}\, \left[S(x,z)\,S(\Bar{y}',\Bar{x})-S(x,\Bar{x})\,S(\Bar{y}',z)\right]\bigg] S(z,y) -  S(x,z)\, S(z,y) - S(\Bar{y}',\Bar{x}) +1 
\end{align}
\item $\mathcal{W}(x',y,x',\Bar{x},\Bar{y}',\Bar{y}')$
    \begin{align}
        \mathcal{W}(x',y,x',\Bar{x},\Bar{y}',\Bar{y}') & = S_{q\Bar{q}g\,q\Bar{q}g}(x',y,x',\Bar{x},\Bar{y}',\Bar{y}') -S_{q\Bar{q}g}(x',y,x') - S^{\dagger}_{q\Bar{q}g}(\Bar{x},\Bar{y}',\Bar{y}') + 1
    \end{align}
At large $N_{c}$ limit, we get
$$S_{q\Bar{q}g\,q\Bar{q}g}(x',y,x',\Bar{x},\Bar{y}',\Bar{y}') = Q(x',x',\Bar{y}',\Bar{x})\, Q(x',y,\Bar{y}',\Bar{y}') = S(\Bar{y}',\Bar{x})\, S(x',y) $$
$$S_{q\Bar{q}g}(x',y,x') = S(x',x')\, S(x',y) = S(x',y)$$
$$S^{\dagger}_{q\Bar{q}g}(\Bar{x},\Bar{y}',\Bar{y}') = S(\Bar{y}',\Bar{x})\, S(\Bar{y}',\Bar{y}') = S(\Bar{y}',\Bar{x}) $$
\boxed{$$\mathcal{W}(x',y,x',\Bar{x},\Bar{y}',\Bar{y}') = S(\Bar{y}',\Bar{x})\, S(x',y) - S(x',y) - S(\Bar{y}',\Bar{x}) +1   $$}\\ \\
Using Gaussian approximation, we get 
\begin{align}
    \mathcal{W}(x',y,x',\Bar{x},\Bar{y}',\Bar{y}') & = S(\Bar{y}',\Bar{x})\, S(x',y) - S(x',y) - S(\Bar{y}',\Bar{x}) +1
\end{align}
\end{itemize}
\begin{comment}
\section{Fourier transforms}
Fourier transform is defined as: 
 \begin{align}\label{eq:Ddef}
        S(a,b) = \int_{l} e^{i\,l\cdot(a-b)}\,\mathcal{D}(l),\hspace{2em} \mathcal{D}(l) = \int d^{2}r\, e^{-i\,l\cdot r}\,S(r),\hspace{2em}\int_{l}\mathcal{D}(l) = S(0) = 1 .
    \end{align}
\textbf{Antiquark emission.} With $c=\zeta x+\eta y+\xi z$, $R=x-y'$, $Y=y-z$ (Jacobian $=1$),
\begin{align}
    x = c+(1-\zeta)R,\hspace{2em} y = c-\zeta R+\frac{\xi}{1-\zeta}Y,\hspace{2em} z = c-\zeta R-\frac{\eta}{1-\zeta}Y,\hspace{2em} y' = c-\zeta R,
\end{align}
so that the dipole arguments are
\begin{align}\label{eq:argqb}
    x-z = R+\frac{\eta}{1-\zeta}\,Y,\hspace{2em} z-y=-Y,\hspace{2em} x-y'=R,
\end{align}
and the phase of the cross section separates as
\begin{align}\label{eq:phaseqb}
    k_{1}\cdot x+k_{2}\cdot y+k_{3}\cdot z = K\cdot c+\Tilde{k}\cdot R-\Tilde{p}\cdot Y,\hspace{1.5em} K=k_{1}+k_{2}+k_{3},\hspace{1.5em} \Tilde{k}=(1-\zeta)k_{1}-\zeta(k_{2}+k_{3}),\hspace{1.5em} \Tilde{p}=\frac{\eta\,k_{3}-\xi\,k_{2}}{1-\zeta}.
\end{align}
\textbf{Quark emission.} With $c$, $\Tilde{R}=x'-y$, $X=x-z$, $x'=\frac{\zeta x+\xi z}{\zeta+\xi}$ (Jacobian $=1$),
\begin{align}
    y = c-(1-\eta)\Tilde{R},\hspace{2em} x = c+\eta\Tilde{R}+\frac{\xi}{1-\eta}X,\hspace{2em} z = c+\eta\Tilde{R}-\frac{\zeta}{1-\eta}X,\hspace{2em} x'=c+\eta\Tilde{R},
\end{align}
\begin{align}\label{eq:argq}
    x-z = X,\hspace{2em} z-y = \Tilde{R}-\frac{\zeta}{1-\eta}\,X,\hspace{2em} x'-y=\Tilde{R},
\end{align}
\begin{align}\label{eq:phaseq}
    k_{1}\cdot x+k_{2}\cdot y+k_{3}\cdot z = K\cdot c+\Tilde{k}_{q}\cdot\Tilde{R}-\Tilde{p}_{q}\cdot X,\hspace{2em} \Tilde{k}_{q}=\eta\,(k_{1}+k_{3})-(1-\eta)\,k_{2},\hspace{2em} \Tilde{p}_{q}=\frac{\zeta\,k_{3}-\xi\,k_{1}}{1-\eta}.
\end{align}
Inverting the Fourier transforms that defined $\mathcal{G},\mathcal{A},\mathcal{G}_{E}$ in Sec.~2,
\begin{align}
    \mathcal{M}^{lr}_{A\,\lambda_{1}\lambda_{2}}(R,Y) = \int\frac{d^{2}\Tilde{k}}{2\pi}\frac{d^{2}\Tilde{p}}{2\pi}\,e^{i\Tilde{k}\cdot R-i\Tilde{p}\cdot Y}\,\Phi^{lr}_{A\,\lambda_{1}\lambda_{2}}(\Tilde{k},\Tilde{p}),\hspace{2em}\int d^{2}R\,d^{2}Y\,e^{-i\Tilde{k}\cdot R+i\Tilde{p}\cdot Y}\,\mathcal{M}_{A}(R,Y) = (2\pi)^{2}\,\Phi_{A}(\Tilde{k},\Tilde{p}),
\end{align}
with $A=1$ (before the shock wave, regular $+$ instantaneous) and $A=2$ (after the shock wave):
\begin{align}\label{eq:Phiqb}
    \Phi_{1}(\Tilde{k},\Tilde{p}) & = \frac{1}{D(\Tilde{k},\Tilde{p})}\left[\frac{(1-\zeta)\,N(\Tilde{k},\Tilde{p})}{\Tilde{k}^{2}+\epsilon^{2}_{f}}-\xi\,\eta\,I\right], \hspace{3em} \Phi_{2}(\Tilde{k},\Tilde{p}) = \frac{N(\Tilde{k},\Tilde{p})}{\zeta(1-\zeta)\,(\Tilde{k}^{2}+\epsilon^{2}_{f})\,(\Tilde{p}^{2}+\omega^{2}_{f})},\notag \\
    D(\Tilde{k},\Tilde{p}) & = \zeta(1-\zeta)^{2}(\Tilde{p}^{2}+\omega^{2}_{f})+\xi\,\eta\,(\Tilde{k}^{2}+\epsilon^{2}_{f}),\notag \\
    N^{lr}_{\lambda_{1}\lambda_{2}}(\Tilde{k},\Tilde{p}) & = \sum_{\lambda}\left[\varphi^{lj}_{\lambda_{1}\lambda}(\zeta)\Tilde{k}^{j}+m_{f}\,\mu^{l}_{\lambda_{1}\lambda}\right]\left[\tau^{rn}_{\lambda\lambda_{2}}(\xi,\eta)\Tilde{p}^{n}+\xi\,\omega_{f}\,\mu^{r}_{\lambda\lambda_{2}}\right],\hspace{2em} I^{lr}_{\lambda_{1}\lambda_{2}} = \left(\delta^{lr}+2i\epsilon^{lr}\lambda_{1}\right)\delta_{\lambda_{1}\lambda_{2}}.
\end{align}
For the quark emission ($\mathcal{\Tilde{M}}_{A}(\Tilde{R},X)=\int\frac{d^{2}\Tilde{k}_{q}}{2\pi}\frac{d^{2}\Tilde{p}_{q}}{2\pi}e^{i\Tilde{k}_{q}\cdot\Tilde{R}-i\Tilde{p}_{q}\cdot X}\,\Tilde{\Phi}_{A}$):
\begin{align}\label{eq:Phiq}
    \Tilde{\Phi}_{1}(\Tilde{k}_{q},\Tilde{p}_{q}) & = \frac{1}{\Tilde{D}}\left[\frac{(1-\eta)\,\Tilde{N}}{\Tilde{k}^{2}_{q}+\Tilde{\epsilon}^{2}_{f}}-\xi\,\zeta\,\Tilde{I}\right], \hspace{3em} \Tilde{\Phi}_{2}(\Tilde{k}_{q},\Tilde{p}_{q}) = \frac{\Tilde{N}}{\eta(1-\eta)\,(\Tilde{k}^{2}_{q}+\Tilde{\epsilon}^{2}_{f})\,(\Tilde{p}^{2}_{q}+\Tilde{\omega}^{2}_{f})},\notag \\
    \Tilde{D} & = \eta(1-\eta)^{2}(\Tilde{p}^{2}_{q}+\Tilde{\omega}^{2}_{f})+\xi\,\zeta\,(\Tilde{k}^{2}_{q}+\Tilde{\epsilon}^{2}_{f}),\notag \\
    \Tilde{N}^{lr}_{\lambda_{1}\lambda_{2}} & = \sum_{\lambda}\left[\tau^{nr}_{\lambda_{1}\lambda}(\xi,\zeta)\Tilde{p}^{n}_{q}+\xi\,\Tilde{\omega}_{f}\,\mu^{r}_{\lambda_{1}\lambda}\right]\left[\varphi^{lj}_{\lambda\lambda_{2}}(\zeta+\xi)\Tilde{k}^{j}_{q}+m_{f}\,\mu^{l}_{\lambda\lambda_{2}}\right],\hspace{2em} \Tilde{I}^{lr}_{\lambda_{1}\lambda_{2}} = \left(\delta^{lr}-2i\epsilon^{lr}\lambda_{1}\right)\delta_{\lambda_{1}\lambda_{2}}.
\end{align}
\begin{align}\label{eq:Phi0}
    \Phi_{0}\equiv\Phi_{1}-\Phi_{2} = -\frac{\xi\,\eta}{D}\left[\frac{N}{\zeta(1-\zeta)(\Tilde{p}^{2}+\omega^{2}_{f})}+I\right],\hspace{3em} \Tilde{\Phi}_{0}\equiv\Tilde{\Phi}_{1}-\Tilde{\Phi}_{2} = -\frac{\xi\,\zeta}{\Tilde{D}}\left[\frac{\Tilde{N}}{\eta(1-\eta)(\Tilde{p}^{2}_{q}+\Tilde{\omega}^{2}_{f})}+\Tilde{I}\right].
\end{align}
All the transforms needed are therefore
\begin{align}\label{eq:Omega}
    \Omega_{1} & \equiv \mathcal{O}_{1}[S(x,z)S(z,y)] = (2\pi)^{2}\int_{l_{1}}\int_{l_{2}}\mathcal{D}(l_{1})\mathcal{D}(l_{2})\;\Phi_{1}\!\left(\Tilde{k}-l_{1},\ \Tilde{p}+\tfrac{\eta}{1-\zeta}l_{1}-l_{2}\right), \notag \\
    \Omega_{2} & \equiv\mathcal{O}_{2}[S(x,y')] = (2\pi)^{2}\int_{l}\mathcal{D}(l)\;\Phi_{2}\!\left(\Tilde{k}-l,\ \Tilde{p}\right),\notag \\
    \Tilde{\Omega}_{1} & \equiv \Tilde{\mathcal{O}}_{1}[S(y,z)S(z,x)] = (2\pi)^{2}\int_{l_{1}}\int_{l_{2}}\mathcal{D}(l_{1})\mathcal{D}(l_{2})\;\Tilde{\Phi}_{1}\!\left(\Tilde{k}_{q}+l_{1},\ \Tilde{p}_{q}+\tfrac{\zeta}{1-\eta}l_{1}-l_{2}\right), \notag \\
    \Tilde{\Omega}_{2} & \equiv\Tilde{\mathcal{O}}_{2}[S(y,x')] = (2\pi)^{2}\int_{l}\mathcal{D}(l)\;\Tilde{\Phi}_{2}\!\left(\Tilde{k}_{q}+l,\ \Tilde{p}_{q}\right),\notag \\
    \Omega^{x}_{1} & \equiv \Tilde{\mathcal{O}}_{1}[S(x,z)S(z,y)] = (2\pi)^{2}\int_{l_{1}}\int_{l_{2}}\mathcal{D}(l_{1})\mathcal{D}(l_{2})\;\Tilde{\Phi}_{1}\!\left(\Tilde{k}_{q}-l_{2},\ \Tilde{p}_{q}+l_{1}-\tfrac{\zeta}{1-\eta}l_{2}\right), \notag \\
    \Omega^{x}_{2} & \equiv\Tilde{\mathcal{O}}_{2}[S(x',y)] = (2\pi)^{2}\int_{l}\mathcal{D}(l)\;\Tilde{\Phi}_{2}\!\left(\Tilde{k}_{q}-l,\ \Tilde{p}_{q}\right),
\end{align}
where in $\Omega_{1}$ and $\Tilde{\Omega}_{1}$, $l_{1}$ belongs to the first dipole and $l_{2}$ to the second. Now, we will move on to the Helicity sums
\noindent\textbf{Antiquark emission} ($N$, $I$):
\begin{align}\label{eq:Kqb}
    \mathcal{K}(a;b) & \equiv\sum N(a)N^{*}(b) = 8\left[\zeta^{2}+(1-\zeta)^{2}\right]\left[\eta^{2}+(1-\zeta)^{2}\right](k_{a}\cdot k_{b})(p_{a}\cdot p_{b})-8\xi(2\zeta-1)\left(2(1-\zeta)-\xi\right)(k_{a}\times k_{b})(p_{a}\times p_{b})\notag \\
    & + 8m^{2}_{f}\left[\eta^{2}+(1-\zeta)^{2}\right](p_{a}\cdot p_{b}) + 8\xi^{2}\omega^{2}_{f}\left[\zeta^{2}+(1-\zeta)^{2}\right](k_{a}\cdot k_{b}) + 8m^{2}_{f}\xi^{2}\omega^{2}_{f} - 8\zeta\eta\,m_{f}\xi\omega_{f}\,(k_{a}-k_{b})\cdot(p_{a}-p_{b}),\notag \\
    \mathcal{J}(a) & \equiv\sum N(a)I^{*} = 8\zeta\eta\,(k_{a}\cdot p_{a})-8m_{f}\xi\omega_{f},\hspace{3em}\sum I\,I^{*}=8 .
\end{align}
\textbf{Quark emission} ($\Tilde{N}$, $\Tilde{I}$):
\begin{align}\label{eq:Kq}
    \Tilde{\mathcal{K}}(a;b) & = 8\left[\zeta^{2}+(1-\eta)^{2}\right]\left[\eta^{2}+(1-\eta)^{2}\right](k_{a}\cdot k_{b})(p_{a}\cdot p_{b})-8\xi(2\eta-1)\left(2(1-\eta)-\xi\right)(k_{a}\times k_{b})(p_{a}\times p_{b})\notag \\
    & + 8m^{2}_{f}\left[\zeta^{2}+(1-\eta)^{2}\right](p_{a}\cdot p_{b}) + 8\xi^{2}\Tilde{\omega}^{2}_{f}\left[\eta^{2}+(1-\eta)^{2}\right](k_{a}\cdot k_{b}) + 8m^{2}_{f}\xi^{2}\Tilde{\omega}^{2}_{f} + 8\zeta\eta\,m_{f}\xi\Tilde{\omega}_{f}\,(k_{a}-k_{b})\cdot(p_{a}-p_{b}),\notag \\
    \Tilde{\mathcal{J}}(a) & = -8\zeta\eta\,(k_{a}\cdot p_{a})-8m_{f}\xi\Tilde{\omega}_{f},\hspace{3em}\sum\Tilde{I}\,\Tilde{I}^{*}=8 .
\end{align}
\textbf{Interference} ($\Tilde{N}(a)$ with $N^{*}(b)$):
\begin{align}\label{eq:Kx}
    \mathcal{K}^{x}(a;b) & = 8\left[1-\xi-2\zeta\eta\right]\left[\xi(1-\xi)+2\zeta\eta\right](k_{a}\cdot k_{b})(p_{a}\cdot p_{b})-8\xi\left[1-2\zeta-\xi\right]^{2}(k_{a}\times k_{b})(p_{a}\times p_{b})\notag \\
    & +8m^{2}_{f}(\xi+2\zeta\eta)(p_{a}\cdot p_{b})-8(\xi+2\zeta\eta)\,\xi^{2}\Tilde{\omega}_{f}\omega_{f}\,(k_{a}\cdot k_{b})-8(\zeta+\xi)^{2}m_{f}\xi\omega_{f}\,(k_{a}\cdot p_{a})\notag \\
    & +8(1-\zeta)^{2}m_{f}\xi\Tilde{\omega}_{f}\,(k_{b}\cdot p_{b})-8\eta^{2}m_{f}\xi\Tilde{\omega}_{f}\,(k_{a}\cdot p_{b})+8\zeta^{2}m_{f}\xi\omega_{f}\,(p_{a}\cdot k_{b}),\notag \\
    \mathcal{J}^{x}(a) & \equiv\sum\Tilde{N}(a)I^{*} = 8(\zeta+\xi)^{2}(k_{a}\cdot p_{a}),\hspace{2em} \Bar{\mathcal{J}}^{x}(b)\equiv\sum\Tilde{I}\,N^{*}(b) = -8(1-\zeta)^{2}(k_{b}\cdot p_{b}),\hspace{2em}\sum\Tilde{I}\,I^{*}=0 .
\end{align}
We will have the following notations: 
\begin{align}\label{eq:hqb}
    \mathfrak{h}_{11}(a;b) & = \frac{1}{D_{a}D_{b}}\left[\frac{(1-\zeta)^{2}\mathcal{K}(a;b)}{A_{a}A_{b}}-(1-\zeta)\,\xi\eta\left(\frac{\mathcal{J}(a)}{A_{a}}+\frac{\mathcal{J}(b)}{A_{b}}\right)+8\xi^{2}\eta^{2}\right],\hspace{2em} \mathfrak{h}_{22}(a;b) = \frac{\mathcal{K}(a;b)}{\zeta^{2}(1-\zeta)^{2}A_{a}B_{a}A_{b}B_{b}},\notag \\
    \mathfrak{h}_{12}(a;b) & = \frac{1}{\zeta(1-\zeta)D_{a}A_{b}B_{b}}\left[\frac{(1-\zeta)\mathcal{K}(a;b)}{A_{a}}-\xi\eta\,\mathcal{J}(b)\right],\hspace{2em} \mathfrak{h}_{21}(a;b) = \frac{1}{\zeta(1-\zeta)D_{b}A_{a}B_{a}}\left[\frac{(1-\zeta)\mathcal{K}(a;b)}{A_{b}}-\xi\eta\,\mathcal{J}(a)\right].
\end{align}
\begin{align}\label{eq:hx}
    \mathfrak{h}^{x}_{11}(a;b) & = \frac{1}{\Tilde{D}_{a}D_{b}}\left[\frac{(1-\eta)(1-\zeta)\mathcal{K}^{x}(a;b)}{\Tilde{A}_{a}A_{b}}-\frac{(1-\eta)\,\xi\eta\,\mathcal{J}^{x}(a)}{\Tilde{A}_{a}}-\frac{(1-\zeta)\,\xi\zeta\,\Bar{\mathcal{J}}^{x}(b)}{A_{b}}\right],\hspace{2em} \mathfrak{h}^{x}_{22}(a;b) = \frac{\mathcal{K}^{x}(a;b)}{\eta(1-\eta)\zeta(1-\zeta)\Tilde{A}_{a}\Tilde{B}_{a}A_{b}B_{b}},\notag \\
    \mathfrak{h}^{x}_{12}(a;b) & = \frac{1}{\zeta(1-\zeta)\Tilde{D}_{a}A_{b}B_{b}}\left[\frac{(1-\eta)\mathcal{K}^{x}(a;b)}{\Tilde{A}_{a}}-\xi\zeta\,\Bar{\mathcal{J}}^{x}(b)\right],\hspace{2em} \mathfrak{h}^{x}_{21}(a;b) = \frac{1}{\eta(1-\eta)D_{b}\Tilde{A}_{a}\Tilde{B}_{a}}\left[\frac{(1-\zeta)\mathcal{K}^{x}(a;b)}{A_{b}}-\xi\eta\,\mathcal{J}^{x}(a)\right].
\end{align}
These are the final results: \\
Direct, anti quark emission:
\begin{align}
    -S(x,z)S(z,y)\  & :\ -\!\int_{l_{1}l_{2}}\!\mathcal{D}\mathcal{D}\,\mathfrak{h}_{11}(a_{l_{1}l_{2}};b), & -S(\Bar{z},\Bar{x})S(\Bar{y},\Bar{z})\  & :\ -\!\int_{l_{1}l_{2}}\!\mathcal{D}\mathcal{D}\,\mathfrak{h}_{11}(b;a_{l_{1}l_{2}}),\notag \\
    -S(x,y')\ & :\ +\!\int_{l}\!\mathcal{D}\,\mathfrak{h}_{21}(a_{l};b), & -S(\Bar{z},\Bar{x})S(\Bar{y},\Bar{z})\  & :\ +\!\int_{l_{1}l_{2}}\!\mathcal{D}\mathcal{D}\,\mathfrak{h}_{21}(b;a_{l_{1}l_{2}}),\notag \\
    -S(x,z)S(z,y)\  & :\ +\!\int_{l_{1}l_{2}}\!\mathcal{D}\mathcal{D}\,\mathfrak{h}_{12}(a_{l_{1}l_{2}};b), & -S(\Bar{y}',\Bar{x})\  & :\ +\!\int_{l}\!\mathcal{D}\,\mathfrak{h}_{12}(b;a_{l}),\notag \\
    -S(x,y')\  & :\ -\!\int_{l}\!\mathcal{D}\,\mathfrak{h}_{22}(a_{l};b), & -S(\Bar{y}',\Bar{x})\  & :\ -\!\int_{l}\!\mathcal{D}\,\mathfrak{h}_{22}(b;a_{l}).
\end{align}
where $\int_{l_{1}l_{2}}\mathcal{D}\mathcal{D}\equiv\int_{l_{1}}\int_{l_{2}}\mathcal{D}(l_{1})\mathcal{D}(l_{2})$.
\begin{itemize}
\item Direct contribution antiquark
\begin{align}\label{eq:resqb}
    \frac{d^{3}N^{\text{direct}}}{d^{3}k_{1}\,d^{3}k_{2}\, d^{3}k_{3}} & = \frac{e^{2}\,g^{2}\,\zeta^{2}\,C_{F}\,N_{c}\;S_{\perp}}{8\,(2\pi)^{6}\,\xi\,q^{+\,2}}\,\delta\left(q^{+}-k^{+}_{1}-k^{+}_{2}-k^{+}_{3}\right)\delta^{(2)}(k_{1}+k_{2}+k_{3})\sum_{f}e^{2}_{f}\notag \\
    & \times(-2)\left\{\int_{l_{1}}\int_{l_{2}}\mathcal{D}(l_{1})\mathcal{D}(l_{2})\Big[\mathfrak{h}_{11}-\mathfrak{h}_{12}\Big](a_{l_{1}l_{2}};b)-\int_{l}\mathcal{D}(l)\Big[\mathfrak{h}_{21}-\mathfrak{h}_{22}\Big](a_{l};b)\right\}
\end{align}
\item Direct contribution quark
\begin{align}\label{eq:resq}
    \frac{d^{3}N^{\text{direct}}}{d^{3}k_{1}\,d^{3}k_{2}\, d^{3}k_{3}} & = \frac{e^{2}\,g^{2}\,\eta^{2}\,C_{F}\,N_{c}\;S_{\perp}}{8\,(2\pi)^{6}\,\xi\,q^{+\,2}}\,\delta\left(q^{+}-k^{+}_{1}-k^{+}_{2}-k^{+}_{3}\right)\delta^{(2)}(k_{1}+k_{2}+k_{3})\sum_{f}e^{2}_{f}\notag \\
    & \times(-2)\left\{\int_{l_{1}}\int_{l_{2}}\mathcal{D}(l_{1})\mathcal{D}(l_{2})\Big[\Tilde{\mathfrak{h}}_{11}-\Tilde{\mathfrak{h}}_{12}\Big](\Tilde{a}_{l_{1}l_{2}};\Tilde{b})-\int_{l}\mathcal{D}(l)\Big[\Tilde{\mathfrak{h}}_{21}-\Tilde{\mathfrak{h}}_{22}\Big](\Tilde{a}_{l};\Tilde{b})\right\},
\end{align}
\item Interference contribution
\begin{align}
    -S(x,z)S(z,y)\  & :\ -\!\int_{l_{1}l_{2}}\!\mathcal{D}\mathcal{D}\,\mathfrak{h}^{x}_{11}(\hat{a}_{l_{1}l_{2}};b), \notag \\
    -S(\Bar{z},\Bar{x})S(\Bar{y},\Bar{z})\  & :\ -\!\int_{l_{1}l_{2}}\!\mathcal{D}\mathcal{D}\,\mathfrak{h}^{x}_{11}(\Tilde{b};a_{l_{1}l_{2}}),\notag \\
    -S(x',y)\  & :\ +\!\int_{l}\!\mathcal{D}\,\mathfrak{h}^{x}_{21}(\hat{a}_{l};b), \notag \\
    -S(\Bar{z},\Bar{x})S(\Bar{y},\Bar{z})\ & :\ +\!\int_{l_{1}l_{2}}\!\mathcal{D}\mathcal{D}\,\mathfrak{h}^{x}_{21}(\Tilde{b};a_{l_{1}l_{2}}),\notag \\
    -S(x,z)S(z,y)\ & :\ +\!\int_{l_{1}l_{2}}\!\mathcal{D}\mathcal{D}\,\mathfrak{h}^{x}_{12}(\hat{a}_{l_{1}l_{2}};b), \notag \\
    -S(\Bar{y}',\Bar{x})\  & :\ +\!\int_{l}\!\mathcal{D}\,\mathfrak{h}^{x}_{12}(\Tilde{b};a_{l}),\notag \\
    -S(x',y)\  & :\ -\!\int_{l}\!\mathcal{D}\,\mathfrak{h}^{x}_{22}(\hat{a}_{l};b), \notag \\
    -S(\Bar{y}',\Bar{x})\  & :\ -\!\int_{l}\!\mathcal{D}\,\mathfrak{h}^{x}_{22}(\Tilde{b};a_{l}),\notag \\
    +S(\Bar{y}',\Bar{x})S(x',y)\  & :\ +\!\int_{l}\int_{l'}\!\mathcal{D}(l)\mathcal{D}(l')\,\mathfrak{h}^{x}_{22}(\hat{a}_{l};a_{l'}),
\end{align}
\begin{align}\label{eq:resx}
    \frac{d^{3}N^{\text{cross}}}{d^{3}k_{1}\,d^{3}k_{2}\, d^{3}k_{3}} & = -\frac{e^{2}\,g^{2}\,\eta\,\zeta\,C_{F}\,N_{c}\;S_{\perp}}{8\,(2\pi)^{6}\,\xi\,q^{+\,2}}\,\delta\left(q^{+}-k^{+}_{1}-k^{+}_{2}-k^{+}_{3}\right)\delta^{(2)}(k_{1}+k_{2}+k_{3})\sum_{f}e^{2}_{f}\notag \\
    & \times2\,\Re\Bigg\{-\int_{l_{1}}\int_{l_{2}}\mathcal{D}(l_{1})\mathcal{D}(l_{2})\Big[\mathfrak{h}^{x}_{11}-\mathfrak{h}^{x}_{12}\Big](\hat{a}_{l_{1}l_{2}};b)\notag \\
    & +\int_{l}\mathcal{D}(l)\Big[\mathfrak{h}^{x}_{21}-\mathfrak{h}^{x}_{22}\Big](\hat{a}_{l};b) - \int_{l_{1}}\int_{l_{2}}\mathcal{D}(l_{1})\mathcal{D}(l_{2})\Big[\mathfrak{h}^{x}_{11}-\mathfrak{h}^{x}_{21}\Big](\Tilde{b};a_{l_{1}l_{2}}) + \int_{l}\mathcal{D}(l)\Big[\mathfrak{h}^{x}_{12}-\mathfrak{h}^{x}_{22}\Big](\Tilde{b};a_{l})\notag \\
    & + \int_{l}\int_{l'}\mathcal{D}(l)\mathcal{D}(l')\,\mathfrak{h}^{x}_{22}(\hat{a}_{l};a_{l'})\Bigg\}.
\end{align}
\end{itemize}
\end{comment}
\section{Fourier transform of the elastic terms}\label{sec:el}
\noindent\textit{New section.} In the Gaussian approximation each quadrupole is written as (elastic part) $-$ $\frac{F}{F}$ $\times$ (inelastic part), e.g.
$$Q(x,z,\Bar{z},\Bar{x}) = \underbrace{S(x,z)\,S(\Bar{z},\Bar{x})}_{\text{elastic}}\ -\ \underbrace{\frac{F(x,z;\Bar{z},\Bar{x})}{F(x,\Bar{z};z,\Bar{x})}\left[S(x,z)\,S(\Bar{z},\Bar{x})-S(x,\Bar{x})\,S(\Bar{z},z)\right]}_{\text{inelastic}} .$$
Here only the \emph{elastic} terms are Fourier transformed: every term of $\mathcal{W}$ that contains dipoles only (no ratio of $F$'s), together with the $+1$.

\subsection{Setting up the Fourier transform}
\textbf{Step 1: the elastic part of every $\mathcal{W}$.} Dropping the $F$-terms in the Gaussian expressions of Sec.~3, each $\mathcal{W}$ becomes a product of an amplitude factor and a conjugate-amplitude factor:
\begin{align}
\text{antiquark:}\hspace{1em} \mathcal{W}^{\text{el}}(x,y,z,\Bar{x},\Bar{y},\Bar{z}) & = S(x,z)S(\Bar{z},\Bar{x})\,S(z,y)S(\Bar{y},\Bar{z})-S(x,z)S(z,y)-S(\Bar{z},\Bar{x})S(\Bar{y},\Bar{z})+1\notag \\ & = \big[S(x,z)S(z,y)-1\big]\big[S(\Bar{z},\Bar{x})S(\Bar{y},\Bar{z})-1\big],\notag \\
\mathcal{W}^{\text{el}}(x,y',y',\Bar{x},\Bar{y},\Bar{z}) & = \big[S(x,y')-1\big]\big[S(\Bar{z},\Bar{x})S(\Bar{y},\Bar{z})-1\big],\notag \\
\mathcal{W}^{\text{el}}(x,y,z,\Bar{x},\Bar{y}',\Bar{y}') & = \big[S(x,z)S(z,y)-1\big]\big[S(\Bar{y}',\Bar{x})-1\big],\notag \\
\mathcal{W}^{\text{el}}(x,y',y',\Bar{x},\Bar{y}',\Bar{y}') & = \big[S(x,y')-1\big]\big[S(\Bar{y}',\Bar{x})-1\big],\notag \\
\text{quark:}\hspace{1em} \mathcal{W}^{\text{el}}(y,x,z,\Bar{y},\Bar{x},\Bar{z}) & = \big[S(y,z)S(z,x)-1\big]\big[S(\Bar{z},\Bar{y})S(\Bar{x},\Bar{z})-1\big],\notag \\
\mathcal{W}^{\text{el}}(y,x',x',\Bar{y},\Bar{x},\Bar{z}) & = \big[S(y,x')-1\big]\big[S(\Bar{z},\Bar{y})S(\Bar{x},\Bar{z})-1\big],\notag \\
\mathcal{W}^{\text{el}}(y,x,z,\Bar{y},\Bar{x}',\Bar{x}') & = \big[S(y,z)S(z,x)-1\big]\big[S(\Bar{x}',\Bar{y})-1\big],\notag \\
\mathcal{W}^{\text{el}}(y,x',x',\Bar{y},\Bar{x}',\Bar{x}') & = \big[S(y,x')-1\big]\big[S(\Bar{x}',\Bar{y})-1\big],\notag \\
\text{interference:}\hspace{1em} \mathcal{W}^{\text{el}}(x,y,z,\Bar{x},\Bar{y},\Bar{z}) & = \big[S(x,z)S(z,y)-1\big]\big[S(\Bar{z},\Bar{x})S(\Bar{y},\Bar{z})-1\big],\notag \\
\mathcal{W}^{\text{el}}(x',y,x',\Bar{x},\Bar{y},\Bar{z}) & = \big[S(x',y)-1\big]\big[S(\Bar{z},\Bar{x})S(\Bar{y},\Bar{z})-1\big],\notag \\
\mathcal{W}^{\text{el}}(x,y,z,\Bar{x},\Bar{y}',\Bar{y}') & = \big[S(x,z)S(z,y)-1\big]\big[S(\Bar{y}',\Bar{x})-1\big],\notag \\
\mathcal{W}^{\text{el}}(x',y,x',\Bar{x},\Bar{y}',\Bar{y}') & = \big[S(x',y)-1\big]\big[S(\Bar{y}',\Bar{x})-1\big].
\end{align}
(The last structure has no $F$-term at all, so it is entirely elastic.) Because of this factorisation, the four terms of each channel combine into $|\ldots|^{2}$. For the antiquark,
\begin{align}\label{eq:elcomb}
& \mathcal{H}[\mathcal{G};\Bar{\mathcal{G}}]\mathcal{W}^{\text{el}}(x,y,z;\ldots)-\mathcal{H}[\mathcal{A};\Bar{\mathcal{G}}]\mathcal{W}^{\text{el}}(x,y',y';\ldots)-\mathcal{H}[\mathcal{G};\Bar{\mathcal{A}}]\mathcal{W}^{\text{el}}(\ldots;\Bar{x},\Bar{y}',\Bar{y}')\notag \\
& +\mathcal{H}[\mathcal{A};\Bar{\mathcal{A}}]\mathcal{W}^{\text{el}}(x,y',y';\Bar{x},\Bar{y}',\Bar{y}')\notag \\
& = \sum_{\lambda_{1}\lambda_{2}}\sum_{l,r}\Big\{\mathcal{M}^{lr}_{1}(x,y,z)\big[S(x,z)S(z,y)-1\big]-\mathcal{M}^{lr}_{2}(x,y,z)\big[S(x,y')-1\big]\Big\}\notag \\
& \hspace{3em}\times\Big\{\mathcal{M}^{lr}_{1}(\Bar{x},\Bar{y},\Bar{z})\big[S(\Bar{x},\Bar{z})S(\Bar{z},\Bar{y})-1\big]-\mathcal{M}^{lr}_{2}(\Bar{x},\Bar{y},\Bar{z})\big[S(\Bar{x},\Bar{y}')-1\big]\Big\}^{*},
\end{align}
where we used $S(a,b)^{*}=S(b,a)$. For the quark channel replace $\mathcal{M}_{1,2}\rightarrow\Tilde{\mathcal{M}}_{1,2}$, $S(x,z)S(z,y)\rightarrow S(y,z)S(z,x)$ and $S(x,y')\rightarrow S(y,x')$ in both braces. For the interference, the first brace is the quark one with $S(x,z)S(z,y)$ and $S(x',y)$, and the second brace is the antiquark one.

\noindent\textbf{Step 2: the dipole in momentum space.} We take the dipole to depend only on the distance between its two points, $S(a,b)=S(|a-b|)$; this is already the case in the Gaussian expressions, where it enters through $\Gamma(x-z)$. Then
\begin{align}\label{eq:Dl}
S(a,b) = \int\frac{d^{2}l}{(2\pi)^{2}}\,e^{i\,l\cdot(a-b)}\,\mathcal{D}(l),\hspace{2em} \mathcal{D}(l) = \int d^{2}r\,e^{-i\,l\cdot r}\,S(r),\hspace{2em} \int\frac{d^{2}l}{(2\pi)^{2}}\,\mathcal{D}(l) = S(0) = 1 .
\end{align}
$\mathcal{D}(l)$ is real and depends only on $|l|$: it is the probability density for the dipole to receive a transverse momentum kick $l$ from the target. For example, $S(r)=e^{-Q^{2}_{s}r^{2}/4}$ gives $\mathcal{D}(l)=\frac{4\pi}{Q^{2}_{s}}e^{-l^{2}/Q^{2}_{s}}$. With $S(a,b)=S(b,a)$, the quark-channel ordering $S(y,z)S(z,x)$ of Sec.~3.2 equals $S(x,z)S(z,y)$, and $S(y,x')=S(x',y)$.

\noindent\textbf{Step 3: change of variables.} For the antiquark emission we use $c=\zeta x+\eta y+\xi z$, $R=x-y'$, $Y=y-z$. The Jacobian is $1$, and
\begin{align}
x = c+(1-\zeta)R,\hspace{1.5em} y = c-\zeta R+\frac{\xi}{1-\zeta}Y,\hspace{1.5em} z = c-\zeta R-\frac{\eta}{1-\zeta}Y,\hspace{1.5em} y' = c-\zeta R,\notag \\ \text{so}\hspace{1em} x-z = R+\frac{\eta}{1-\zeta}Y,\hspace{1em} z-y=-Y,\hspace{1em} x-y' = R .
\end{align}
For the quark emission we use $c$, $\Tilde{R}=x'-y$, $X=x-z$. The Jacobian is again $1$, and
\begin{align}
y = c-(1-\eta)\Tilde{R},\hspace{1.5em} x = c+\eta\Tilde{R}+\frac{\xi}{1-\eta}X,\hspace{1.5em} z = c+\eta\Tilde{R}-\frac{\zeta}{1-\eta}X,\hspace{1.5em} x' = c+\eta\Tilde{R},\notag \\ \text{so}\hspace{1em} x-z = X,\hspace{1em} z-y = \Tilde{R}-\frac{\zeta}{1-\eta}X,\hspace{1em} x'-y=\Tilde{R} .
\end{align}
The phase of the cross section becomes
\begin{align}\label{eq:phase}
k_{1}\cdot x+k_{2}\cdot y+k_{3}\cdot z = (k_{1}+k_{2}+k_{3})\cdot c+\Tilde{k}\cdot R-\Tilde{p}\cdot Y = (k_{1}+k_{2}+k_{3})\cdot c+\Tilde{k}_{q}\cdot\Tilde{R}-\Tilde{p}_{q}\cdot X,
\end{align}
with the relative momenta
\begin{align}\label{eq:relmom}
\Tilde{k} = (1-\zeta)k_{1}-\zeta(k_{2}+k_{3}),\hspace{1.5em} \Tilde{p}=\frac{\eta\,k_{3}-\xi\,k_{2}}{1-\zeta},\hspace{1.5em} \Tilde{k}_{q} = \eta(k_{1}+k_{3})-(1-\eta)k_{2},\hspace{1.5em} \Tilde{p}_{q} = \frac{\zeta\,k_{3}-\xi\,k_{1}}{1-\eta}.
\end{align}
In the elastic terms nothing depends on $c$ and $\Bar{c}$ except the phase. Hence $\int d^{2}c\,d^{2}\Bar{c}\,e^{-i(k_{1}+k_{2}+k_{3})\cdot(c-\Bar{c})} = S_{\perp}(2\pi)^{2}\delta^{(2)}(k_{1}+k_{2}+k_{3})$, where $S_{\perp}=\int d^{2}b$ is the transverse area of the target. (We keep, as in Sec.~3, photon transverse momentum $q=0$; for $q\neq0$ replace $\delta^{(2)}(k_{1}+k_{2}+k_{3})\rightarrow\delta^{(2)}(k_{1}+k_{2}+k_{3}-q)$.) The elastic terms therefore carry no net transverse momentum transfer from the target. On this support,
$$\Tilde{k}=k_{1},\hspace{2em}\Tilde{p}=k_{3}+\frac{\xi}{1-\zeta}\,k_{1},\hspace{2em}\Tilde{k}_{q}=-k_{2},\hspace{2em}\Tilde{p}_{q}=k_{3}+\frac{\xi}{1-\eta}\,k_{2}.$$

\noindent\textbf{Step 4: the wave functions in momentum space.} Undoing the Fourier transforms of Sec.~2 (which defined $\mathcal{G}^{(\alpha\beta)}$, $\mathcal{A}^{(\alpha\beta)}$, $\mathcal{G}_{E}$),
\begin{align}
\mathcal{M}^{lr}_{1,2}(R,Y) = \int\frac{d^{2}\Tilde{k}}{2\pi}\,\frac{d^{2}\Tilde{p}}{2\pi}\,e^{i\Tilde{k}\cdot R-i\Tilde{p}\cdot Y}\,\Phi^{lr}_{1,2}(\Tilde{k},\Tilde{p}),\hspace{2em} \int d^{2}R\,d^{2}Y\,e^{-i\Tilde{k}\cdot R+i\Tilde{p}\cdot Y}\,\mathcal{M}^{lr}_{1,2}(R,Y) = (2\pi)^{2}\,\Phi^{lr}_{1,2}(\Tilde{k},\Tilde{p}),
\end{align}
with (helicity indices $\lambda_{1}\lambda_{2}$ suppressed on the left)
\begin{align}\label{eq:Phiel}
\Phi^{lr}_{1}(k,p) & = \sum_{\lambda}\Big[\,\underbrace{\varphi^{lj}_{\lambda_{1}\lambda}(\zeta)\tau^{rn}_{\lambda\lambda_{2}}(\xi,\eta)\,k^{j}p^{n}}_{\text{no mass}}\notag \\ & \hspace{4em}+\underbrace{\xi\omega_{f}\,\varphi^{lj}_{\lambda_{1}\lambda}(\zeta)\mu^{r}_{\lambda\lambda_{2}}\,k^{j}+m_{f}\,\mu^{l}_{\lambda_{1}\lambda}\tau^{rn}_{\lambda\lambda_{2}}(\xi,\eta)\,p^{n}+m_{f}\xi\omega_{f}\,\mu^{l}_{\lambda_{1}\lambda}\mu^{r}_{\lambda\lambda_{2}}}_{\text{mass}}\Big]\,f_{1}(k,p)\notag \\ & -\xi\eta\left(\delta^{lr}+2i\epsilon^{lr}\lambda_{1}\right)\delta_{\lambda_{1}\lambda_{2}}\,g(k,p),\notag \\
\Phi^{lr}_{2}(k,p) & = \sum_{\lambda}\Big[\text{same bracket}\Big]\,f_{2}(k,p),
\end{align}
and the scalar functions
\begin{align}\label{eq:f12}
f_{1}(k,p) & = \frac{1-\zeta}{(k^{2}+\epsilon^{2}_{f})\left[\zeta(1-\zeta)^{2}(p^{2}+\omega^{2}_{f})+\xi\eta(k^{2}+\epsilon^{2}_{f})\right]},\hspace{2em} g(k,p) = \frac{1}{\zeta(1-\zeta)^{2}(p^{2}+\omega^{2}_{f})+\xi\eta(k^{2}+\epsilon^{2}_{f})},\notag \\
f_{2}(k,p) & = \frac{1}{\zeta(1-\zeta)(k^{2}+\epsilon^{2}_{f})(p^{2}+\omega^{2}_{f})},\hspace{2em} f_{1}(k,p)-f_{2}(k,p) = -\frac{\xi\eta}{\zeta(1-\zeta)(p^{2}+\omega^{2}_{f})\left[\zeta(1-\zeta)^{2}(p^{2}+\omega^{2}_{f})+\xi\eta(k^{2}+\epsilon^{2}_{f})\right]}.
\end{align}
Quark emission, with $\Tilde{\mathcal{M}}_{1,2}(\Tilde{R},X)=\int\frac{d^{2}\Tilde{k}_{q}}{2\pi}\frac{d^{2}\Tilde{p}_{q}}{2\pi}e^{i\Tilde{k}_{q}\cdot\Tilde{R}-i\Tilde{p}_{q}\cdot X}\,\Tilde{\Phi}_{1,2}(\Tilde{k}_{q},\Tilde{p}_{q})$:
\begin{align}\label{eq:Phielq}
\Tilde{\Phi}^{lr}_{1}(k,p) & = \sum_{\lambda}\Big[\underbrace{\tau^{nr}_{\lambda_{1}\lambda}(\xi,\zeta)\varphi^{lj}_{\lambda\lambda_{2}}(\zeta+\xi)\,k^{j}p^{n}}_{\text{no mass}}\notag \\ & \hspace{4em}+\underbrace{\xi\Tilde{\omega}_{f}\,\mu^{r}_{\lambda_{1}\lambda}\varphi^{lj}_{\lambda\lambda_{2}}(\zeta+\xi)\,k^{j}+m_{f}\,\tau^{nr}_{\lambda_{1}\lambda}(\xi,\zeta)\mu^{l}_{\lambda\lambda_{2}}\,p^{n}+m_{f}\xi\Tilde{\omega}_{f}\,\mu^{r}_{\lambda_{1}\lambda}\mu^{l}_{\lambda\lambda_{2}}}_{\text{mass}}\Big]\,\Tilde{f}_{1}(k,p)\notag \\ & -\xi\zeta\left(\delta^{lr}-2i\epsilon^{lr}\lambda_{1}\right)\delta_{\lambda_{1}\lambda_{2}}\,\Tilde{g}(k,p),\notag \\
\Tilde{\Phi}^{lr}_{2}(k,p) & = \sum_{\lambda}\Big[\text{same bracket}\Big]\,\Tilde{f}_{2}(k,p),
\end{align}
where $\Tilde{f}_{1},\Tilde{f}_{2},\Tilde{g}$ are $f_{1},f_{2},g$ with $\zeta\rightarrow\eta$ in the explicit factors ($1-\zeta\rightarrow1-\eta$, $\xi\eta\rightarrow\xi\zeta$, $\zeta(1-\zeta)\rightarrow\eta(1-\eta)$) and $\epsilon_{f}\rightarrow\Tilde{\epsilon}_{f}$, $\omega_{f}\rightarrow\Tilde{\omega}_{f}$:
\begin{gather*}\Tilde{f}_{1}(k,p) = \frac{1-\eta}{(k^{2}+\Tilde{\epsilon}^{2}_{f})\left[\eta(1-\eta)^{2}(p^{2}+\Tilde{\omega}^{2}_{f})+\xi\zeta(k^{2}+\Tilde{\epsilon}^{2}_{f})\right]},\hspace{1.5em} \Tilde{f}_{2}(k,p) = \frac{1}{\eta(1-\eta)(k^{2}+\Tilde{\epsilon}^{2}_{f})(p^{2}+\Tilde{\omega}^{2}_{f})},\\ \Tilde{g}(k,p) = \frac{1}{\eta(1-\eta)^{2}(p^{2}+\Tilde{\omega}^{2}_{f})+\xi\zeta(k^{2}+\Tilde{\epsilon}^{2}_{f})}.\end{gather*}

\subsection{Evaluating the Fourier transform}
\textbf{Step 5: the dipole integrals.} Insert Eq.~\eqref{eq:Dl} for each dipole. A dipole $S(a,b)$ with $a-b=\alpha R+\beta Y$ contributes $e^{il\cdot(\alpha R+\beta Y)}$, which combines with $e^{-i\Tilde{k}\cdot R+i\Tilde{p}\cdot Y}$. The $R$ and $Y$ integrals then give $(2\pi)^{2}\Phi$ at shifted momenta:
\begin{align}
\int d^{2}R\,d^{2}Y\,e^{-i\Tilde{k}\cdot R+i\Tilde{p}\cdot Y}\,\mathcal{M}_{1}\,S(x,z)S(z,y) & = (2\pi)^{2}\int\frac{d^{2}l_{1}}{(2\pi)^{2}}\frac{d^{2}l_{2}}{(2\pi)^{2}}\,\mathcal{D}(l_{1})\mathcal{D}(l_{2})\;\Phi_{1}\Big(\Tilde{k}-l_{1},\ \Tilde{p}+\frac{\eta}{1-\zeta}\,l_{1}-l_{2}\Big),\notag \\
\int d^{2}R\,d^{2}Y\,e^{-i\Tilde{k}\cdot R+i\Tilde{p}\cdot Y}\,\mathcal{M}_{2}\,S(x,y') & = (2\pi)^{2}\int\frac{d^{2}l_{1}}{(2\pi)^{2}}\,\mathcal{D}(l_{1})\;\Phi_{2}\Big(\Tilde{k}-l_{1},\ \Tilde{p}\Big),\notag \\
\int d^{2}\Tilde{R}\,d^{2}X\,e^{-i\Tilde{k}_{q}\cdot\Tilde{R}+i\Tilde{p}_{q}\cdot X}\,\Tilde{\mathcal{M}}_{1}\,S(x,z)S(z,y) & = (2\pi)^{2}\int\frac{d^{2}l_{1}}{(2\pi)^{2}}\frac{d^{2}l_{2}}{(2\pi)^{2}}\,\mathcal{D}(l_{1})\mathcal{D}(l_{2})\;\Tilde{\Phi}_{1}\Big(\Tilde{k}_{q}-l_{2},\ \Tilde{p}_{q}+l_{1}-\frac{\zeta}{1-\eta}\,l_{2}\Big),\notag \\
\int d^{2}\Tilde{R}\,d^{2}X\,e^{-i\Tilde{k}_{q}\cdot\Tilde{R}+i\Tilde{p}_{q}\cdot X}\,\Tilde{\mathcal{M}}_{2}\,S(x',y) & = (2\pi)^{2}\int\frac{d^{2}l_{2}}{(2\pi)^{2}}\,\mathcal{D}(l_{2})\;\Tilde{\Phi}_{2}\Big(\Tilde{k}_{q}-l_{2},\ \Tilde{p}_{q}\Big),
\end{align}
and the terms $-1$ give $-(2\pi)^{2}\Phi(\Tilde{k},\Tilde{p})$. Here $l_{1}$ is the momentum exchanged by the quark--gluon dipole $S(x,z)$ (or by $S(x,y')$), and $l_{2}$ the one exchanged by the gluon--antiquark dipole $S(z,y)$ (or by $S(x',y)$). The shifted arguments are just the relative momenta of the partons \emph{before} the scattering, $k_{1}-l_{1}$, $k_{2}+l_{2}$, $k_{3}+l_{1}-l_{2}$:
\begin{align}\label{eq:kappa}
\kappa & \equiv \Tilde{k}-l_{1} = k_{1}-l_{1}, & \rho & \equiv \Tilde{p}+\frac{\eta}{1-\zeta}\,l_{1}-l_{2} = (k_{3}+l_{1}-l_{2})+\frac{\xi}{1-\zeta}(k_{1}-l_{1}),\notag \\
\kappa_{q} & \equiv \Tilde{k}_{q}-l_{2} = -(k_{2}+l_{2}), & \rho_{q} & \equiv \Tilde{p}_{q}+l_{1}-\frac{\zeta}{1-\eta}\,l_{2} = (k_{3}+l_{1}-l_{2})+\frac{\xi}{1-\eta}(k_{2}+l_{2}).
\end{align}
For the after-shock-wave terms only one dipole is present; we may still write $\int\frac{d^{2}l_{1}}{(2\pi)^{2}}\frac{d^{2}l_{2}}{(2\pi)^{2}}\mathcal{D}(l_{1})\mathcal{D}(l_{2})$ in front, because $\int\frac{d^{2}l}{(2\pi)^{2}}\mathcal{D}(l)=1$. Their arguments are $(\kappa,\Tilde{p})$ and $(\kappa_{q},\Tilde{p}_{q})$.

\noindent\textbf{Step 6: the elastic amplitudes.} Collecting the four pieces of each brace in Eq.~\eqref{eq:elcomb}, with $\int_{l_{1}l_{2}}\equiv\int\frac{d^{2}l_{1}}{(2\pi)^{2}}\frac{d^{2}l_{2}}{(2\pi)^{2}}\mathcal{D}(l_{1})\mathcal{D}(l_{2})$:
\begin{align}
\text{brace}\ = (2\pi)^{2}\,\mathcal{A}^{lr}_{\lambda_{1}\lambda_{2}},\hspace{2em} \mathcal{A} = \int_{l_{1}l_{2}}\Big[\Phi_{1}(\kappa,\rho)-\Phi_{2}(\kappa,\Tilde{p})\Big]-\Big[\Phi_{1}(\Tilde{k},\Tilde{p})-\Phi_{2}(\Tilde{k},\Tilde{p})\Big]\hspace{2em}\text{(antiquark)},
\end{align}
and $\Tilde{\mathcal{A}}$ is the same with $\Tilde{\Phi}$, $(\kappa_{q},\rho_{q})$, $(\kappa_{q},\Tilde{p}_{q})$, $(\Tilde{k}_{q},\Tilde{p}_{q})$. Since every $\Phi$ is a sum of a few fixed spinor structures times scalar functions, all the $l$-integrals sit in the following explicit integrals.

\noindent\textbf{Antiquark emission:}
\begin{align}\label{eq:Tqb}
T^{jn} & = \int_{l_{1}l_{2}}\Big[\kappa^{j}\rho^{n}\,f_{1}(\kappa,\rho)-\kappa^{j}\Tilde{p}^{n}\,f_{2}(\kappa,\Tilde{p})\Big]-\Tilde{k}^{j}\Tilde{p}^{n}\Big[f_{1}(\Tilde{k},\Tilde{p})-f_{2}(\Tilde{k},\Tilde{p})\Big],\notag \\
T^{j}_{k} & = \int_{l_{1}l_{2}}\Big[\kappa^{j}\,f_{1}(\kappa,\rho)-\kappa^{j}\,f_{2}(\kappa,\Tilde{p})\Big]-\Tilde{k}^{j}\Big[f_{1}(\Tilde{k},\Tilde{p})-f_{2}(\Tilde{k},\Tilde{p})\Big],\notag \\
T^{n}_{p} & = \int_{l_{1}l_{2}}\Big[\rho^{n}\,f_{1}(\kappa,\rho)-\Tilde{p}^{n}\,f_{2}(\kappa,\Tilde{p})\Big]-\Tilde{p}^{n}\Big[f_{1}(\Tilde{k},\Tilde{p})-f_{2}(\Tilde{k},\Tilde{p})\Big],\notag \\
T_{0} & = \int_{l_{1}l_{2}}\Big[f_{1}(\kappa,\rho)-f_{2}(\kappa,\Tilde{p})\Big]-\Big[f_{1}(\Tilde{k},\Tilde{p})-f_{2}(\Tilde{k},\Tilde{p})\Big],\hspace{3em} T_{I} = \int_{l_{1}l_{2}}g(\kappa,\rho)-g(\Tilde{k},\Tilde{p}).
\end{align}
\textbf{Quark emission:} $\Tilde{T}^{jn},\Tilde{T}^{j}_{k},\Tilde{T}^{n}_{p},\Tilde{T}_{0},\Tilde{T}_{I}$, the same formulas with $(\kappa,\rho)\rightarrow(\kappa_{q},\rho_{q})$, $\Tilde{p}\rightarrow\Tilde{p}_{q}$, $\Tilde{k}\rightarrow\Tilde{k}_{q}$ and $f_{1},f_{2},g\rightarrow\Tilde{f}_{1},\Tilde{f}_{2},\Tilde{g}$:
\begin{align}\label{eq:Tq}
\Tilde{T}^{jn} & = \int_{l_{1}l_{2}}\Big[\kappa^{j}_{q}\rho^{n}_{q}\,\Tilde{f}_{1}(\kappa_{q},\rho_{q})-\kappa^{j}_{q}\Tilde{p}^{n}_{q}\,\Tilde{f}_{2}(\kappa_{q},\Tilde{p}_{q})\Big]-\Tilde{k}^{j}_{q}\Tilde{p}^{n}_{q}\Big[\Tilde{f}_{1}(\Tilde{k}_{q},\Tilde{p}_{q})-\Tilde{f}_{2}(\Tilde{k}_{q},\Tilde{p}_{q})\Big],\ \ \text{etc.}
\end{align}
In terms of these,
\begin{align}\label{eq:Ael}
\mathcal{A}^{lr}_{\lambda_{1}\lambda_{2}} & = \sum_{\lambda}\Big[\varphi^{lj}_{\lambda_{1}\lambda}(\zeta)\tau^{rn}_{\lambda\lambda_{2}}(\xi,\eta)\,T^{jn}\notag \\ & \hspace{4em}+\xi\omega_{f}\,\varphi^{lj}_{\lambda_{1}\lambda}(\zeta)\mu^{r}_{\lambda\lambda_{2}}\,T^{j}_{k}+m_{f}\,\mu^{l}_{\lambda_{1}\lambda}\tau^{rn}_{\lambda\lambda_{2}}(\xi,\eta)\,T^{n}_{p}+m_{f}\xi\omega_{f}\,\mu^{l}_{\lambda_{1}\lambda}\mu^{r}_{\lambda\lambda_{2}}\,T_{0}\Big]\notag \\ & -\xi\eta\left(\delta^{lr}+2i\epsilon^{lr}\lambda_{1}\right)\delta_{\lambda_{1}\lambda_{2}}\,T_{I},\notag \\
\Tilde{\mathcal{A}}^{lr}_{\lambda_{1}\lambda_{2}} & = \sum_{\lambda}\Big[\tau^{nr}_{\lambda_{1}\lambda}(\xi,\zeta)\varphi^{lj}_{\lambda\lambda_{2}}(\zeta+\xi)\,\Tilde{T}^{jn}\notag \\ & \hspace{4em}+\xi\Tilde{\omega}_{f}\,\mu^{r}_{\lambda_{1}\lambda}\varphi^{lj}_{\lambda\lambda_{2}}(\zeta+\xi)\,\Tilde{T}^{j}_{k}+m_{f}\,\tau^{nr}_{\lambda_{1}\lambda}(\xi,\zeta)\mu^{l}_{\lambda\lambda_{2}}\,\Tilde{T}^{n}_{p}+m_{f}\xi\Tilde{\omega}_{f}\,\mu^{r}_{\lambda_{1}\lambda}\mu^{l}_{\lambda\lambda_{2}}\,\Tilde{T}_{0}\Big]\notag \\ & -\xi\zeta\left(\delta^{lr}-2i\epsilon^{lr}\lambda_{1}\right)\delta_{\lambda_{1}\lambda_{2}}\,\Tilde{T}_{I}.
\end{align}
The coordinate integrals are now done:
\begin{align}
\int_{x,y,z,\Bar{x},\Bar{y},\Bar{z}}e^{-ik_{1}\cdot(x-\Bar{x})-ik_{2}\cdot(y-\Bar{y})-ik_{3}\cdot(z-\Bar{z})}\,\Big[\text{elastic part of the bracket}\Big] = S_{\perp}(2\pi)^{2}\delta^{(2)}(k_{1}+k_{2}+k_{3})\,(2\pi)^{4}\sum_{\lambda_{1}\lambda_{2}}\sum_{l,r}\mathcal{A}^{lr}_{\lambda_{1}\lambda_{2}}\mathcal{A}^{lr\,*}_{\lambda_{1}\lambda_{2}},
\end{align}
and correspondingly $\Tilde{\mathcal{A}}\Tilde{\mathcal{A}}^{*}$ (quark) and $\Tilde{\mathcal{A}}\mathcal{A}^{*}$ (interference).

\noindent\textbf{Step 7: helicity and polarisation sums.} Notation: $|T|^{2}\equiv\sum_{j,n}T^{jn}T^{jn*}$, $\ \mathrm{tr}\,T\equiv T^{11}+T^{22}$, $\ T_{k}\cdot T^{*}_{p}\equiv\sum_{j}T^{j}_{k}T^{j*}_{p}$, and
$$\epsilon^{ij}\epsilon^{mn}\,T^{im}U^{jn*} = T^{11}U^{22*}+T^{22}U^{11*}-T^{12}U^{21*}-T^{21}U^{12*}.$$
\textbf{Antiquark emission.} $\sum\mathcal{A}\mathcal{A}^{*} = H^{\text{no mass}}+H^{\text{mass}}$:
\begin{align}\label{eq:Hqb}
H^{\text{no mass}} & = 8\left[\zeta^{2}+(1-\zeta)^{2}\right]\left[\eta^{2}+(1-\zeta)^{2}\right]|T|^{2}-8\xi(2\zeta-1)\left(2(1-\zeta)-\xi\right)\epsilon^{ij}\epsilon^{mn}T^{im}T^{jn*}\notag \\ & -16\,\zeta\xi\eta^{2}\,\Re\big[\mathrm{tr}\,T\;T^{*}_{I}\big]+8\xi^{2}\eta^{2}\,|T_{I}|^{2},\notag \\
H^{\text{mass}} & = 8m^{2}_{f}\left[\eta^{2}+(1-\zeta)^{2}\right]T_{p}\cdot T^{*}_{p}+8\xi^{2}\omega^{2}_{f}\left[\zeta^{2}+(1-\zeta)^{2}\right]T_{k}\cdot T^{*}_{k}+8m^{2}_{f}\xi^{2}\omega^{2}_{f}\,|T_{0}|^{2}\notag \\
& -16\,\zeta\eta\,m_{f}\xi\omega_{f}\,\Re\big[\mathrm{tr}\,T\;T^{*}_{0}-T_{k}\cdot T^{*}_{p}\big]+16\,\xi\eta\,m_{f}\xi\omega_{f}\,\Re\big[T_{0}\,T^{*}_{I}\big].
\end{align}
\textbf{Quark emission.} $\sum\Tilde{\mathcal{A}}\Tilde{\mathcal{A}}^{*} = \Tilde{H}^{\text{no mass}}+\Tilde{H}^{\text{mass}}$:
\begin{align}\label{eq:Hq}
\Tilde{H}^{\text{no mass}} & = 8\left[\zeta^{2}+(1-\eta)^{2}\right]\left[\eta^{2}+(1-\eta)^{2}\right]|\Tilde{T}|^{2}-8\xi(2\eta-1)\left(2(1-\eta)-\xi\right)\epsilon^{ij}\epsilon^{mn}\Tilde{T}^{im}\Tilde{T}^{jn*}\notag \\ & +16\,\zeta^{2}\xi\eta\,\Re\big[\mathrm{tr}\,\Tilde{T}\;\Tilde{T}^{*}_{I}\big]+8\xi^{2}\zeta^{2}\,|\Tilde{T}_{I}|^{2},\notag \\
\Tilde{H}^{\text{mass}} & = 8m^{2}_{f}\left[\zeta^{2}+(1-\eta)^{2}\right]\Tilde{T}_{p}\cdot\Tilde{T}^{*}_{p}+8\xi^{2}\Tilde{\omega}^{2}_{f}\left[\eta^{2}+(1-\eta)^{2}\right]\Tilde{T}_{k}\cdot\Tilde{T}^{*}_{k}+8m^{2}_{f}\xi^{2}\Tilde{\omega}^{2}_{f}\,|\Tilde{T}_{0}|^{2}\notag \\
& +16\,\zeta\eta\,m_{f}\xi\Tilde{\omega}_{f}\,\Re\big[\mathrm{tr}\,\Tilde{T}\;\Tilde{T}^{*}_{0}-\Tilde{T}_{k}\cdot\Tilde{T}^{*}_{p}\big]+16\,\xi\zeta\,m_{f}\xi\Tilde{\omega}_{f}\,\Re\big[\Tilde{T}_{0}\,\Tilde{T}^{*}_{I}\big].
\end{align}
\textbf{Interference.} $\sum\Tilde{\mathcal{A}}\mathcal{A}^{*} = H^{x,\,\text{no mass}}+H^{x,\,\text{mass}}$:
\begin{align}\label{eq:Hx}
H^{x,\,\text{no mass}} & = 8\left[1-\xi-2\zeta\eta\right]\left[\xi(1-\xi)+2\zeta\eta\right]\sum_{j,n}\Tilde{T}^{jn}T^{jn*}-8\xi\left[1-2\zeta-\xi\right]^{2}\epsilon^{ij}\epsilon^{mn}\Tilde{T}^{im}T^{jn*}\notag \\ & -8\xi\eta(\zeta+\xi)^{2}\,\mathrm{tr}\,\Tilde{T}\;T^{*}_{I}+8\xi\zeta(1-\zeta)^{2}\,\Tilde{T}_{I}\;\mathrm{tr}\,T^{*},\notag \\
H^{x,\,\text{mass}} & = 8m^{2}_{f}(\xi+2\zeta\eta)\,\Tilde{T}_{p}\cdot T^{*}_{p}-8(\xi+2\zeta\eta)\,\xi^{2}\Tilde{\omega}_{f}\omega_{f}\,\Tilde{T}_{k}\cdot T^{*}_{k}\notag \\ & -8(\zeta+\xi)^{2}m_{f}\xi\omega_{f}\,\mathrm{tr}\,\Tilde{T}\;T^{*}_{0}+8(1-\zeta)^{2}m_{f}\xi\Tilde{\omega}_{f}\,\Tilde{T}_{0}\;\mathrm{tr}\,T^{*}\notag \\
& -8\eta^{2}m_{f}\xi\Tilde{\omega}_{f}\,\Tilde{T}_{k}\cdot T^{*}_{p}+8\zeta^{2}m_{f}\xi\omega_{f}\,\Tilde{T}_{p}\cdot T^{*}_{k}.
\end{align}
(There is no $\Tilde{T}_{0}T^{*}_{0}$, no $\Tilde{T}_{I}T^{*}_{I}$ and no mass $\times$ instantaneous term in the interference: those helicity sums vanish.)

\subsection{Results}
\begin{align}\label{eq:resel}
\frac{d^{3}N^{\text{direct}}_{\text{el}}}{d^{3}k_{1}\,d^{3}k_{2}\,d^{3}k_{3}}\Bigg|_{\Bar{q}\text{ emits}} & = \frac{e^{2}g^{2}\zeta^{2}C_{F}N_{c}\;S_{\perp}}{8\,(2\pi)^{6}\,\xi\,q^{+\,2}}\,\delta\left(q^{+}-k^{+}_{1}-k^{+}_{2}-k^{+}_{3}\right)\delta^{(2)}(k_{1}+k_{2}+k_{3})\sum_{f}e^{2}_{f}\,\Big[H^{\text{no mass}}+H^{\text{mass}}\Big],\notag \\
\frac{d^{3}N^{\text{direct}}_{\text{el}}}{d^{3}k_{1}\,d^{3}k_{2}\,d^{3}k_{3}}\Bigg|_{q\text{ emits}} & = \frac{e^{2}g^{2}\eta^{2}C_{F}N_{c}\;S_{\perp}}{8\,(2\pi)^{6}\,\xi\,q^{+\,2}}\,\delta\left(q^{+}-k^{+}_{1}-k^{+}_{2}-k^{+}_{3}\right)\delta^{(2)}(k_{1}+k_{2}+k_{3})\sum_{f}e^{2}_{f}\,\Big[\Tilde{H}^{\text{no mass}}+\Tilde{H}^{\text{mass}}\Big],\notag \\
\frac{d^{3}N^{\text{cross}}_{\text{el}}}{d^{3}k_{1}\,d^{3}k_{2}\,d^{3}k_{3}} & = -\frac{e^{2}g^{2}\eta\zeta\,C_{F}N_{c}\;S_{\perp}}{8\,(2\pi)^{6}\,\xi\,q^{+\,2}}\,\delta\left(q^{+}-k^{+}_{1}-k^{+}_{2}-k^{+}_{3}\right)\delta^{(2)}(k_{1}+k_{2}+k_{3})\sum_{f}e^{2}_{f}\;2\,\Re\Big[H^{x,\,\text{no mass}}+H^{x,\,\text{mass}}\Big].
\end{align}
The prefactors are those of Sec.~3 times $S_{\perp}(2\pi)^{2}\cdot(2\pi)^{4}$ from Steps 3 and 6. The sum of the two direct terms and the interference is the full elastic trijet cross section, $\propto\sum|\zeta\,\mathcal{A}-\eta\,\Tilde{\mathcal{A}}|^{2}$. The momenta in the $T$'s are on the support $k_{1}+k_{2}+k_{3}=0$, i.e.\ $\Tilde{k}=k_{1}$, $\Tilde{p}=k_{3}+\frac{\xi}{1-\zeta}k_{1}$, $\Tilde{k}_{q}=-k_{2}$, $\Tilde{p}_{q}=k_{3}+\frac{\xi}{1-\eta}k_{2}$.

\noindent\textbf{What ``no mass'' and ``mass'' mean.} $H^{\text{no mass}}$ collects the helicity-conserving terms, whose numerators carry no explicit $m_{f}$. $H^{\text{mass}}$ collects the helicity-flip terms; every one of them carries $m^{2}_{f}$ (note $\omega_{f}=\xi m_{f}/(1-\zeta)$). The quark mass also enters the $T$'s through the scales $\epsilon^{2}_{f}=\zeta(1-\zeta)Q^{2}+m^{2}_{f}$ and $\omega_{f}$ in $f_{1},f_{2},g$. For $m_{f}\rightarrow0$: $H^{\text{mass}}\rightarrow0$, $\epsilon^{2}_{f}\rightarrow\zeta(1-\zeta)Q^{2}$, $\omega_{f}\rightarrow0$, and $H^{\text{no mass}}$ gives the massless elastic cross section.

\noindent\textbf{Checks.} (i) Without a target ($S=1$, $\mathcal{D}(l)=(2\pi)^{2}\delta^{(2)}(l)$), $\kappa=\Tilde{k}$, $\rho=\Tilde{p}$ and all $T$'s vanish, as they must. (ii) The factorisation of the elastic parts, the dipole shifts and the change of variables were verified symbolically. The convolution formula was tested numerically against the coordinate integral. Formulas \eqref{eq:Hqb}--\eqref{eq:Hx} were compared with brute-force sums of $\mathcal{A}$, $\Tilde{\mathcal{A}}$ built from Eqs.~\eqref{eq:Phiel}, \eqref{eq:Phielq} over all helicities and polarisations, for random momenta, kicks, masses and fractions (agreement $10^{-15}$; Appendix~\ref{app:checks4}).

\appendix
\section{Corrections to the supplied file}\label{app:corr}
Each item gives \textbf{[equation number(s) and page(s) in this PDF]}, then the line number(s) in the supplied \texttt{Trijet\_massive.tex}. In the marked version the changes are printed in {\color{red}red}. The new Gaussian (large-$N_{c}$) expansions of all twelve $\mathcal{W}$'s in Sec.~3 were checked and are correct. Items 1--4 are errors that were corrected in the previous rounds and are back in the supplied file.

\noindent\textbf{Errors}
\begin{enumerate}
    \item \textbf{[Eq.~\eqref{chg:pFT2}, p.~\pageref{chg:pFT2}]} (L224): in the after-shock-wave transform the exponent must be $e^{-i\Tilde{p}\cdot Y}$, not $e^{+i\Tilde{p}\cdot Y}$. This is the phase of the state (L143) and of the $\Tilde{p}$ transform of L221. With $e^{+i\Tilde{p}\cdot Y}$ the $\mathcal{A}^{(01)}$ term would carry $+iY^{n}$ instead of $-iY^{n}$.
    \item \textbf{[Eq.~\eqref{chg:omq}, p.~\pageref{chg:omq}]} (L566): in the double-flip term of the quark-channel after-shock-wave vertex $\Tilde{\Psi}[\Tilde{\mathcal{A}}]$ the gluon-flip factor is $\xi\Tilde{\omega}_{f}=\xi^{2}m_{f}/(1-\eta)$, not $\xi\omega_{f}=\xi^{2}m_{f}/(1-\zeta)$. It must match $\Tilde{\Psi}[\Tilde{\mathcal{G}}]$ on the same line and the Dirac-spinor result for the quark line.
    \item \textbf{[Eq.~\eqref{chg:dena}, p.~\pageref{chg:dena}; Eq.~\eqref{chg:denb}, p.~\pageref{chg:denb}]} (L213, L344): in the energy denominator, $-E_{q}(k'')+E_{\Bar{q}}(p'')$ must be $-E_{q}(k'')-E_{\Bar{q}}(p'')$. The reduced denominators used afterwards were already correct.
    \item \textbf{[Eq.~\eqref{chg:dupM}, p.~\pageref{chg:dupM}]} (L538): $\Tilde{\mathcal{M}}_{1}(\Tilde{\mathcal{G}})$ is written twice in a row, which reads as a product $\Tilde{\mathcal{M}}_{1}\Tilde{\mathcal{M}}_{1}$.
\end{enumerate}
\noindent\textbf{Wording and notation}
\begin{enumerate}\setcounter{enumi}{4}
    \item \textbf{[p.~\pageref{chg:ba}]} (L241--L242): ``$(b)=(a)|_{y,z\rightarrow y'}$'' holds literally only for $m_{f}=0$. With masses, $\mathcal{A}^{(\alpha\beta)}$ and $\mathcal{G}^{(\alpha\beta)}$ are different functions; only their $Y\rightarrow0$ limits agree.
    \item \textbf{[p.~\pageref{chg:Y2RY}]} (L265--L266): the sentence refers to the factor $\frac{Y^{2}}{R\cdot Y}$ of the massless effective vertex, which no longer appears, so the sentence was rephrased.
    \item \textbf{[p.~\pageref{chg:Drule}]} (L297): ``$D\rightarrow\Tilde{Q}\Tilde{R}$'' is the massless after-shock-wave rule. Also, for the quark channel the photon scale is $\eta(1-\eta)Q^{2}$, not $\Tilde{Q}^{2}=\zeta(1-\zeta)Q^{2}$.
    \item \textbf{[Eq.~\eqref{chg:PsiAq}, p.~\pageref{chg:PsiAq}]} \textbf{[Eq.~\eqref{chg:M2q}, p.~\pageref{chg:M2q}]} \textbf{[p.~\pageref{chg:Gqarg}]} (L348, L541, L548): $[\Tilde{A}]$ and $(\Tilde{A})$ become $[\Tilde{\mathcal{A}}]$, and the argument of $\Tilde{\mathcal{G}}^{(\alpha\beta)}$ is $(\Tilde{R},X)$, not $(R,X)$.
    \item \textbf{[p.~\pageref{chg:ty}; Eq.~\eqref{chg:ty2}, p.~\pageref{chg:ty2}; Eq.~\eqref{chg:ty3}, p.~\pageref{chg:ty3}]} (L551, L645, L649): the typo ``Explicilty'', and ``,'' $\rightarrow$ ``;'' in $\mathcal{H}^{x}[\cdot\,;\cdot]$.
\end{enumerate}
\noindent\textbf{Not an error, but an assumption to state:} the quark-direct structures are written as $\mathcal{W}(y,x,z,\Bar{y},\Bar{x},\Bar{z})$ etc. These are the complex conjugates of the physical colour structures of the quark-emission amplitudes (Appendix~\ref{app:checks4}, check J). They equal the physical structures only for a C-even target, i.e.\ when $\langle\mathcal{W}\rangle$ is real. The Gaussian averages used here are C-even. In Sec.~\ref{sec:el} we use $S(a,b)=S(|a-b|)$, so $S(y,z)S(z,x)=S(x,z)S(z,y)$ and the question does not arise.

\noindent\textbf{Overall normalisation:} the prefactors $\frac{e^{2}g^{2}\zeta^{2}C_{F}N_{c}}{8(2\pi)^{12}\xi q^{+2}}$ etc.\ are those of the massless note; the mass does not change them. The checks verify everything relative to these prefactors (signs and weights of all terms, channels and helicity paths), but not the absolute normalisation itself.

\noindent\textbf{The old ``Fourier transforms'' section} (L720--L844) is inside \verb|\begin{comment}...\end{comment}| and does not appear in the PDF. It is replaced by Sec.~\ref{sec:el}.

\section{How the results were checked}\label{app:checks4}
The scripts are in \texttt{trijet/checks/massive/} of the repository.
\begin{enumerate}
    \item[E--I] (Secs.~1--3) The massive two-component vertices were fitted to Lepage--Brodsky spinors. The complete amplitude (both emitters, regular $+$ instantaneous) was compared with the Feynman/Dirac-spinor amplitude for all helicities and polarisations. The radial functions $\mathcal{G}^{(\alpha\beta)}$, $\mathcal{A}^{(\alpha\beta)}$, $\mathcal{G}_{E}$ were compared with an independent Hankel evaluation (agreement $10^{-13}$). All helicity sums $\mathcal{H}$, $\Tilde{\mathcal{H}}$, $\mathcal{H}^{x}$ were computed symbolically.
    \item[J] \texttt{checkJ\_colour.py}: every large-$N_{c}$ Wilson-line structure of Sec.~3 against the exact colour trace $\frac{1}{C_{F}N_{c}}\sum_{a}\text{Tr}[C^{a}_{A}C^{a\dagger}_{B}]$ of the amplitude colour matrices, for random $SU(N)$ matrices, $N=4,\ldots,32$. The difference falls like $1/N^{2}$.
    \item[K] \texttt{checkK\_dipole\_ft.py}: the changes of variables (Jacobian $=1$), the phase decomposition \eqref{eq:phase}, the dipole arguments and momentum shifts \eqref{eq:kappa}, symbolically. Also a numerical test of the convolution formula of Step~5 with Gaussian test functions: coordinate integral vs.\ momentum convolution agree to $2\times10^{-11}$.
    \item[N] \texttt{checkN\_elastic.py}: (a) symbolically, the $F$-free part of each of the twelve Gaussian expressions of Sec.~3 equals the factorised form of Step~1. (b) For a discrete normalised kick distribution $\mathcal{D}$ and random momenta, the amplitudes $\mathcal{A}$, $\Tilde{\mathcal{A}}$ were built directly from Eqs.~\eqref{eq:Phiel}, \eqref{eq:Phielq} and the shifted arguments \eqref{eq:kappa}, and $\sum\mathcal{A}\mathcal{A}^{*}$, $\sum\Tilde{\mathcal{A}}\Tilde{\mathcal{A}}^{*}$, $\sum\Tilde{\mathcal{A}}\mathcal{A}^{*}$ were summed by brute force over $\lambda_{1},\lambda_{2},l,r$. These sums were compared with $H^{\text{no mass}}+H^{\text{mass}}$, Eqs.~\eqref{eq:Hqb}--\eqref{eq:Hx}, for $(\zeta,\xi,Q,m_{f})=(0.3,0.2,1.1,0.8)$, $(0.45,0.15,0.7,1.6)$ and $(0.2,0.5,1.4,0)$. Agreement is $10^{-15}$ in every channel.
\end{enumerate}

\end{document}